\documentclass[11pt,a4paper]{article}
\usepackage[T1]{fontenc}
\usepackage[utf8]{inputenc}
\usepackage{lmodern,amsmath,amssymb,amsthm,mathtools,microtype,mathrsfs}
\usepackage[margin=25mm,top=23mm,bottom=25mm]{geometry}
\usepackage{booktabs,array,longtable,enumitem,xcolor,fancyhdr}
\usepackage[colorlinks=true,linkcolor=blue!45!black,urlcolor=blue!45!black,citecolor=blue!45!black]{hyperref}
\hypersetup{pdftitle={Reading the Alpoge--Claude RH and Dirichlet-GRH proof: working version 0.52},pdfauthor={ChatGPT},pdfsubject={Partial human-readable reconstruction with explicit arithmetic inputs},pdfkeywords={Riemann hypothesis, Lean, theta lifts, prime sums, Arakelov slopes}}
\setlength{\parindent}{0pt}
\setlength{\parskip}{5pt plus 1pt}
\setlist{topsep=3pt,itemsep=3pt,leftmargin=*}
\setlength{\headheight}{14pt}
\pagestyle{fancy}\fancyhf{}
\fancyhead[L]{\small RH and Dirichlet GRH: a reconstruction}
\fancyhead[R]{\small Working version 0.52}
\fancyfoot[C]{\thepage}
\renewcommand{\headrulewidth}{0.3pt}
\numberwithin{equation}{section}
\newtheorem{theorem}{Theorem}[section]
\newtheorem{proposition}[theorem]{Proposition}
\newtheorem{lemma}[theorem]{Lemma}
\newtheorem{corollary}[theorem]{Corollary}
\theoremstyle{remark}\newtheorem{remark}[theorem]{Remark}
\newcommand{\Rea}{\operatorname{Re}}
\newcommand{\HS}{\mathrm{HS}}
\newcommand{\tr}{\operatorname{tr}}
\newcommand{\codim}{\operatorname{codim}}
\newcommand{\covol}{\operatorname{covol}}
\newcommand{\OK}{\mathcal O_K}
\newcommand{\src}[2]{\href{https://github.com/CocoSgt/riemann-hypothesis-lean/blob/7cbf7a0c32fca1badf05931879a37c3f8a2b03b4/src/#1}{\texttt{#2}}}
\newcommand{\status}[1]{\textbf{Status:} #1}
\begin{document}
\begin{center}
{\LARGE\bfseries RH and Dirichlet GRH: new human proofs}\\[7pt]
{\large Version 0.52 --- 23 November 2026}\\[5pt]
{\large Bessel nonvanishing, section laws and Selmer signatures}
\end{center}
\textbf{Scope.} This edition contains the current audit and all new chapters.
It is a partial human reconstruction of the accepted formal proof.
Earlier named claims are in the cumulative manuscript and preserved components.
Each result below states its original inputs and exact remaining producers.
\setcounter{tocdepth}{1}
\tableofcontents
\clearpage
\begingroup\emergencystretch=3em
\section{Current completion audit: Bessel nonvanishing and the comparison steps}
\label{sec:dependency-audit-v52}

\subsection{Status of the human reconstruction}

We accept the pinned Alpoge--Claude formal proof as correct. The task
is to write ordinary mathematical proofs and preserve their exact
connections. Version 0.52 adds 87 numbered claims to the 3192
preserved predecessor claims, giving 3279 in all. The selected
sufficient graph has 121 nodes and 161 edges. There remain 24
whole open obligations: 17 class C and 7 class B. Of these,
17 are shared by RH and unrestricted Dirichlet GRH and
7 are additional Dirichlet-GRH inputs. These are inventories,
not completion percentages or estimates of the remaining effort.

\subsection{The four new chapters}

\paragraph{GSL1--GSL26: bases that work after every coefficient change.}
The algebraic core constructs finite matrices satisfying
\[
 D_0B=0,\qquad CB=rI,\qquad BC+HD_0=rI.
\]
Whenever $r$ becomes a unit, the columns of $B$ are a basis of the
actual kernel, with inverse coordinate map $r^{-1}C$. The equations
remain valid after every ring map, including nonflat maps and maps
to rings with nilpotents. This supplies one basis before all later
tests are chosen.

A bounded flat section complex, descended from its supplied
Noetherian cube model, has a finite-projective model valid after
every tensor product. Its fixed chart functions give every field of
the literal \texttt{Perf} record. The geometric rank, vanishing
and translated-power arguments are proved from their exact named
inputs. Applied to the actual bundle, the universal matrix argument
then gives the whole conditional \texttt{Sec\allowbreak Base\allowbreak Change\allowbreak Any/\allowbreak St} and
\texttt{Sec\allowbreak Twist\allowbreak St}. Over a local ring the denominator nonzero in
the residue field is already invertible, so every residue section
lifts on the original twist bundle.

The actual geometric suppliers remain: the section-cube model and
finite cohomology/smooth descent, polarization base change, Serre
and multiplication/invariant descent, geometric integrality, initial rank and section
decomposition/growth, and theorem-of-square generators. Supplied
\texttt{Moduli} and the inherited uniformization foundations remain
explicit. GSL26 gives their exact installation in the same
TDR/NFR/NFL family; it does not assert these global inputs exist.

\paragraph{AFE1--AFE21: the full arithmetic entries on the same frame.}
The character CRT decomposition, image-primary decomposition,
quadratic Gauss sums and dyadic conductor calculation supply all
nineteen \texttt{Gen} fields for both the Dat and Host branches.
The Host argument retains $\ell=3$ and supplies a local witness
at every specified place. The auxiliary \texttt{Is\allowbreak Host\allowbreak Frame}
and full \texttt{exists\_\allowbreak sigma} are also constructed.

The CAF modulus is chosen before the LFC field and witness prime.
The same field, frame, $C_0$ and marked $g_1$ then satisfy the
whole $P\Rightarrow\texttt{Gd\allowbreak Frame\allowbreak Ex}$ and
$P\Rightarrow\texttt{Host\allowbreak Frame\allowbreak Ex}$ statements, including the
actual character seed. All five marking conditions hold.

The ordinary-point representation descends through the actual
Galois quotient over every original local Noetherian domain.
The exact triangular characters, polarization, residual lift and
adic continuity survive. No completeness or flatness premise is
added. The full \texttt{Fin\allowbreak Pt}\!$\Rightarrow$\!\texttt{Fin\allowbreak End}
theorem chooses the one subgroup $G_L$ before every coefficient
algebra and point. The original $P$, actual finite-point producer
and automorphic seeds/endpoints remain separate suppliers.

\paragraph{BQS1--BQS23: the actual quotients and full Bessel nonvanishing.}
The standard and dual Euler identities give the exact eigenspace
projections. The highest-weight argument controls the full integral
root cone and its extremal line. A compactly averaged trace-kernel
identity transfers the scalar-model calculation to the actual
representation. Minimal-type extraction works for every nonzero
unitary image, including targets that are not complete, with the
original $h-2$ margin.

The generic quotient construction is applied to the actual tensor
coefficient map and the simultaneous kernel of all the prescribed
Bessel periods. It gives the literal function space, both directions
of its smooth $L^2$ image relation, stability under the actual compact
group, the shifted lowest type and Casimir, and period injectivity.
These are both full original \texttt{Quot\allowbreak P} and \texttt{Quot\allowbreak M}.
The inherited evaluated periods now supply the full
\texttt{Step\allowbreak Statement} and \texttt{Bessel\allowbreak NVGap}.
The setting return still has its original \texttt{Forms0}
premise. Orbital comparisons, uniqueness, transport and the global
normalized spectral estimate remain distinct work.

\paragraph{SSC1--SSC17: the actual Selmer-signature comparison.}
Trace equality gives conjugacy over the original local coefficient
ring. The local bridges transport crystalline cocycles. For an
arbitrary injective local map of DVRs $O\to O'$, a basis of
$O'/\pi O'$ lifts compatibly through every quotient by $\pi^n$.
It may be infinite; each continuous cocycle into a discrete torsion
module has finite support by compactness. The coordinates therefore
compare all the original local spans and yield the full coefficient
change theorem, including infinite module lengths.

The same-character relation produces a common DVR. The original
Selmer module injects into the signature module with its literal
local conditions. A common-place valuation argument and tame inertia
give the conductor-root bound. One attached signature consequently
has every original \texttt{Sel\allowbreak Core} field, with the same twist
orientation and the exact $1/\ell^2$ budget required by
\texttt{Kill\allowbreak At}.

Thus the whole original \texttt{Signature\allowbreak Choice} supplier is now
available on the DFD representation family. The setup implication
DFD15 retains only the coherent Wach \texttt{HW/\allowbreak Fon/\allowbreak Cmp} supplier.
The subsequent local row bounds must still construct one
\texttt{Concl\allowbreak F} family. DNE already proves its full export to
\texttt{Adm\allowbreak F} and \texttt{Link\allowbreak N}, with
$K=2000(|C_a|+1)^3$ chosen before the family.

\subsection{The exact frontier}

No additional whole selected obligation is closed. The new complete
subtheorems and their remaining producer premises are recorded
separately. The v0.49 \texttt{GLOBAL\_\allowbreak LOCAL\_\allowbreak GAIN} closure
remains intact. The full original whole contracts below are preserved
from the selected graph; their version-0.52 progress paragraphs
govern the current status. Historical qualifications remain visible
for traceability. Every new claim has an independent mathematical
review bound to exact proof and source bytes. No new Lean replay
or mechanically extracted transitive Lean audit is asserted.
\paragraph{1. GLOBAL\_\allowbreak DENOMINATOR — C — Actual LinkN local table and binding producer}

Closed(Des $\mathcal C$ 3) LinkN: the actual etaSum is bounded with exact dZ0 and same-host vvTerm, subject only to the literal ShaFinOdd and summability guards.

\emph{Shared objects and choice order.} one FamP; one class $\mathcal A$ and its target-compatible subclass; one coefficient field per family datum; one x and i at each application; original selected rows, words and levels; thresholds chosen before the common family; actual D.nd; actual etaJ and Aj; sha\_\allowbreak eq

\emph{Completed inputs (version 0.52).} SSC1-SSC17 (full original Selmer SignatureChoice and family setup retaining only coherent Wach comparison), DFD1-DFD15, DNE1-DNE14 (full original ConclF-to-LinkN consumer)

\emph{Remaining producer scope.} SSC proves the full original Selmer SignatureChoice supplier on the actual DFD complement family: integral trace conjugacy, crystalline local comparison, arbitrary injective local DVR coefficient change, the original signature inclusion, normalized finite/infinite lengths and the exact tame-conductor term. The signature uses the same oriented character as TwGood and gives DFD KillAt with its exact 1/ell\textasciicircum{}2 budget. SSC17 reduces the remaining DFD15 setup supplier to one coherent HW/Fon/Cmp on the actual canonical datum. Construct the all-layer Wach lattices, actual etale and Wach modules, Herr class maps and comparison, retaining the Bhatt-Scholze input at its stated scope. Then prove every original local row/table bound on one common family to obtain ConclF, retaining ShaFinOdd, exceptional-set scope, vvTerm summability and uniform thresholds. DNE already supplies the full ConclF-to-AdmF-to-LinkN consumer with K=2000(|Ca|+1)\textasciicircum{}3 chosen before the family. GLOBAL\_\allowbreak DENOMINATOR remains open until its actual common ConclF producer is supplied.

\emph{Source.} \href{https://github.com/CocoSgt/riemann-hypothesis-lean/blob/7cbf7a0c32fca1badf05931879a37c3f8a2b03b4/src/Formalism/SharpRestriction/GlobalTheoremInputs.lean#L45-L48}{LinkN}.

\emph{Source.} \href{https://github.com/CocoSgt/riemann-hypothesis-lean/blob/7cbf7a0c32fca1badf05931879a37c3f8a2b03b4/src/Formalism/SharpRestriction/GlobalTheoremInputs.lean#L40-L44}{LinkN\_\allowbreak at}.

\emph{Source.} \href{https://github.com/CocoSgt/riemann-hypothesis-lean/blob/7cbf7a0c32fca1badf05931879a37c3f8a2b03b4/src/Development/AutomorphicArithmetic/CharacterBoxes/DivinNOfHypV.lean#L117-L134}{divinN\_\allowbreak of\_\allowbreak hypV}.



\paragraph{2. GLOBAL\_\allowbreak RETAINED\_\allowbreak DEFECT — C — Original retained-row defect bound LinkIRR}

Rows.DelpNonneg and delp$\le$exp(-x\textasciicircum{}(ez/7)) uniformly for every late actual row datum. CURRENT VERSION0.47 QUALIFICATION: Finish the same-D/common-family adelic setting after ASG: normalized nonzero Bessel W and decay, measured-field block, original Siegel-Weil rescaling and measure/character ties, Bessel uniqueness/local ties, then the original AutTie/RowDictTie/LabTie. RowsFinDim is now proved universally by RFD. Supply the actual CountIdempotentsClass theorem or its three sufficient arithmetic cylinder counts, with one simultaneous family threshold; reuse the completed ARO structure. Completed new original producers: ARO1-ARO30, RFD1-RFD9, ASG1-ASG18. Any earlier narrative naming these completed producers as outstanding is historical; this qualification is the current scope authority. CURRENT VERSION0.48 QUALIFICATION: Finish the same-D/common-family adelic setting after ASG and AMD: normalized nonzero Bessel W and decay, original orbital/fiber and Siegel-Weil measure/character ties, Bessel uniqueness/local ties, and the full original transport theorem, then AutTie/RowDictTie/LabTie. The complete MExists measured-data block and unipotent covolume normalization are now supplied, and universal RowsFinDim remains available. Supply the actual CountIdempotentsClass theorem or its three sufficient arithmetic cylinder counts with one simultaneous family threshold. Completed current producers or conditional consumers: ARO1-ARO30, RFD1-RFD9, ASG1-ASG18, AMD1-AMD17. Earlier frontier descriptions are historical; this qualification is the current scope authority. Original whole-contract premises and quantifiers are retained. CURRENT VERSION0.49 QUALIFICATION: Prove the exact original BesselNVGap nonzero integral on its supplied discrete-series, compact-representation and character data. CoefWeightedLoc, WlocExists with its original nonvanishing antecedent, WFromLoc, the normalized global product and all-later-measured-block decay are complete. On the inherited Forms0 scope this leaves BesselNVGap for ConcreteSettingExists. Continue original orbital/fiber and Siegel-Weil measure/character ties, the required Bessel uniqueness/local ties and the full original transport theorem, then AutTie/RowDictTie/LabTie. Supply the actual CountIdempotentsClass theorem or its three sufficient arithmetic cylinder counts with one simultaneous family threshold. Completed current producers or conditional consumers: ARO1-ARO30, RFD1-RFD9, ASG1-ASG18, AMD1-AMD17, BWP1-BWP24 at their exact BesselNVGap/Forms0 scope. Earlier frontier descriptions are historical; this qualification is the current scope authority. Original whole-contract premises and quantifiers are retained.

\emph{Shared objects and choice order.} one FamP; one class $\mathcal A$ and its target-compatible subclass; one coefficient field per family datum; one x and i at each application; original selected rows, words and levels; thresholds chosen before the common family; AllRows A; LabelCountGoodInjective; AllJoinA; CountAllRows

\emph{Completed inputs (version 0.52).} SSC full Selmer-signature comparison, DFD common-family subresults, DNE full separate denominator consumer, WRS1-WRS10 full selected-class local gain, Inherited first DivDef3 and row-pair integrability

\emph{Remaining producer scope.} Complete the original F3 local coefficient bridges and F4 common-overfield readings, with valuation/norm/GoodCoef/iq and the entire coefficient field, and actual complex Hecke standard/no-Jordan plus arithmetic character data on the same form datum. Construct the actual base rows, host/frame lengths, nonzero chains and group ties, MultEq support, eigensystem/carrier/parts, all-embedding logLam, HD/d3b and NEb. Keep the stronger second DivRows3 scope distinct from the completed first DivDef3 and row-pair integrability. The full arithmetic entry and conditional common-frame construction do not by themselves install HOST\_\allowbreak NUMBER\_\allowbreak S, HOST\_\allowbreak FRAME\_\allowbreak LENGTH, TASK\_\allowbreak HOST\_\allowbreak VAN\_\allowbreak E\_\allowbreak INSTALL, TASK\_\allowbreak HOST\_\allowbreak XPP\_\allowbreak INSTALL or TASK\_\allowbreak COUNT\_\allowbreak V\_\allowbreak INSTALL. Continue the distinct retained-defect, period-height and normalized spectral producers, and the original compatible second-layer selection on one common arithmetic/geometric family. GLOBAL\_\allowbreak LOCAL\_\allowbreak GAIN remains fully closed at its original v0.49 scope. Complete local comparison theorems do not themselves construct the whole common family. Common design existence, NBUILT compatibility, original true-column records, shared fields/embeddings, Fits and eta agreement, simultaneous rows and eventual-family thresholds remain explicit. A stronger historical WPt/TSel construction may serve other producers but is not needed to reopen the proved selected-class local-gain estimate.

\emph{Source.} \href{https://github.com/CocoSgt/riemann-hypothesis-lean/blob/7cbf7a0c32fca1badf05931879a37c3f8a2b03b4/src/Formalism/SharpRestriction/GlobalTheoremInputs.lean#L55-L58}{LinkIRR}.

\emph{Source.} \href{https://github.com/CocoSgt/riemann-hypothesis-lean/blob/7cbf7a0c32fca1badf05931879a37c3f8a2b03b4/src/Development/AbelianVarietyTheory/Proofs/GlobalInterpolation/LinkIRROfJoinEx.lean#L17-L38}{linkIRR\_\allowbreak of\_\allowbreak join\_\allowbreak ex}.

\emph{Source.} \href{https://github.com/CocoSgt/riemann-hypothesis-lean/blob/7cbf7a0c32fca1badf05931879a37c3f8a2b03b4/src/Development/AbelianVarietyTheory/Proofs/GlobalInterpolation/LinkIRROfJoinEx.lean#L143-L166}{link\_\allowbreak irr}.



\paragraph{3. GLOBAL\_\allowbreak PERIOD\_\allowbreak HEIGHT — C — Full actual period and height comparison LinkDIVt}

Closed conditional TuSize -> DGamSize -> (LinkEta -> DIVinN -> LinkDIV), with every original LinkDIV\_\allowbreak at clause.

\emph{Shared objects and choice order.} one FamP; one class $\mathcal A$ and its target-compatible subclass; one coefficient field per family datum; one x and i at each application; original selected rows, words and levels; thresholds chosen before the common family; same RR target; same Bu/besS; same Nu and Bk

\emph{Source.} \href{https://github.com/CocoSgt/riemann-hypothesis-lean/blob/7cbf7a0c32fca1badf05931879a37c3f8a2b03b4/src/Formalism/SharpRestriction/GlobalTheoremInputs.lean#L77-L90}{LinkDIV\_\allowbreak at}.

\emph{Source.} \href{https://github.com/CocoSgt/riemann-hypothesis-lean/blob/7cbf7a0c32fca1badf05931879a37c3f8a2b03b4/src/Formalism/SharpRestriction/GlobalTheoremInputs.lean#L94-L96}{LinkDIVt}.

\emph{Source.} \href{https://github.com/CocoSgt/riemann-hypothesis-lean/blob/7cbf7a0c32fca1badf05931879a37c3f8a2b03b4/src/Development/AbelianVarietyTheory/Proofs/GlobalInterpolation/LinkBrest.lean#L7-L25}{GlobalTheory.Top.link\_\allowbreak brest}.

\emph{Source.} \href{https://github.com/CocoSgt/riemann-hypothesis-lean/blob/7cbf7a0c32fca1badf05931879a37c3f8a2b03b4/src/Development/AbelianVarietyTheory/Proofs/GlobalInterpolation/WireSel.lean#L11-L41}{GlobalTheory.Top.wire\_\allowbreak sel}.



\paragraph{4. GLOBAL\_\allowbreak SPECTRAL — C — Original normalized kernel/Gram estimates LinkS}

Closed TuSize -> LinkS with exact Gamma normalization, off-diagonal Hilbert-Schmidt bound, diagonal error and Dall/Dstar estimate, errors exp(-x\textasciicircum{}(epsilon/8)). CURRENT VERSION0.48 QUALIFICATION: Complete the original same-family spectral/Bessel and Siegel-Weil inputs with all measure, character, transport and common-object ties. The complete MExists measured block and prescribed-measure covolume normalization are supplied; these are not the full normalized spectral theorem. Completed current producers or conditional consumers: ASG1-ASG18, AMD1-AMD17, GCF1-GCF18 at their stated scopes. Earlier frontier descriptions are historical; this qualification is the current scope authority. Original whole-contract premises and quantifiers are retained. CURRENT VERSION0.49 QUALIFICATION: Complete the original same-family spectral/Bessel and Siegel-Weil inputs with every measure, character, transport and common-object tie. The normalized packet and decay are now constructed conditional on the exact BesselNVGap, and the analytic setting uses the original supplied Forms0 scope. These do not imply the full normalized spectral theorem. Completed current producers or conditional consumers: ASG1-ASG18, AMD1-AMD17, GCF1-GCF18 at their stated scopes, BWP1-BWP24 at their exact BesselNVGap/Forms0 scope, WRS1-WRS10 original local gain. Earlier frontier descriptions are historical; this qualification is the current scope authority. Original whole-contract premises and quantifiers are retained.

\emph{Shared objects and choice order.} one FamP; one class $\mathcal A$ and its target-compatible subclass; one coefficient field per family datum; one x and i at each application; original selected rows, words and levels; thresholds chosen before the common family; actual KD; same S.Cu spaces; GamOf; Dstar; Dall

\emph{Completed inputs (version 0.52).} BQS1-BQS23 (full QuotP/QuotM, StepStatement and BesselNVGap; setting at Forms0), BST1-BST11, BSE1-BSE9, BNV1-BNV23, BSP1-BSP10, BWP1-BWP24 at original Forms0 scope, AMD1-AMD24

\emph{Remaining producer scope.} BQS completes the original Euler annihilators, K2 weight-cone and highest-line suppliers, K3 minimal-type extraction, generic quotient theorem and both full QuotP/QuotM statements. The actual smooth L2 image and the simultaneous kernel of the prescribed periods are used. Together with BST/BSE/BNV/BSP this proves the full StepStatement and BesselNVGap. ConcreteSettingExists still has the original Forms0 premise. Next construct that actual global setting and the orbital/fiber and Siegel-Weil character/measure comparisons, Bessel uniqueness and local ties, full transport, AutTie/RowDictTie/LabTie and simultaneous arithmetic cylinder counts for CountIdempotentsClass. The local nonvanishing theorem does not yet supply the whole normalized GLOBAL\_\allowbreak SPECTRAL estimate.

\emph{Source.} \href{https://github.com/CocoSgt/riemann-hypothesis-lean/blob/7cbf7a0c32fca1badf05931879a37c3f8a2b03b4/src/Formalism/SharpRestriction/GlobalTheoremInputs.lean#L145-L149}{LinkS\_\allowbreak at}.

\emph{Source.} \href{https://github.com/CocoSgt/riemann-hypothesis-lean/blob/7cbf7a0c32fca1badf05931879a37c3f8a2b03b4/src/Formalism/SharpRestriction/GlobalTheoremInputs.lean#L152-L157}{LinkS}.

\emph{Source.} \href{https://github.com/CocoSgt/riemann-hypothesis-lean/blob/7cbf7a0c32fca1badf05931879a37c3f8a2b03b4/src/Development/AbelianVarietyTheory/Proofs/GlobalInterpolation/LinkAlt.lean#L16-L50}{GlobalTheory.Top.link\_\allowbreak alt}.



\paragraph{5. AUT\_\allowbreak UNRAM\_\allowbreak COEF — C — Actual unramified-coefficient automorphy}

For each actual Dat with coefficient ramification one and fixed iota, choose one marked g1 supplying FinA, AutB and LemW together. CURRENT VERSION0.48 QUALIFICATION: Construct the actual AutomorphyUnramifiedCoefficients package: for every original Dat with coefficient ramification one and fixed iota, choose one marked g1 that supplies FinA, AutB and LemW together. The original InpZ record, full RigII, CNR3, finite PropGK consumer and now the entire OneChar implication are supplied. Continue RH Cases I/V and the independent additional Dirichlet-GRH CaseIII/promodularity, SeedIV, tensor-factor and rank-four-chain tasks. Completed current producers or conditional consumers: OCH1-OCH30 (the full consumer implication), KIZ original InpZ, KTR1-KTR13, RNC16, PGK1-PGK6. Earlier frontier descriptions are historical; this qualification is the current scope authority. Original whole-contract premises and quantifiers are retained. CURRENT VERSION0.49 QUALIFICATION: Construct one common GdEnd for each original datum and fixed iota: the same marked g1 must supply FinEnd, AutEnd3 and AutEnd2. UCE now proves all ordinary local/curve/coefficient consumers and the exact implication to FinA/AutB/LemW, including arbitrary entry maps and every original outer-corner extension. The producer still needs the actual common abelian-CM GdFrameEx and LinkedEndE with both rank-two and rank-three runs, full ModUAlt/PfromAut and the two actual seeds, followed by the original local/automorphic descent. At this pin the alternative upstream route gives GdEndAlt; its additional local/descent fields must be matched to the actual original points, so no name-only GdEndAlt-to-GdEnd identification is allowed. Keep the supplied InpZ, RigII, CNR3, finite PropGK and whole OneChar implication. Continue RH Cases I/V and the independent additional Dirichlet-GRH CaseIII/promodularity, SeedIV, tensor-factor and rank-four-chain tasks. Completed current producers or conditional consumers: UCE1-UCE19 (complete GdEnd consumer), OCH1-OCH30, KIZ original InpZ, KTR1-KTR13, RNC16, PGK1-PGK6. Earlier frontier descriptions are historical; this qualification is the current scope authority. Original whole-contract premises and quantifiers are retained.

\emph{Completed inputs (version 0.52).} AFE1-AFE21 (full Gen entries, conditional frame joins, NRep descent and FinPt-to-FinEnd), LFC1-LFC23, CAF1-CAF11, MFC1-MFC14, UCE1-UCE19 (full original GdEnd consumer), OCH1-OCH30, KIZ original InpZ, KTR1-KTR13, RNC16, PGK1-PGK6

\emph{Remaining producer scope.} AFE proves full gen\_\allowbreak gd and gen\_\allowbreak host, all nineteen fields in each, the full implications from P to GdFrameEx and to HostFrameEx, the auxiliary IsHostFrame and exists\_\allowbreak sigma, all Marks fields, and full exists\_\allowbreak nrep3 at the original local Noetherian-domain scope. The whole finEnd\_\allowbreak of\_\allowbreak finPt implication is complete, with G\_\allowbreak L fixed before every coefficient algebra and point. The remaining producers are the original P where not already supplied, actual FinPt on this frame, AutPt/AutEnd3/AutEnd2 and the actual rank-two and rank-three automorphic seeds, followed by full ModUAlt/PfromAut and local/automorphic descent to GdEnd. UCE already supplies the complete GdEnd consumer. The alternative GdEndAlt route retains its additional local/descent fields. Do not reopen Gen entries, the common-frame join or the FinPt-to-FinEnd transport; do not replace the Host ell=3 scope by the stronger Dat guard.

\emph{Source.} \href{https://github.com/CocoSgt/riemann-hypothesis-lean/blob/7cbf7a0c32fca1badf05931879a37c3f8a2b03b4/src/Development/AutomorphicArithmetic/AutomorphicVarieties/AutomorphicCases/AutomorphyUnramifiedCoefficients.lean#L24-L27}{AutomorphyUnramifiedCoefficients}.

\emph{Source.} \href{https://github.com/CocoSgt/riemann-hypothesis-lean/blob/7cbf7a0c32fca1badf05931879a37c3f8a2b03b4/src/Development/AutomorphicArithmetic/AutomorphicVarieties/AutomorphicCases/AutomorphyUnramifiedCoefficients.lean#L12-L13}{FinA}.

\emph{Source.} \href{https://github.com/CocoSgt/riemann-hypothesis-lean/blob/7cbf7a0c32fca1badf05931879a37c3f8a2b03b4/src/Development/AutomorphicArithmetic/AutomorphicVarieties/AutomorphicCases/AutomorphyUnramifiedCoefficients.lean#L15-L17}{AutB}.

\emph{Source.} \href{https://github.com/CocoSgt/riemann-hypothesis-lean/blob/7cbf7a0c32fca1badf05931879a37c3f8a2b03b4/src/Development/AutomorphicArithmetic/AutomorphicVarieties/AutomorphicCases/AutomorphyUnramifiedCoefficients.lean#L19-L22}{LemW}.



\paragraph{6. TASK\_\allowbreak CASE3\_\allowbreak AUT\_\allowbreak JOIN — B — Complete the original Case-III automorphy/image join}

Install the actual Case-III ProModAlt, SeedIVAlt and shared PMExAlt with the inherited CohomologyHyp/Ferrero-Washington/HFit/SeedII facts into the original whole Case-III automorphy/image conclusion.

\emph{Source.} \href{https://github.com/CocoSgt/riemann-hypothesis-lean/blob/7cbf7a0c32fca1badf05931879a37c3f8a2b03b4/src/Development/AutomorphicArithmetic/AutomorphicVarieties/ModularityArgument/ThresholdArguments/DenseThresholdCases/PromodularityCaseThree.lean#L43-L44}{AutomorphyRankFourDense}.

\emph{Source.} \href{https://github.com/CocoSgt/riemann-hypothesis-lean/blob/7cbf7a0c32fca1badf05931879a37c3f8a2b03b4/src/Development/AutomorphicArithmetic/AutomorphicVarieties/ModularityArgument/ThresholdArguments/DenseThresholdCases/CuFinal.lean#L44-L44}{main4\_\allowbreak alt}.



\paragraph{7. TASK\_\allowbreak CASE3\_\allowbreak PROMOD\_\allowbreak JOIN — C — Reconstruct the actual Case-III chain-to-promodularity proof}

Use the original tensor-factor tower, ChainRankFour, ChainDesc4, seeded rank-four chain reduction, actual PMExAlt and SeedIVAlt to construct ProModAlt on every original Ctx3 star point at the same fixed iota, preserving all finite-coefficient, inertia and support maps.

\emph{Source.} \href{https://github.com/CocoSgt/riemann-hypothesis-lean/blob/7cbf7a0c32fca1badf05931879a37c3f8a2b03b4/src/Development/AutomorphicArithmetic/AutomorphicVarieties/ModularityArgument/ThresholdArguments/DenseThresholdCases/PromodularityCaseThreeOfLemmas.lean#L131-L153}{UpperTheory.V.Cu.promodularity\_\allowbreak case\_\allowbreak three\_\allowbreak of\_\allowbreak chain\_\allowbreak reduction}.

\emph{Source.} \href{https://github.com/CocoSgt/riemann-hypothesis-lean/blob/7cbf7a0c32fca1badf05931879a37c3f8a2b03b4/src/Development/AutomorphicArithmetic/AutomorphicVarieties/ModularityArgument/ThresholdArguments/DenseThresholdCases/PromodularityCaseThree.lean#L40-L42}{PromodularityCaseThreeDense}.



\paragraph{8. SEED\_\allowbreak IV\_\allowbreak ALT — C — Actual Case-III seed with prechosen residual datum}

Base Ctx3/Four tower, iota and Amb -> w,s before every later Dies tower -> compatible finite-coefficient seed with exact residual support/root entries.

\emph{Source.} \href{https://github.com/CocoSgt/riemann-hypothesis-lean/blob/7cbf7a0c32fca1badf05931879a37c3f8a2b03b4/src/Development/AutomorphicArithmetic/AutomorphicVarieties/ModularityArgument/ThresholdArguments/DenseThresholdCases/PromodularityCaseThree.lean#L19-L25}{SeedIVAlt}.



\paragraph{9. TENSOR\_\allowbreak FACTOR\_\allowbreak TOWER — B — Actual tensor-factor tower and characters}

Original finite-field, tame, factor and character constraints; source-level construction is sizeable and has not been reconstructed here.

\emph{Source.} \href{https://github.com/CocoSgt/riemann-hypothesis-lean/blob/7cbf7a0c32fca1badf05931879a37c3f8a2b03b4/src/Development/AutomorphicArithmetic/AutomorphicVarieties/ModularityArgument/ThresholdArguments/DenseThresholdCases/PromodularityCaseThreeOfLemmas.lean#L131-L153}{UpperTheory.V.Cu.promodularity\_\allowbreak case\_\allowbreak three\_\allowbreak of\_\allowbreak chain\_\allowbreak reduction}.



\paragraph{10. CHAIN\_\allowbreak RANK\_\allowbreak FOUR — C — Actual SomeIn-to-AllIn propagation}

Same source Space and PM ideal, residual point Good4, original bound 30(C+1)\textasciicircum{}2 <= d; geometry must preserve all identities.

\emph{Source.} \href{https://github.com/CocoSgt/riemann-hypothesis-lean/blob/7cbf7a0c32fca1badf05931879a37c3f8a2b03b4/src/Development/AutomorphicArithmetic/AutomorphicVarieties/ModularityArgument/ThresholdArguments/DenseThresholdCases/PromodularityCaseThree.lean#L26-L32}{ChainRankFour}.



\paragraph{11. CHAIN\_\allowbreak DESC4 — B — Original residual chain descent}

Finite residual descent retaining source Chain4 compatibility.

\emph{Source.} \href{https://github.com/CocoSgt/riemann-hypothesis-lean/blob/7cbf7a0c32fca1badf05931879a37c3f8a2b03b4/src/Development/AutomorphicArithmetic/AutomorphicVarieties/ModularityArgument/ThresholdArguments/DenseThresholdCases/PromodularityCaseThreeOfLemmas.lean#L131-L153}{UpperTheory.V.Cu.promodularity\_\allowbreak case\_\allowbreak three\_\allowbreak of\_\allowbreak chain\_\allowbreak reduction}.



\paragraph{12. RANK4\_\allowbreak CHAIN\_\allowbreak REDUCTION — C — Seeded reduction to the original residual chain}

Select T1 before every later Deep tower; retain one compatible seed and Chain4 for every original residual lift over Q.

\emph{Source.} \href{https://github.com/CocoSgt/riemann-hypothesis-lean/blob/7cbf7a0c32fca1badf05931879a37c3f8a2b03b4/src/Development/AutomorphicArithmetic/AutomorphicVarieties/ModularityArgument/ThresholdArguments/DenseThresholdCases/PromodularityCaseThree.lean#L33-L39}{RankFourChainReductionDense}.

\emph{Source.} \href{https://github.com/CocoSgt/riemann-hypothesis-lean/blob/7cbf7a0c32fca1badf05931879a37c3f8a2b03b4/src/Development/AutomorphicArithmetic/AutomorphicVarieties/ModularityArgument/ThresholdArguments/DenseThresholdCases/CuPChainIV.lean#L16-L82}{UpperTheory.V.Cu.chainIV\_\allowbreak alt}.



\paragraph{13. TASK\_\allowbreak COUNT\_\allowbreak V\_\allowbreak INSTALL — B — Complete the original CountV finite assembly}

From exactly the listed source host/cohomology inputs, complete countV\_\allowbreak of\_\allowbreak finiteness\_\allowbreak duality\_\allowbreak euler\_\allowbreak host, including original host/frame formation, normalized length, the CaseV guard and R1<=2\textasciicircum{}c2.

\emph{Completed inputs (version 0.52).} AFE full Host entry/auxiliary frame/local witnesses including ell=3, LFC1-LFC23, CAF9-CAF11, MFC1-MFC14, Inherited CountV consumers

\emph{Remaining producer scope.} AFE proves full gen\_\allowbreak gd and gen\_\allowbreak host, all nineteen fields in each, the full implications from P to GdFrameEx and to HostFrameEx, the auxiliary IsHostFrame and exists\_\allowbreak sigma, all Marks fields, and full exists\_\allowbreak nrep3 at the original local Noetherian-domain scope. The whole finEnd\_\allowbreak of\_\allowbreak finPt implication is complete, with G\_\allowbreak L fixed before every coefficient algebra and point. The remaining producers are the original P where not already supplied, actual FinPt on this frame, AutPt/AutEnd3/AutEnd2 and the actual rank-two and rank-three automorphic seeds, followed by full ModUAlt/PfromAut and local/automorphic descent to GdEnd. UCE already supplies the complete GdEnd consumer. The alternative GdEndAlt route retains its additional local/descent fields. Do not reopen Gen entries, the common-frame join or the FinPt-to-FinEnd transport; do not replace the Host ell=3 scope by the stronger Dat guard. Complete the original F3 local coefficient bridges and F4 common-overfield readings, with valuation/norm/GoodCoef/iq and the entire coefficient field, and actual complex Hecke standard/no-Jordan plus arithmetic character data on the same form datum. Construct the actual base rows, host/frame lengths, nonzero chains and group ties, MultEq support, eigensystem/carrier/parts, all-embedding logLam, HD/d3b and NEb. Keep the stronger second DivRows3 scope distinct from the completed first DivDef3 and row-pair integrability. The full arithmetic entry and conditional common-frame construction do not by themselves install HOST\_\allowbreak NUMBER\_\allowbreak S, HOST\_\allowbreak FRAME\_\allowbreak LENGTH, TASK\_\allowbreak HOST\_\allowbreak VAN\_\allowbreak E\_\allowbreak INSTALL, TASK\_\allowbreak HOST\_\allowbreak XPP\_\allowbreak INSTALL or TASK\_\allowbreak COUNT\_\allowbreak V\_\allowbreak INSTALL.

\emph{Source.} \href{https://github.com/CocoSgt/riemann-hypothesis-lean/blob/7cbf7a0c32fca1badf05931879a37c3f8a2b03b4/src/Development/AutomorphicArithmetic/AutomorphicVarieties/HeightAndVanishing/Counts.lean#L69-L69}{CountV}.

\emph{Source.} \href{https://github.com/CocoSgt/riemann-hypothesis-lean/blob/7cbf7a0c32fca1badf05931879a37c3f8a2b03b4/src/Development/AutomorphicArithmetic/AutomorphicVarieties/ModularityArgument/ThresholdArguments/DenseThresholdCases/CountVOf.lean#L33-L60}{UpperTheory.V.UnitaryThreeInertia.countV\_\allowbreak of\_\allowbreak finiteness\_\allowbreak duality\_\allowbreak euler\_\allowbreak host}.



\paragraph{14. THM\_\allowbreak S6\_\allowbreak S — C — Original unitary-three ordinary support and automorphy}

Original host/frame and fixed iota -> one marked g1 -> every universal R,u -> weight-zero finiteness and every Type-I BPrime\_\allowbreak s point.

\emph{Source.} \href{https://github.com/CocoSgt/riemann-hypothesis-lean/blob/7cbf7a0c32fca1badf05931879a37c3f8a2b03b4/src/Development/AutomorphicArithmetic/AutomorphicVarieties/ModularityArgument/ThresholdArguments/DenseThresholdCases/ThmS6CoreS.lean#L17-L26}{ThmS6\_\allowbreak s}.

\emph{Source.} \href{https://github.com/CocoSgt/riemann-hypothesis-lean/blob/7cbf7a0c32fca1badf05931879a37c3f8a2b03b4/src/Development/AutomorphicArithmetic/AutomorphicVarieties/ModularityArgument/ThresholdArguments/DenseThresholdCases/ThmS6CoreS.lean#L12-L16}{BPrime\_\allowbreak s}.



\paragraph{15. HOST\_\allowbreak FRAME\_\allowbreak LENGTH — B — Original Case-V host/frame and normalized-length installation}

Actual HostData, frame, local-condition equality, abstract count and normalized original Selmer length must all refer to the same objects.

\emph{Completed inputs (version 0.52).} AFE/LFC full entry and conditional common frame, CAF and MFC, DNE weight/host-degree dictionary at supplied tables

\emph{Remaining producer scope.} Complete the original F3 local coefficient bridges and F4 common-overfield readings, with valuation/norm/GoodCoef/iq and the entire coefficient field, and actual complex Hecke standard/no-Jordan plus arithmetic character data on the same form datum. Construct the actual base rows, host/frame lengths, nonzero chains and group ties, MultEq support, eigensystem/carrier/parts, all-embedding logLam, HD/d3b and NEb. Keep the stronger second DivRows3 scope distinct from the completed first DivDef3 and row-pair integrability. The full arithmetic entry and conditional common-frame construction do not by themselves install HOST\_\allowbreak NUMBER\_\allowbreak S, HOST\_\allowbreak FRAME\_\allowbreak LENGTH, TASK\_\allowbreak HOST\_\allowbreak VAN\_\allowbreak E\_\allowbreak INSTALL, TASK\_\allowbreak HOST\_\allowbreak XPP\_\allowbreak INSTALL or TASK\_\allowbreak COUNT\_\allowbreak V\_\allowbreak INSTALL.

\emph{Source.} \href{https://github.com/CocoSgt/riemann-hypothesis-lean/blob/7cbf7a0c32fca1badf05931879a37c3f8a2b03b4/src/Development/AutomorphicArithmetic/AutomorphicVarieties/ModularityArgument/ThresholdArguments/DenseThresholdCases/CountVOf.lean#L33-L60}{UpperTheory.V.UnitaryThreeInertia.countV\_\allowbreak of\_\allowbreak finiteness\_\allowbreak duality\_\allowbreak euler\_\allowbreak host}.



\paragraph{16. HOST\_\allowbreak NUMBER\_\allowbreak S — C — Uniform original-host rank-three catalogue}

One exponent before all closures, coefficient rings, levels, characters, frames, embeddings and marked points; one separating trace function controls all later finite pairwise nonisomorphic collections.

\emph{Completed inputs (version 0.52).} AFE full Host entry at original scope, LFC full conditional common frame, CAF and MFC, Inherited host-count consumers

\emph{Remaining producer scope.} Complete the original F3 local coefficient bridges and F4 common-overfield readings, with valuation/norm/GoodCoef/iq and the entire coefficient field, and actual complex Hecke standard/no-Jordan plus arithmetic character data on the same form datum. Construct the actual base rows, host/frame lengths, nonzero chains and group ties, MultEq support, eigensystem/carrier/parts, all-embedding logLam, HD/d3b and NEb. Keep the stronger second DivRows3 scope distinct from the completed first DivDef3 and row-pair integrability. The full arithmetic entry and conditional common-frame construction do not by themselves install HOST\_\allowbreak NUMBER\_\allowbreak S, HOST\_\allowbreak FRAME\_\allowbreak LENGTH, TASK\_\allowbreak HOST\_\allowbreak VAN\_\allowbreak E\_\allowbreak INSTALL, TASK\_\allowbreak HOST\_\allowbreak XPP\_\allowbreak INSTALL or TASK\_\allowbreak COUNT\_\allowbreak V\_\allowbreak INSTALL.

\emph{Source.} \href{https://github.com/CocoSgt/riemann-hypothesis-lean/blob/7cbf7a0c32fca1badf05931879a37c3f8a2b03b4/src/Development/AutomorphicArithmetic/AutomorphicVarieties/ModularityArgument/ThresholdArguments/DenseThresholdCases/HostNumberS.lean#L18-L28}{HostNumber\_\allowbreak s}.



\paragraph{17. TASK\_\allowbreak HOST\_\allowbreak VAN\_\allowbreak E\_\allowbreak INSTALL — B — Transport K/duality to literal HostVanE}

Use the already supplied LemmaK/local-global duality and finiteness on the actual CountV host/frame and exact original coefficient modules.

\emph{Source.} \href{https://github.com/CocoSgt/riemann-hypothesis-lean/blob/7cbf7a0c32fca1badf05931879a37c3f8a2b03b4/src/Development/AutomorphicArithmetic/AutomorphicVarieties/UnitaryThreeDeformations/HostNumber.lean#L9-L12}{HostVanE}.

\emph{Source.} \href{https://github.com/CocoSgt/riemann-hypothesis-lean/blob/7cbf7a0c32fca1badf05931879a37c3f8a2b03b4/src/Development/AutomorphicArithmetic/AutomorphicVarieties/LevelCohomology/SoHostVanE.lean#L11-L18}{UpperTheory.V.X.So.hostVanE}.



\paragraph{18. TASK\_\allowbreak HOST\_\allowbreak XPP\_\allowbreak INSTALL — B — Transport CorX to literal HostXpp}

Preserve Psi\textasciicircum{}-2 eta, the prime-to-ell quotient, actual local conditions and the original CountV host while installing BR7.

\emph{Source.} \href{https://github.com/CocoSgt/riemann-hypothesis-lean/blob/7cbf7a0c32fca1badf05931879a37c3f8a2b03b4/src/Development/AutomorphicArithmetic/AutomorphicVarieties/UnitaryThreeDeformations/HostNumber.lean#L14-L17}{HostXpp}.

\emph{Source.} \href{https://github.com/CocoSgt/riemann-hypothesis-lean/blob/7cbf7a0c32fca1badf05931879a37c3f8a2b03b4/src/Development/AutomorphicArithmetic/AutomorphicVarieties/LevelCohomology/Xpp.lean#L11-L84}{UpperTheory.V.X.xpp}.



\paragraph{19. DESIGN\_\allowbreak CHOICES\_\allowbreak COMMON — C — Construct the original common design family and bundle}

For the one target HostExW, construct the original compatible DesignChoices on Cdes=WithTw(WithSm Cprime), where Cprime=WithChN(A intersect target class), A is the one Built construction class, and cT(eps,y)=cTmin(y). Retain every original field, tuple, word, metric and common bundle duty. DSZ10 and RMP11 give size/rank properties of class members, not member existence. The original RecExP/RecAtG and common object/row route remains coupled. Construct ONE original common design family R and its recursive/geometric half with original Inputs/Pins/Holds/ExAt/InGeo and required CoreP. PKT now reconstructs full exists\_\allowbreak pkT, rowPk\_\allowbreak nonempty and exists\_\allowbreak moved with Copy/Coh on those original given inputs and a supplied fitting later stage: one catalogue, both ChT, complete LfT, actual integral lattices and all metric/measure/period constants are retained; one numerical record precedes every eta and V/S/N/G. RDR supplies one Side0 before all readers, ReadPk/pins and full ReadL with exact o.eta=q.eta. The initial actual geometry, Fits/division/central conditions, primitive rational tests and eta agreement must still be produced. Inherited CSR supplies finite columns, one ReadC before later twists, period/evaluation and scalar/error transports, RowSideDetermined and same-R Bundle/BTw/BWit/BunU consumers. RPI supplies exact RowPairInt for actual pairs and finite translates with original mu\_\allowbreak b/fd0; first literal DivDef3 at the same updated R,F,Y,Z,C,sEE follows from its retained original readings, Pins/Holds, stored nonzero IsItW chain, five real-group predicates and group tie. Remaining common fields include host/frame, MultEq nonzero-part support, that original nonzero chain/group tie, eigensystem/carrier/parts and all-embedding logLam, HD/d3b, NEb, actual rational tests and base rows. Retain every simultaneous family hypothesis and eventual threshold. Stronger second DivRows3 remains separate; universal quotient-height CuspRapidDecay and actual common member existence are not inferred from the comparison constructors. Version 0.41 supplies the finite phase-correct rational theta route: exact coefficient extension, split representatives, square-law/PhiIs, Pins from CoversAt, class-sum/point-identity transfer to OneMapOverL, actual-good-projector RatG and both-level LiftTestB. Its original higher covers, all-tuple class-sum uniqueness including phase rescaling, one PointIdentCore before O, base ClassSumsSpan and geometric LatFull remain explicit producers. MLat is available from the original LatFull input. This sharpens the rational-test work without claiming actual common geometry, Fits or eta agreement. Version0.46 proves finite generation of the original whole integral cusp lattice by actual integral-frame evaluation. It also proves the exact identity-cover, analytic unfolding, original-field coefficient descent, global clearing-scalar and both-level lattice-transport implications. Full spanning retains the supplied original IsNab connection and the five explicit principal-chart, level-system, raw-Gaussian, short-rule/dictionary and ramified-jet producers. The higher tuple covers, actual common class sums, PointIdentCore and ClassSumsSpan remain separate. CURRENT VERSION0.47 QUALIFICATION: Construct one original common geometric family and the higher CoversAt/all-tuple class sums and uniqueness with original phase rescalings, one PointIdentCore with common lambda,z, and base ClassSumsSpan. CNC now supplies the full original IsNab; PCL supplies the conditional trace-dual principal-lattice and quantitative normalization. Actual uniformization/HasPrinc and the five principal-chart, finite-level-system, Gaussian/reproduction, short-rule/simultaneous-dictionary and ramified-jet producer groups remain. Keep all original both-level, coefficient-field, Fits and eta agreements. Completed new original producers: ICL1-ICL21, CNC1-CNC16, PCL1-PCL12. Any earlier narrative naming these completed producers as outstanding is historical; this qualification is the current scope authority. CURRENT VERSION0.48 QUALIFICATION: Construct the original finite holomorphic FamilyUnif/N1c, normalized theta families and local lifting, then the common principal charts and finite-level theta systems. PUN now supplies the exact fixed-space models, all-family center lattices, A=1 exponent matching and individual primitive points with the full point-equality criterion; CNC supplies the full canonical connection and PCL the quantitative trace-dual normalization. Complete higher CoversAt and all-tuple class sums/uniqueness with the original phase rescalings, one PointIdentCore with common lambda,z, and base ClassSumsSpan. The Gaussian/reproduction, short-rule/simultaneous-dictionary and ramified-jet producers, and all original both-level, coefficient-field, Fits and eta agreements remain explicit. Completed current producers or conditional consumers: ICL1-ICL21, CNC1-CNC16, PCL1-PCL12, PUN1-PUN24, GCF1-GCF18 at their stated scopes. Earlier frontier descriptions are historical; this qualification is the current scope authority. Original whole-contract premises and quantifiers are retained. CURRENT VERSION0.49 QUALIFICATION: Construct the adjoining simultaneous principal framed points and test-ring dictionary. Full finite holomorphic FamilyUnif/N1c, normalized theta families and local lifts are now supplied by HUF. The next bounded target is allTest\_\allowbreak of\_\allowbreak secBaseChange: represent FrPt by the Lie-eigenframe bundles and theta-section frame line, use the full n=1 SecBaseChangeSt on every nonflat coefficient base change, and prove finite-generated subalgebra descent with the exact TRing/FromTest transport. Then supply the full FramesOK/CoverTri/FibTri and IsUnf/hasPrinc package, including adapted mod-eight symplectic basis, normalized theta multiplier/weight and arithmetic invariance. Continue higher CoversAt, all-tuple class sums and uniqueness with original phase rescalings, one PointIdentCore with common lambda,z, base ClassSumsSpan, Gaussian/reproduction, short-rule dictionaries and ramified jets, retaining both-level, coefficient-field, Fits and eta agreements. Common-family nonemptiness remains distinct from compatibility of an already supplied member. Completed current producers or conditional consumers: ICL1-ICL21, CNC1-CNC16, PCL1-PCL12, PUN1-PUN24, HUF1-HUF21 (full original FamilyUnif/N1c), GCF1-GCF18 and WRS1-WRS10 at their stated scopes. Earlier frontier descriptions are historical; this qualification is the current scope authority. Original whole-contract premises and quantifiers are retained.

\emph{Completed inputs (version 0.52).} GSL1-GSL26 (full conditional section laws, universal basis and projective-model algebra), NFL1-NFL22, TDR1-TDR13, NFR1-NFR13, HUF1-HUF21, PUN1-PUN24, ICL1-ICL21, CNC1-CNC16, PCL1-PCL12, GCF/WRS at original supplied-member scope

\emph{Remaining producer scope.} GSL proves the full conditional SecBaseChangeAny/St and SecTwistSt at their original scopes, with one actual basis before every affine test and every residue section. Its unconditional algebra gives the universal kernel certificate, arbitrary tensor transfer, finite-projective models from the supplied Noetherian section cube, and all Perf coordinates. The remaining geometric suppliers are the actual Noetherian section-cube model and finite cohomology/smooth descent, PolBaseChangeSt, Serre and multiplication/invariant descent, geometric integrality, principal rank one and section decomposition/growth, and the theorem-of-square generators. These must hold on the same polarized abelian scheme before applying GSL26 to the completed TDR/NFR/NFL uniformization. The supplied Moduli and inherited HUF/PUN foundations remain explicit. Then construct the finite level system, Gaussian/reproduction, short-rule and simultaneous dictionaries, ramified jets, higher CoversAt, all-tuple class sums and uniqueness with phases, one PointIdentCore with common lambda and z, and base ClassSumsSpan. Preserve both levels, the entire coefficient field, Fits and eta agreements, and the common-family quantifier order. Common design nonemptiness remains required.



\paragraph{20. LOCAL\_\allowbreak NBUILT\_\allowbreak COMMON — C — Actual local specialization on the common selected family}

Construct, with the required closed-class property on the same Des(C,3), one eventual family of actual NBuilt certificates for the original Fm.D(x,i), retaining the certificate before forgetting to TSel. HLP now reconstructs full HeckeLift0, all-index commutation, HeckeStd0 and HeckeIntTw from the original MtInput and complex Hecke facts, in the full ambient GenAt range k<=nV. Every maximal ideal is an actual supported residual character class; literal MtSub is the O-span of projections of the original MtLat generators, is finite free, and localizes onto the exact NCov.part. Literal Psi has its exact range and off-space projection error, both original span assertions hold, and ActEx/ActHK are constructed on the same lattice under original semisimplicity. Original MjLat relation quotients, universal coefficient-map relation fields, selected tame generators, coefficient embeddings and all identity transports give the complete nonempty\_\allowbreak spec constructor on its full printed input list. A single compositum combines the four original coefficient requirements with big fixed beforehand. Whole ndSpec\_\allowbreak exists still needs actual common geometric/choice MtInput, complex HeckeStdPlain and the no-Jordan producer, original F3 bridges with valuation/GoodCoef/iq, original F4 common-overfield reading of the entire coefficient field, and exact arithmetic character/reciprocity/admissibility suppliers. Original AutSpec, Galois specialization and the remaining NBuilt fields on that same installation also remain. Do not replace the actual projected lattice or maps by abstract objects with the same dimensions.

\emph{Completed inputs (version 0.52).} GSL/NFL full conditional section laws and uniformization, AFE/LFC full entry and conditional common frame, TDR/NFR and CAF/MFC, WRS selected built-class extraction at membership

\emph{Remaining producer scope.} Continue the distinct retained-defect, period-height and normalized spectral producers, and the original compatible second-layer selection on one common arithmetic/geometric family. GLOBAL\_\allowbreak LOCAL\_\allowbreak GAIN remains fully closed at its original v0.49 scope. Complete local comparison theorems do not themselves construct the whole common family. Common design existence, NBUILT compatibility, original true-column records, shared fields/embeddings, Fits and eta agreement, simultaneous rows and eventual-family thresholds remain explicit. A stronger historical WPt/TSel construction may serve other producers but is not needed to reopen the proved selected-class local-gain estimate. GSL proves the full conditional SecBaseChangeAny/St and SecTwistSt at their original scopes, with one actual basis before every affine test and every residue section. Its unconditional algebra gives the universal kernel certificate, arbitrary tensor transfer, finite-projective models from the supplied Noetherian section cube, and all Perf coordinates. The remaining geometric suppliers are the actual Noetherian section-cube model and finite cohomology/smooth descent, PolBaseChangeSt, Serre and multiplication/invariant descent, geometric integrality, principal rank one and section decomposition/growth, and the theorem-of-square generators. These must hold on the same polarized abelian scheme before applying GSL26 to the completed TDR/NFR/NFL uniformization. The supplied Moduli and inherited HUF/PUN foundations remain explicit. Then construct the finite level system, Gaussian/reproduction, short-rule and simultaneous dictionaries, ramified jets, higher CoversAt, all-tuple class sums and uniqueness with phases, one PointIdentCore with common lambda and z, and base ClassSumsSpan. Preserve both levels, the entire coefficient field, Fits and eta agreements, and the common-family quantifier order.



\paragraph{21. SECOND\_\allowbreak LAYER\_\allowbreak COMPATIBLE\_\allowbreak SELECTION — C — Original second-layer selector and class duties}

For the same A, supply Duty on A and a Closed common original TSel family whose selected records are the original true-column rec7 records (or have the entire dependent tuple transported), with their real witness and original coefficient-field root data. Existence/compatibility remains arithmetic; merely abstractly isomorphic fields or independently selected columns do not suffice. CSR21–22 now construct the same-R Bundle/BTw/BWit/BunU outputs from the exact selected realization/row-side packet and NEb, with the original EvP/ParEv quantifier order. Actual common selection, the five stated realization packets and NEb remain upstream; none is replaced by an independently chosen compatible record. Version 0.38 closes the selected analytic extension and reconstructs the first updated DivDef3 consumer; it leaves the original common selection, geometric readings, rational tests, base rows, host/parts/NEb and exact retained DI premises upstream, with unchanged EvP/ParEv thresholds. Version 0.39 supplies the original given-data reader constructors and exact RowPairInt/first DivDef3 on original premises. It leaves one actual common geometric/comparison/test member, eta agreement, compatible selection, rational tests, base rows, host/parts/NEb and original EvP/ParEv thresholds upstream. Version 0.40 also supplies original comparison/RowPk/fitting-stage Moved algorithms and exact common data transport. The actual initial geometric member, fitting tests, eta agreement and compatible second-layer selector with original eventual thresholds remain upstream. Version 0.41 supplies the finite phase-correct rational theta route: exact coefficient extension, split representatives, square-law/PhiIs, Pins from CoversAt, class-sum/point-identity transfer to OneMapOverL, actual-good-projector RatG and both-level LiftTestB. Its original higher covers, all-tuple class-sum uniqueness including phase rescaling, one PointIdentCore before O, base ClassSumsSpan and geometric LatFull remain explicit producers. MLat is available from the original LatFull input. This sharpens the rational-test work without claiming actual common geometry, Fits or eta agreement. CURRENT VERSION0.49 QUALIFICATION: Continue the distinct global denominator, period/height, normalized spectral and second-layer compatible-selection producers, together with existence of the original common arithmetic/geometric family. GLOBAL\_\allowbreak LOCAL\_\allowbreak GAIN and its actual GroupCompatibilityFamily supplier are now complete on the unchanged selected built class; no additional WPt/TSel realization is needed for that estimate. Do not reopen local first-growth, mean, auxiliary admission, selected D3 tops/Adm2F, the single GIs2 stage or the exact comparison/kernel transport. A full RecT or TSel may remain useful for other tasks, but is not claimed by this local-gain closure. Completed current producers or conditional consumers: GCF1-GCF18 at their stated scopes, WRS1-WRS10 original selected-class compatibility/local gain. Earlier frontier descriptions are historical; this qualification is the current scope authority. Original whole-contract premises and quantifiers are retained.

\emph{Completed inputs (version 0.52).} WRS1-WRS10 original selected-class compatibility/local gain, DNE full denominator consumer, DFD/SSC canonical-data and signature comparisons, GSL/NFL and AFE/LFC conditional constructors, Inherited exact same-record realization/transport consumers

\emph{Remaining producer scope.} Continue the distinct retained-defect, period-height and normalized spectral producers, and the original compatible second-layer selection on one common arithmetic/geometric family. GLOBAL\_\allowbreak LOCAL\_\allowbreak GAIN remains fully closed at its original v0.49 scope. Complete local comparison theorems do not themselves construct the whole common family. Common design existence, NBUILT compatibility, original true-column records, shared fields/embeddings, Fits and eta agreement, simultaneous rows and eventual-family thresholds remain explicit. A stronger historical WPt/TSel construction may serve other producers but is not needed to reopen the proved selected-class local-gain estimate.



\paragraph{22. BASE\_\allowbreak THETA\_\allowbreak READINGS — C — Actual base reading on the common bundle}

For every family in the original class, eventually on the SAME selected Bundle and for each i<m, supply the actual R,F,Read,ReadL, weak models/Weil data and phase-one pins. Identify their LiftTie.theta, mtRead, mjRead and ThDef3 basis with that Bundle.P.A3, P.Zs, P.S0 and its stored maps. Exact field, embeddings, metrics, measures, row family and weak choices are retained. This remains one part of the coupled BTH12 producer. PKT now supplies the complete original comparison packet, RowPk and fitting-stage Moved with Copy/Coh on original given geometry, Inputs/Pins/Holds/ExAt/InGeo and Fits. Its ChT/LfT/Reps, actual lattice carriers, original Haar and norm constants, period maps, representatives and all-place RrT are constructed together; its one numerical record precedes every eta/V/S/N/G. RDR gives the original level-zero side, moved-Read and ReadPk/pin constructors on those objects, with one Side0 before every reader and all twelve Read ties. Full ReadL retains the exact o.eta(j,v)=q.eta(j,v) guard for every 1<=j<=m and place v; arbitrary Over does not imply it. Inherited CSR fixes one ReadC before all TwArgs, transports both metrics and all thirteen reading maps, and gives actual period/evaluation with kappa=(c*nP)\textasciicircum{}(-1/2). RPI supplies exact row-pair integrability, so first updated DivDef3 follows from the retained original Pins/Holds, nonzero chain, five real-group predicates and group tie. Initial theta extension remains complete. Remaining work is actual initial common geometry and its original predicates, Fits/primitive rational tests, eta agreement and outer family/eventual/same-Bundle duties, reFace compatibility, MultEq support, original chain/group data, actual base rows and any stronger DivRows3 theta square. No new compatible reading, family threshold or Haar normalization is substituted. Version 0.41 supplies the finite phase-correct rational theta route: exact coefficient extension, split representatives, square-law/PhiIs, Pins from CoversAt, class-sum/point-identity transfer to OneMapOverL, actual-good-projector RatG and both-level LiftTestB. Its original higher covers, all-tuple class-sum uniqueness including phase rescaling, one PointIdentCore before O, base ClassSumsSpan and geometric LatFull remain explicit producers. MLat is available from the original LatFull input. This sharpens the rational-test work without claiming actual common geometry, Fits or eta agreement.

\emph{Completed inputs (version 0.52).} GSL1-GSL26, NFL1-NFL22, TDR1-TDR13, NFR1-NFR13, HUF1-HUF21, PUN1-PUN24, CNC1-CNC16, Inherited actual-reading consumers at original inputs

\emph{Remaining producer scope.} GSL proves the full conditional SecBaseChangeAny/St and SecTwistSt at their original scopes, with one actual basis before every affine test and every residue section. Its unconditional algebra gives the universal kernel certificate, arbitrary tensor transfer, finite-projective models from the supplied Noetherian section cube, and all Perf coordinates. The remaining geometric suppliers are the actual Noetherian section-cube model and finite cohomology/smooth descent, PolBaseChangeSt, Serre and multiplication/invariant descent, geometric integrality, principal rank one and section decomposition/growth, and the theorem-of-square generators. These must hold on the same polarized abelian scheme before applying GSL26 to the completed TDR/NFR/NFL uniformization. The supplied Moduli and inherited HUF/PUN foundations remain explicit. Then construct the finite level system, Gaussian/reproduction, short-rule and simultaneous dictionaries, ramified jets, higher CoversAt, all-tuple class sums and uniqueness with phases, one PointIdentCore with common lambda and z, and base ClassSumsSpan. Preserve both levels, the entire coefficient field, Fits and eta agreements, and the common-family quantifier order. Complete the original F3 local coefficient bridges and F4 common-overfield readings, with valuation/norm/GoodCoef/iq and the entire coefficient field, and actual complex Hecke standard/no-Jordan plus arithmetic character data on the same form datum. Construct the actual base rows, host/frame lengths, nonzero chains and group ties, MultEq support, eigensystem/carrier/parts, all-embedding logLam, HD/d3b and NEb. Keep the stronger second DivRows3 scope distinct from the completed first DivDef3 and row-pair integrability. The full arithmetic entry and conditional common-frame construction do not by themselves install HOST\_\allowbreak NUMBER\_\allowbreak S, HOST\_\allowbreak FRAME\_\allowbreak LENGTH, TASK\_\allowbreak HOST\_\allowbreak VAN\_\allowbreak E\_\allowbreak INSTALL, TASK\_\allowbreak HOST\_\allowbreak XPP\_\allowbreak INSTALL or TASK\_\allowbreak COUNT\_\allowbreak V\_\allowbreak INSTALL.



\paragraph{23. BASE\_\allowbreak THETA\_\allowbreak RATIONAL\_\allowbreak TESTS — C — Simultaneous original base rational tests}

On the SAME original Bundle/reading/theta map, for every i<m supply one finite separating complex-linear test family whose values lie in the fixed embedded coefficient field on theta of EVERY original rational source vector and on EVERY stored target-coordinate basis vector. BTH8/BTH11 use this same LiftTestB for rational adjoint image and matrix coordinates. RTG now reconstructs prescribed-field coefficient extension, original split class representatives/RepsOK, the self-dual split-lattice square law and finite residue PhiIs; every Pins field follows from the actual CoversAt. With the base coefficient record literally fixed, exact class-sum existence/uniqueness including phase-rescaled O, PointIdentCore with common lambda,z before every O, and base ClassSumsSpan imply all-tuple finite lift span and OneMapOverL. The phase is exactly cj(twAt embLL A3.A.v0)(R.phs i). RatG uses the actual good projector eOp(y), and full LiftTestB follows at BOTH original levels. These finite implication proofs are complete. Remaining original producers are the prescribed higher covers, those actual class sums/uniqueness, PointIdentCore, base ClassSumsSpan and the base/common geometric datum; exact MLat is supplied only from original LatFull, whose geometric production remains separate. Finite dimension or an unrelated rational theta model cannot replace these same-object identities. Original outer eventual-family/Bundle quantifiers remain. Version0.46 proves finite generation of the original whole integral cusp lattice by actual integral-frame evaluation. It also proves the exact identity-cover, analytic unfolding, original-field coefficient descent, global clearing-scalar and both-level lattice-transport implications. Full spanning retains the supplied original IsNab connection and the five explicit principal-chart, level-system, raw-Gaussian, short-rule/dictionary and ramified-jet producers. The higher tuple covers, actual common class sums, PointIdentCore and ClassSumsSpan remain separate. CURRENT VERSION0.47 QUALIFICATION: Construct one original common geometric family and the higher CoversAt/all-tuple class sums and uniqueness with original phase rescalings, one PointIdentCore with common lambda,z, and base ClassSumsSpan. CNC now supplies the full original IsNab; PCL supplies the conditional trace-dual principal-lattice and quantitative normalization. Actual uniformization/HasPrinc and the five principal-chart, finite-level-system, Gaussian/reproduction, short-rule/simultaneous-dictionary and ramified-jet producer groups remain. Keep all original both-level, coefficient-field, Fits and eta agreements. Completed new original producers: ICL1-ICL21, CNC1-CNC16, PCL1-PCL12. Any earlier narrative naming these completed producers as outstanding is historical; this qualification is the current scope authority. CURRENT VERSION0.48 QUALIFICATION: Construct the original finite holomorphic FamilyUnif/N1c, normalized theta families and local lifting, then the common principal charts and finite-level theta systems. PUN now supplies the exact fixed-space models, all-family center lattices, A=1 exponent matching and individual primitive points with the full point-equality criterion; CNC supplies the full canonical connection and PCL the quantitative trace-dual normalization. Complete higher CoversAt and all-tuple class sums/uniqueness with the original phase rescalings, one PointIdentCore with common lambda,z, and base ClassSumsSpan. The Gaussian/reproduction, short-rule/simultaneous-dictionary and ramified-jet producers, and all original both-level, coefficient-field, Fits and eta agreements remain explicit. Completed current producers or conditional consumers: ICL1-ICL21, CNC1-CNC16, PCL1-PCL12, PUN1-PUN24. Earlier frontier descriptions are historical; this qualification is the current scope authority. Original whole-contract premises and quantifiers are retained. CURRENT VERSION0.49 QUALIFICATION: Construct the adjoining simultaneous principal framed points and test-ring dictionary. Full finite holomorphic FamilyUnif/N1c, normalized theta families and local lifts are now supplied by HUF. The next bounded target is allTest\_\allowbreak of\_\allowbreak secBaseChange: represent FrPt by the Lie-eigenframe bundles and theta-section frame line, use the full n=1 SecBaseChangeSt on every nonflat coefficient base change, and prove finite-generated subalgebra descent with the exact TRing/FromTest transport. Then supply the full FramesOK/CoverTri/FibTri and IsUnf/hasPrinc package, including adapted mod-eight symplectic basis, normalized theta multiplier/weight and arithmetic invariance. Continue higher CoversAt, all-tuple class sums and uniqueness with original phase rescalings, one PointIdentCore with common lambda,z, base ClassSumsSpan, Gaussian/reproduction, short-rule dictionaries and ramified jets, retaining both-level, coefficient-field, Fits and eta agreements. Completed current producers or conditional consumers: ICL1-ICL21, CNC1-CNC16, PCL1-PCL12, PUN1-PUN24, HUF1-HUF21 (full original FamilyUnif/N1c and theta family). Earlier frontier descriptions are historical; this qualification is the current scope authority. Original whole-contract premises and quantifiers are retained.

\emph{Completed inputs (version 0.52).} GSL1-GSL26, NFL1-NFL22, TDR1-TDR13, NFR1-NFR13, HUF1-HUF21, PUN1-PUN24, ICL1-ICL21, CNC1-CNC16, PCL1-PCL12, Inherited both-level RatG/LiftTestB consumers

\emph{Remaining producer scope.} GSL proves the full conditional SecBaseChangeAny/St and SecTwistSt at their original scopes, with one actual basis before every affine test and every residue section. Its unconditional algebra gives the universal kernel certificate, arbitrary tensor transfer, finite-projective models from the supplied Noetherian section cube, and all Perf coordinates. The remaining geometric suppliers are the actual Noetherian section-cube model and finite cohomology/smooth descent, PolBaseChangeSt, Serre and multiplication/invariant descent, geometric integrality, principal rank one and section decomposition/growth, and the theorem-of-square generators. These must hold on the same polarized abelian scheme before applying GSL26 to the completed TDR/NFR/NFL uniformization. The supplied Moduli and inherited HUF/PUN foundations remain explicit. Then construct the finite level system, Gaussian/reproduction, short-rule and simultaneous dictionaries, ramified jets, higher CoversAt, all-tuple class sums and uniqueness with phases, one PointIdentCore with common lambda and z, and base ClassSumsSpan. Preserve both levels, the entire coefficient field, Fits and eta agreements, and the common-family quantifier order.



\paragraph{24. SLR\_\allowbreak BASE\_\allowbreak ROWS\_\allowbreak PRODUCER — C — Actual base row identities on that same bundle}

For the SAME selected Bundle, common original fields and readings used by the other three BTH12 inputs, supply actual RowsHold0\_\allowbreak 3 uniformly eventually for every family. Its original primitive row realization, spectrum/type and basis identifications must be retained. BTH supplies a theta consumer only; it does not construct the base row identity. No separate witness may be chosen merely to satisfy this row requirement.



\clearpage
\section{Universal section bases and the two global section laws}
\label{sec:global-section-laws-v052}

\status{GSL1--GSL26 reconstruct the finite-complex and geometric
arguments supplying the original section laws.  The conclusions concern
arbitrary coefficient rings and the actual line-bundle sections, including
every residue-section lift in \texttt{Sec\allowbreak Twist\allowbreak St}.  The precise upstream
geometric premises are recorded below; they are not silently promoted to
proved global laws.  All source references use the accepted formal proof
at commit \texttt{7cbf7a0c32fca1badf05931879a37c3f8a2b03b4}.}

The purpose of this section is to remove the linear-algebra and
section-transport work from the premises retained in NFL22 and TDR13.
Its central construction is an identity of matrices over the original
ring.  Consequently, neither the coefficient ring nor a later base
change is required to be Noetherian, reduced or flat.  The final
statements preserve the original polarized abelian scheme, its line
bundle, the chosen point and every subsequent affine test map.

Here and below, a line bundle $M$ is the source's actual cocycle
\texttt{LB}: an open cover $U_i$, units $g_{ij}$ on intersections,
and $g_{ij}g_{jk}=g_{ik}$.  A section is an actual family $s_i$ with
$s_i=g_{ij}s_j$.  We write $\Gamma(X,M)$ for this module; this is
\texttt{LB.\allowbreak Sec} with its module structure, or \texttt{Sec\allowbreak M} after
the specified \texttt{LPow} chart conversion.  We identify
$\Gamma(\operatorname{Spec}R,\mathcal O)$ with $R$ by the canonical
ring isomorphism, and transport back through that isomorphism whenever
the original statement uses its literal coefficient ring.

For $f:X\to\operatorname{Spec}R$, let $c_f:R\to\Gamma(X,\mathcal O)$
be the structure map.  The source's $\operatorname{pw}(M,t)$ has
transition functions $g_{ij}^{t}$, for any integer $t$; thus negative
powers cause no difficulty.  The source's
$\operatorname{tw}(M,H,t)$ is $M\otimes H^t$, on the intersection
cover.  All isomorphisms used below are represented by actual units on
the appropriate common charts.  An equality of classes in a Picard group
alone is not used as a section comparison.

We abbreviate \texttt{First\allowbreak Cohomology\allowbreak Vanishes} by $\mathrm{V}(M)$.
Its meaning is the following explicit assertion, for \emph{every}
open cover $V_a\subseteq U_{i(a)}$: if
\[
 \phi_{ab}+g_{i(a)i(b)}\phi_{bd}=\phi_{ad}
 \quad\hbox{on }V_a\cap V_b\cap V_d,
\]
then $\phi_{ab}=\psi_a-g_{i(a)i(b)}\psi_b$ for suitable actual
functions $\psi_a\in\Gamma(V_a,\mathcal O)$.  This notation does
not introduce an extra cohomology comparison theorem.

\subsection{One kernel basis before every base change}

% BEGIN GSL1
\paragraph{GSL1: section coordinates, pullback and uniqueness.}
\label{claim:GSL1}
Compatible functions on any open refinement of a line-bundle cover
determine a unique section.  The construction is additive, commutes with
multiplication by a function from the base, and commutes with pullback.
For a finite linear combination of pulled line bundles, these assertions
hold for the literal \texttt{Lin\allowbreak Comb} charts and the source's
\texttt{Is\allowbreak Pull\allowbreak Sec\allowbreak LC}; for a single power they give \texttt{Is\allowbreak Pull\allowbreak Sec}.

\begin{proof}
Suppose $W_a\subseteq U_{i(a)}$ cover $X$ and the prescribed coordinates
are $v_a$.  On $W_a\cap U_j$, the coordinate in chart $j$ must be
$g_{j,i(a)}v_a$.  The cocycle identity shows that these functions agree
on overlaps.  The sheaf axiom glues them on $U_j$; the same identity gives
the section compatibility between $U_j$ and $U_k$.  This also proves
uniqueness, since equality on the $W_a$ forces equality after changing
to every other chart.  Addition and scalar multiplication commute with
restrictions and multiplication by transition units, so uniqueness gives
the asserted linearity.

If $\pi:X'\to X$, pull back every coordinate and every transition unit.
The functions still satisfy the same compatibility equation, now on
$\pi^{-1}U_i$.  Repeating a pullback gives the pullback by the composite,
because the maps of structure sheaves compose.  For
$\bigotimes_{k\in K}\phi_k^*L^{\epsilon_k}$ the chart indexed by
$j:K\to\operatorname{Charts}(L)$ is
$\bigcap_k\phi_k^{-1}U_{j(k)}$, and its transition is the product
$\prod_k\phi_k^*g_{j(k)j'(k)}^{\epsilon_k}$.  Pullback of this
product is the transition on the corresponding literal common chart.
The coordinate equality just established is exactly
\texttt{Is\allowbreak Pull\allowbreak Sec\allowbreak LC}.  Taking $K=\{1\}$ gives the power version.
The redundant one-element index in \texttt{LB.\allowbreak pull} is removed by the
mutually inverse chart maps $j\mapsto(\ast\mapsto j)$ and
$J\mapsto J(\ast)$, so no choice of a different section is involved.
\end{proof}
% END GSL1

% BEGIN GSL2
\paragraph{GSL2: a denominator and a left inverse over the original ring.}
\label{claim:GSL2}
Let $R$ be a commutative ring, $K$ a field, $\gamma:R\to K$ a ring
map, and $A:R^l\to R^m$ a matrix.  If $A_K$ is injective, there are
$r\in R$ and a matrix $T:R^m\to R^l$ such that
\[
 \gamma(r)\ne0,\qquad TA=rI_l. \tag{GSL.1}
\]

\begin{proof}
Choose $l$ rows whose restriction $N$ of $A_K$ has nonzero determinant;
they exist because $A_K$ has rank $l$.  Let $J:R^m\to R^l$ select
these same rows over $R$, set $N=JA$, $r=\det N$, and
$T=\operatorname{adj}(N)J$.  The adjugate identity over an arbitrary
commutative ring gives $TA=rI_l$.  Determinants commute with ring maps,
so $\gamma(r)=\det N_K\ne0$.  For $l=0$ use the empty determinant
$r=1$ and the unique empty matrix.  No cancellation in $R$ has been used.
\end{proof}
% END GSL2

% BEGIN GSL3
\paragraph{GSL3: the universal kernel certificate.}
\label{claim:GSL3}
Suppose
\[
 R^a\xrightarrow{D_0}R^b\xrightarrow{D_1}R^c,
 \qquad D_1D_0=0,
\]
is exact at $K^b$ after applying a ring map $\gamma:R\to K$ to a
field.  There are $n\ge0$, $r\in R$, and matrices
$B:R^n\to R^a$, $C:R^a\to R^n$, $H:R^b\to R^a$ with
\[
 \gamma(r)\ne0,\qquad
 D_0B=0,\qquad CB=rI_n,\qquad BC+HD_0=rI_a. \tag{GSL.2}
\]

\begin{proof}
Select columns of $(D_0)_K$ that form a basis of its image, and let
$S:R^k\to R^a$ be their coordinate inclusion.  GSL2 applied to $D_0S$
gives $r_0,T_0$ with $\gamma(r_0)\ne0$ and
$T_0D_0S=r_0I_k$.  The stacked matrix
\[
 A=\begin{pmatrix}T_0\\D_1\end{pmatrix}:R^b\longrightarrow R^{k+c}
\]
is injective over $K$.  Indeed, if $A_Kw=0$, exactness puts $w$ in
$\operatorname{im}(D_0)_K=\operatorname{im}(D_0S)_K$; write
$w=(D_0S)_Kv$.  Then $0=(T_0)_Kw=\gamma(r_0)v$, hence $v=w=0$.
Apply GSL2 again to obtain $T_1A=r_1I_b$, with $\gamma(r_1)\ne0$.

Set $Q=r_0I_a-ST_0D_0$.  Then
\[
 QS=0,\qquad T_0D_0Q=0,\qquad D_1D_0Q=0.
\]
Multiplying the last two equalities by $T_1$ gives $r_1D_0Q=0$;
this is why the second denominator is retained even when $R$ has
zero divisors.  Let $E:R^n\to R^a$ include the complementary
coordinates to those selected by $S$, where $n=a-k$.  Thus
$E^tE=I_n$, $E^tS=0$ and $SS^t+EE^t=I_a$.  Consequently
$QEE^t=Q$ and $E^tQE=r_0I_n$.  Define
\[
 r=r_1r_0,\qquad B=r_1QE,\qquad C=E^t,\qquad H=r_1ST_0.
\]
The first three identities in (GSL.2) now follow directly, and
\[
 BC+HD_0=r_1Q+r_1ST_0D_0=r_1r_0I_a.
\]
Both field images of $r_0,r_1$ are nonzero.  The construction also
covers $k=0$ and $n=0$, using empty matrices throughout.
\end{proof}
% END GSL3

% BEGIN GSL4
\paragraph{GSL4: a kernel basis after every permitted ring map.}
\label{claim:GSL4}
Choose $r,n,B,C,H$ as in GSL3, once and for all.  For every commutative
ring $R'$ and every $\varphi:R\to R'$ for which $\varphi(r)$ is a
unit, the map
\[
 R'^n\longrightarrow\ker(D_0)_{R'},\qquad z\longmapsto B_{R'}z
 \tag{GSL.3}
\]
is an isomorphism, where $\ker(D_0)_{R'}$ means the kernel of the
\emph{base-changed matrix}.  Its inverse on a vector $v$ in that kernel
is $\varphi(r)^{-1}C_{R'}v$.

If $R$ is local, $\kappa$ its residue field, $K/\kappa$ a field
extension and $\gamma:R\to\kappa\hookrightarrow K$, then the
same $r$ is already a unit of $R$.  The basis therefore works over $R$
and $\kappa$, and every vector in $\ker(D_0)_\kappa$ lifts to a
vector in $\ker D_0$.

\begin{proof}
All identities (GSL.2) remain true after applying $\varphi$.  Thus
$D_0B=0$ makes (GSL.3) well defined.  If $D_0v=0$, then
$BCv=\varphi(r)v$, which gives existence of the displayed coordinates.
The identity $CB=\varphi(r)I_n$ gives uniqueness.  These identities
also prove the assertion if $R'$ is the zero ring.  Nothing asserts
that tensoring an arbitrary kernel is exact; the identities prove the
needed conclusion for these particular matrices.

In the local case, $\gamma(r)\ne0$ implies that the residue of $r$
is nonzero, hence $r$ lies outside the maximal ideal and is a unit.
Express a residue vector uniquely as $\overline B\,\overline z$.
Lift its finitely many coordinates along the surjection
$R\to\kappa$, obtaining $z\in R^n$.  Then $Bz$ is a cycle over $R$
and reduces to the prescribed vector.  The basis is chosen before the
residue vector, not anew for each one.
\end{proof}
% END GSL4

\subsection{Finite complexes with the required universal coordinates}

% BEGIN GSL5
\paragraph{GSL5: replacing three finite projectives by three finite matrices.}
\label{claim:GSL5}
Let $P^0\xrightarrow{d_0}P^1\xrightarrow{d_1}P^2$ be a complex
of finite projective $R$-modules.  There are finite free modules
$F^0,F^1,F^2$, matrices $D_0,D_1$ with $D_1D_0=0$, and a fixed
map $p:F^0\to P^0$ such that, after \emph{every} ring map out of $R$,
$p$ induces an isomorphism on the degree-zero kernels.  Exactness at
the middle term of the projective complex implies exactness at the
middle term of the free complex.  The same assertions hold after
tensoring with an arbitrary $R$-module.

\begin{proof}
Choose split presentations
$p_i:E_i=R^{n_i}\to P^i$, $s_i:P^i\to E_i$ with $p_is_i=1$.
Put $F^0=E_0$, $F^1=E_1\oplus E_0$,
$F^2=E_2\oplus E_1\oplus E_0$, and set
\begin{align*}
 D_0(v)&=(s_1d_0p_0v,\ v-s_0p_0v),\\
 D_1(w_1,w_0)&=(s_2d_1p_1w_1,\ w_1-s_1p_1w_1,\ s_0p_0w_0).
\end{align*}
Each component of $D_1D_0$ vanishes, by $d_1d_0=0$ and the
splitting identities.  A vector in $\ker D_0$ satisfies
$v=s_0p_0v$ and $d_0p_0v=0$, so $p_0$ identifies this kernel with
$\ker d_0$, with inverse $s_0$.

Suppose $D_1(w_1,w_0)=0$.  Then $w_1=s_1p_1w_1$,
$d_1p_1w_1=0$, and $p_0w_0=0$ (apply $p_0$ to the third
component).  If the projective complex is exact, choose $x\in P^0$
with $d_0x=p_1w_1$.  The vector $v=s_0x+w_0$ satisfies
$D_0v=(w_1,w_0)$.  Conversely $D_1D_0=0$ gives the other inclusion.
All displayed maps and identities tensor to the same formulas; the
splittings survive arbitrary tensor products.  Finally, if $e_k$ are
the standard basis of $E_0$, the fixed elements $u_k=p_0e_k$
give the kernel comparison by
$v\mapsto\sum_kv_k\otimes u_k$ after base change.  These fixed
coordinates will become the \texttt{Perf.\allowbreak u} functions.
\end{proof}
% END GSL5

% BEGIN GSL6
\paragraph{GSL6: a finite complex over a Noetherian ring.}
\label{claim:GSL6}
Let $R$ be Noetherian and $C^\bullet$ a nonnegative cochain complex,
zero in degrees above $d$, whose cohomology modules are finite.
There exist a complex $K^\bullet$, also zero above $d$, and a
quasi-isomorphism $\phi:K^\bullet\to C^\bullet$, with every $K^p$
finite and every $K^p$ for $p\ge1$ finite free.  No finiteness of the
individual $C^p$ is assumed.

\begin{proof}
We give the finite lifting step underlying a descending construction.
Suppose finite $F^1$, maps $\phi_1:F^1\to C^1$ and
$\delta_1:F^1\to F^2$ have already been chosen.  In
$C^0\oplus F^1$ consider
\[
 N=\{(c,k):d_Cc=\phi_1k,\ \delta_1k=0\}.
\]
Its projection $J$ to $F^1$ is finite because $R$ is Noetherian.
Choose finitely many generators of $J$ and lift them to $N$.  Also
choose cycles in $C^0$ lifting generators of
$\ker d_C/(B\cap\ker d_C)$, where $B$ is the boundary submodule
at the current stage.  This quotient is finite by the assumed
cohomology finiteness.  Take a finite free module on both sets of
chosen lifts, with maps $\phi_0$ and $\delta_0$ given by their first
and second coordinates (the second coordinate of the cycle lifts is
zero).  Then
\[
 d_C\phi_0=\phi_1\delta_0,\quad
 \delta_1\delta_0=0,
\]
and for every $(c,k)\in N$ there is $x$ with
$\delta_0x=k$ and $c-\phi_0x\in B$.  Indeed first match $k$ using
the selected generators of $J$; the remaining first coordinate is a
cycle, and the selected cohomology generators match it modulo $B$.

Start with zero modules above degree $d$ and apply this step downward,
using $B=\operatorname{im}d_C^{p-1}$ at degree $p>0$, and $B=0$
at degree zero.  It produces finite free modules $F^p$, a chain map
$\psi:F^\bullet\to C^\bullet$, and the following lifting
properties:
\begin{enumerate}
 \item If $d_Cc=\psi_1k$ and $\delta_1k=0$ in degree zero,
 there is $x\in F^0$ with $\psi_0x=c$, $\delta_0x=k$.
 \item At degree $p+1$, if $d_Cc=\psi_{p+2}k$ and
 $\delta_{p+2}k=0$, there are $c'\in C^p$, $k'\in F^{p+1}$
 such that $d_Cc'+\psi_{p+1}k'=c$ and $\delta_{p+1}k'=k$.
\end{enumerate}
These are just the preceding lifting step applied successively to the
shifted complex; the boundary error at a positive degree supplies $c'$.

Replace $F^0$ by
$K^0=\operatorname{im}(\psi_0,\delta_0)\subset C^0\oplus F^1$,
and set $K^p=F^p$ for $p\ge1$.  Define $\phi_0$ and $d_K^0$
by the two projections.  The module $K^0$ is finite since it is an
image of $F^0$.  The map $(\phi_0,d_K^0)$ is injective by its
definition; the first lifting property says that its image is the
kernel of the next cone differential.  The second property gives
exactness at all subsequent cone terms.  Explicitly, the resulting
exact sequence is
\[
 0\longrightarrow K^0\xrightarrow{(\phi_0,d_K)}
 C^0\oplus K^1\xrightarrow{(c,k)\mapsto(d_Cc-\phi k,-d_Kk)}
 C^1\oplus K^2\longrightarrow\cdots . \tag{GSL.4}
\]
The same lifting properties show surjectivity and injectivity on each
cohomology module: a cycle lifts modulo a boundary by the second
property with $k=0$, and a cycle mapping to a boundary is a boundary
by that property with the given $k$.  In degree zero the first
property and the initial injection apply.  Thus $\phi$ is a
quasi-isomorphism, as claimed.
\end{proof}
% END GSL6

% BEGIN GSL7
\paragraph{GSL7: flat cones give a universally valid projective model.}
\label{claim:GSL7}
Under GSL6, assume in addition that every $C^p$ is flat over $R$.
Then $K^0$ is finite projective, all $K^p$ are finite projective,
and $\phi$ remains a quasi-isomorphism after tensoring with
\emph{any} $R$-module.  In particular it remains one after every
ring map $R\to R'$.

\begin{proof}
For a short exact sequence $0\to A\to B\to C\to0$ with $B,C$
flat, $A$ is flat and this short exact sequence remains exact after
every tensor product.  To see the latter statement, the only possible
failure is injectivity on the left, measured by
$\operatorname{Tor}_1(C,-)=0$.  The same long exact Tor sequence
gives $\operatorname{Tor}_1(A,-)=0$ since both adjacent higher Tor
groups of $B,C$ vanish.  These are precisely the elementary flatness
criteria used here.

All terms $C^p\oplus K^{p+1}$ in (GSL.4) are flat and the sequence
is bounded on the right.  Starting with its terminal zero image and
working leftward, apply the preceding fact to each short exact sequence
formed by an image, a cone term and the next image.  Every image is
flat, including the initial copy of $K^0$.  Every such short exact
sequence is tensor-exact.  Therefore (GSL.4) stays exact after any
tensor product, and its lifting characterization again says that
$1\otimes\phi$ is a quasi-isomorphism.

Finally $K^0$ is finite over the Noetherian ring $R$, hence finitely
presented.  A finitely presented flat module is projective: localize
at a prime, lift a basis of its residue vector space, use flatness to
show that the kernel of the resulting finite free surjection has
zero residue module, and apply Nakayama to this finite kernel.
Thus the localization is free.  Finite presentation spreads this
basis to a principal neighborhood.  On a finite such cover, clear
denominators in the local dual-basis identities; powers of the
covering elements still generate the unit ideal.  A linear
combination of these identities is a global dual-basis identity,
which splits a finite free surjection onto the module.
This proves projectivity and applies to $K^0$.  The positive
terms were already finite free.  Thus the exact universal
finite-projective-complex theorem holds, including its
tensor-by-any-module clause.
\end{proof}
% END GSL7

% BEGIN GSL8
\paragraph{GSL8: arbitrary-base transfer from the original Noetherian cube model.}
\label{claim:GSL8}
Let $I$ be a finite linearly ordered set, $R_0\to R$ a ring map
with $R_0$ Noetherian, and $C_0$ and $A$ cubes of modules over
$R_0$ and $R$, respectively.  Suppose a cube map
$\theta:R\otimes_{R_0}C_0\to A$ is an isomorphism at every
nonempty subset of $I$, each corresponding $C_0(J)$ is flat,
and every cohomology module of the unaugmented \v Cech complex
of $C_0$ is finite over $R_0$.  There is a finite-projective
nonnegative complex $K$ over $R$ and a map
$\phi:K\to\check C(A)$ that is a quasi-isomorphism after every
ring map out of $R$.  Moreover $K^p=0$ for $p\ge |I|$.

\begin{proof}
The degree-$p$ term of $\check C(C_0)$ is the finite direct sum of
$C_0(J)$ over the $(p+1)$-element subsets $J$.  It is flat and
vanishes for $p\ge |I|$.  If $I$ is empty, take the zero complex
and zero map.  Otherwise GSL6--GSL7 provide a finite-projective
complex $K_0$ and a universally tensor-valid quasi-isomorphism
$K_0\to\check C(C_0)$, with the same upper bound.  Put
$K=R\otimes_{R_0}K_0$.  Finite projectives remain finite
projective under arbitrary extension of scalars, since their split
finite free presentations tensor to split presentations.

Tensor product commutes with these finite direct sums and their
alternating restriction differentials.  Hence $\theta$ supplies
an actual isomorphism
$R\otimes_{R_0}\check C(C_0)\simeq\check C(A)$.  The empty
vertex is not a term of this unaugmented complex, so no assertion
about $\theta(\varnothing)$ is needed.  Compose the base-changed
map with this isomorphism.  If $R\to R'$ is any subsequent map,
associativity identifies the resulting map with the one obtained by
tensoring $K_0\to\check C(C_0)$ directly by $R'$ over $R_0$.
GSL7 proves it is a quasi-isomorphism.  This proves the source's
full \texttt{exists\_\allowbreak finite\_\allowbreak projective\_\allowbreak complex\_\allowbreak of\_\allowbreak noetherian\_\allowbreak model}
conclusion without a flatness hypothesis on either ring map.
\end{proof}
% END GSL8

% BEGIN GSL9
\paragraph{GSL9: the actual section \v Cech complex under arbitrary base change.}
\label{claim:GSL9}
Let $X\to\operatorname{Spec}R$ have a finite ordered open cover
$D_i$, with every nonempty finite intersection affine, and let $M$
be an actual line bundle.  The complex
\[
 C^p=\bigoplus_{i_0<\cdots<i_p}
       \Gamma(D_{i_0}\cap\cdots\cap D_{i_p},M)
 \tag{GSL.5}
\]
with alternating restriction differential commutes with every ring
base change.  Its degree-zero kernel is the module of actual global
sections.  If $\mathrm V(M)$ holds, it is exact in degree one.
All these identifications respect restrictions and the actual local
coordinates of sections.

\begin{proof}
On an affine intersection $D=\operatorname{Spec}A$, the line bundle
is represented by its invertible $A$-module $Q=\Gamma(D,M)$.
For any $R\to R'$, its pullback to
$\operatorname{Spec}(R'\otimes_R A)$ is represented by
\[
 (R'\otimes_R A)\otimes_AQ\simeq R'\otimes_RQ.
\]
This follows directly by localizing on charts where $Q$ is free of
rank one and gluing the transition units; no flatness of $R'$ is
needed.  These comparisons commute with restriction to smaller affine
intersections, so they give an isomorphism of the whole complexes.

A degree-zero cycle is a family of sections on the $D_i$ agreeing
on pairwise intersections.  GSL1 glues it uniquely to a global
section.  For the degree-one statement, extend an alternating
\v Cech cocycle $\alpha_{ij}$ to all ordered pairs, putting
$\alpha_{ii}=0$ and $\alpha_{ji}=-\alpha_{ij}$.  Its cocycle
equation reads $\alpha_{ij}+\alpha_{jk}=\alpha_{ik}$ as sections
of the same line bundle on a triple intersection.  Refine each $D_i$
by $D_i\cap U_a$, where $U_a$ is an $M$-chart, and express
$\alpha_{ij}$ in the first chart.  This is exactly a cocycle in
the definition of $\mathrm V(M)$.  Its solution gives sections
$\beta_{ia}$ on the refinement.  For fixed $i$, the differences
between $\beta_{ia}$ and $\beta_{ib}$ are $\alpha_{ii}=0$,
so they glue to $\beta_i$ on $D_i$.  Then
$\alpha_{ij}=\beta_i-\beta_j$, proving exactness.  This argument
also explains why the literal all-cover version of $\mathrm V$
is sufficient even when the $D_i$ themselves do not trivialize $M$.
\end{proof}
% END GSL9

% BEGIN GSL10
\paragraph{GSL10: constructing every field of the literal \texttt{Perf} record.}
\label{claim:GSL10}
In GSL9 suppose there is a nonnegative complex $P^\bullet$ of
finite projective $R$-modules and a map
$\phi:P^\bullet\to C^\bullet$ that is a quasi-isomorphism after
every ring map out of $R$.  Then the source's actual
$\texttt{Perf}(f,M)$ record exists.  In particular this conclusion
holds from the Noetherian cube-model data of GSL8 for the actual
section cube of GSL9.

\begin{proof}
Apply GSL5 to $P^0\to P^1\to P^2$, obtaining matrices
$D_0,D_1$ and fixed elements $u_k\in P^0$.  For each cover index
$i$ and each line-bundle chart $j$, use
\[
 W_{ij}=D_i\cap U_j,\qquad \operatorname{ch}(i,j)=j.
\]
These opens cover $X$.  The $i$-component of $\phi^0(u_k)$ is
an actual section of $M$ on $D_i$; let $u_{k,ij}$ be its coordinate
function on $W_{ij}$.  These functions, fixed over the original
ring, are the \texttt{Perf.\allowbreak u} data.

For any $R\to R'$ and any specified Cartesian model
$X'\to X$ of the base change, GSL5 identifies
$\ker(D_0)_{R'}$ with $\ker(d_P^0)_{R'}$.  The assumed
quasi-isomorphism identifies the latter with $\ker d_C^0$
(degree zero has no incoming boundary).  GSL9 identifies that kernel
with $\Gamma(X',M')$, where $M'$ is the actual pullback.  Their
composite is the required set equivalence $E_{R'}$.  On the inverse
image of $W_{ij}$ its coordinate is precisely
\[
 E_{R'}(v)|_{W_{ij}'}
   =\sum_k c_{f'}(v_k)\,\pi^*u_{k,ij}. \tag{GSL.6}
\]
The equality holds because $\phi$ is a linear map, its scalar
extension acts by tensoring, and the GSL9 comparison uses the
actual pulled coordinates.  Thus it is the printed
\texttt{Perf.\allowbreak sec} equality, not merely an abstract dimension
comparison.

If $\mathrm V(M')$ holds, GSL9 gives $H^1(C_{R'})=0$.
The universal quasi-isomorphism gives exactness of
$P^0_{R'}\to P^1_{R'}\to P^2_{R'}$, and GSL5 gives exactness
of the two matrices.  This is exactly \texttt{Perf.\allowbreak exact}.
The fields $a,b,c,D_0,D_1,\mathrm{comp},I,W,\mathrm{ch},
\mathrm{le},\mathrm{cover},u,\mathrm{sec},\mathrm{exact}$ have
all been constructed.  Every assertion holds for arbitrary $R'$ and
for every specified pullback square, since its unique comparison
isomorphism with the standard tensor-product model transports the
same chart equalities.

To apply this result to a \texttt{Very\allowbreak Ample\allowbreak By} presentation,
take $D_i=D(s_i)$.  On $D_i$ the nonvanishing section $s_i$
trivializes the presenting line bundle, so
$D_i\cap D_j$ is the principal open defined by $s_j/s_i$ in
the affine $D_i$.  More generally every nonempty finite
intersection is obtained by inverting a product of these
ratios in one $D_i$, hence is affine.  The finite cover and
all hypotheses of GSL9 therefore hold.  Consequently the
source's universally valid finite-projective model for this
actual section complex, whenever supplied, proves the whole
\texttt{nonempty\_\allowbreak perf\_\allowbreak of\_\allowbreak smooth\_\allowbreak is\allowbreak Proj\allowbreak By} implication.
\end{proof}
% END GSL10

% BEGIN GSL11
\paragraph{GSL11: linearity and naturality forced by \texttt{Perf}'s coordinates.}
\label{claim:GSL11}
Let an actual $\texttt{Perf}(f,M)$ be given, without choosing a
preferred family of its set equivalences.  Its coordinate equation
forces every equivalence $E_{R'}$ to be linear and uniquely determines
it.  For a subsequent ring map $R'\to R''$, the pullback of
$E_{R'}(v)$ is $E_{R''}(v_{R''})$.  The same statements hold after
the literal \texttt{LPow} and \texttt{Lin\allowbreak Comb} chart conversions.

\begin{proof}
The inverse images of the fixed $W_i$ cover every pullback scheme.
The right side of (GSL.6) is additive in $v$ and linear over the
base ring.  GSL1 therefore identifies the section for $v+v'$ with
the sum of the sections for $v,v'$, and the section for $av$ with
$a$ times the section for $v$.  It also proves that two purported
sections with those coordinates are equal.  Thus the set equivalence
is both linear and uniquely determined by its coordinate formula.

Pull (GSL.6) through the subsequent map.  Pullback commutes with the
finite sum, products, structure maps and the fixed functions $u_{ki}$.
The resulting expression is exactly (GSL.6) for the base-changed
vector.  Uniqueness gives the asserted naturality.  The chart maps
in GSL1 preserve this equation by restrictions, hence preserve its
uniqueness and all these operations.  In particular, the argument
does not assume that an arbitrarily selected equivalence in the
definition of \texttt{Perf} was natural to begin with.
\end{proof}
% END GSL11

% BEGIN GSL12
\paragraph{GSL12: a single local section basis valid on every affine test.}
\label{claim:GSL12}
Let $F=\texttt{Perf}(f,M)$ and $\gamma:R\to K$ be a field point
at which the actual pulled line bundle has $\mathrm V(M_K)$.
There are $r\in R$ with $\gamma(r)\ne0$ and actual sections
$b_1,\ldots,b_m\in\Gamma(X,M)$ whose restrictions form a basis
after \emph{every} map $R\to R'$ inverting $r$.  Here
$m=\dim_K\Gamma(X_K,M_K)$.  If $R$ is local and $\gamma$ factors
through an extension of its residue field, they already form a basis
over $R$.

\begin{proof}
The \texttt{Perf.\allowbreak exact} field gives middle exactness at $K$ for
$D_0,D_1$.  Apply GSL3--GSL4, and let $b_i$ be the actual section
represented by the $i$th column of $B$.  This column is a cycle
over $R$ itself, since $D_0B=0$ before localization.  GSL11 says
that any actual pulled section satisfying \texttt{Is\allowbreak Pull\allowbreak Sec} or
\texttt{Is\allowbreak Pull\allowbreak Sec\allowbreak LC} is represented by the pulled column.  GSL4
gives unique coordinates in these columns for every kernel vector;
GSL11 transfers them, with their uniqueness, to the actual section
module.  Thus the sections are independent and span on every test.
At $K$ they form a basis of length $m$, so that length is the stated
dimension.  The final assertion is the local-unit clause of GSL4.
\end{proof}
% END GSL12

% BEGIN GSL13
\paragraph{GSL13: lifting actual sections through a surjective coefficient map.}
\label{claim:GSL13}
Let $F=\texttt{Perf}(f,M)$, where $M$ may be the literal
$\operatorname{LinComb}(L,\phi,\epsilon)$, and let
$\rho:R\to\overline R$ be a surjective ring map.  Suppose a fixed
matrix $B$ has $D_0B=0$ over $R$ and its columns span
$\ker(D_0)_{\overline R}$.  Then every section of the actual pulled
$M$ over $\overline R$ lifts to a section over $R$.  For a
\texttt{Lin\allowbreak Comb}, the lift satisfies the printed
\texttt{Is\allowbreak Pull\allowbreak Sec\allowbreak LC} coordinate equalities.

\begin{proof}
Use the coordinate equivalence for the specified pullback square
to represent a given section $\overline s$ by
$\overline v\in\ker(D_0)_{\overline R}$.  Write
$\overline v=\overline B\,\overline z$ and lift each coordinate
of $\overline z$ to $z\in R^m$.  The vector $v=Bz$ is a cycle
over $R$.  The identity-base-change equivalence, followed by the
inverse of the harmless identity-pullback chart conversion in GSL1,
gives an actual section $s$ of $M$.  By GSL11 its pullback has the
same coordinates as $\overline s$ on the inverse images of all
$W_i$, hence equals $\overline s$.

For clarity, if $M=\operatorname{LinComb}(L,\phi,\epsilon)$ and
$\pi$ is the map on total spaces, the target in the source is the
line bundle with maps $\pi\phi_k$ on the same finite index set, not
the syntactically nested bundle $M.\mathrm{pull}(\pi)$.  The chart
conversions of GSL1 identify these bundles by the identity map on
each common coordinate function.  Accordingly the equality just
proved reads, for every such chart $j$,
\[
 \overline s_j
   =\operatorname{res}\bigl(\pi^*s_j\bigr).
 \tag{GSL.7}
\]
This is exactly \texttt{Is\allowbreak Pull\allowbreak Sec\allowbreak LC}.  Surjectivity of $\rho$ and
the single fixed spanning matrix are all that the argument needs;
neither locality nor a field hypothesis enters this lifting lemma.
\end{proof}
% END GSL13

\subsection{The geometric inputs and their proved reductions}

The next arguments use the original \texttt{Pol\allowbreak AV} record $P/S$.
It contains a commutative abelian scheme $A/S$, its relative dimension
$g$, and an actual rigidified symmetric relatively ample line bundle
$L$.  Its \texttt{KTriv} condition tests the translation kernel over
every affine base: local triviality of $t_x^*L\otimes L^{-1}$ forces
$x$ to be the identity.  In particular we do not replace this record
by the data of an arbitrary ample bundle or by complex fibers alone.

We state the upstream inputs when they are used.  The original
\texttt{Pol\allowbreak Base\allowbreak Change\allowbreak St}, denoted $\mathcal P$, says: for every map
$t:S'\to S$ and every $P/S$, there is a polarized record $P'/S'$
and a Cartesian map $G:A'\to A$ with $g'=g$, preserving multiplication
on every test scheme, together with an actual generated
line-bundle isomorphism $G^*L\simeq L'$.  This entire record remains
an input here.

% BEGIN GSL14
\paragraph{GSL14: transport of sections, bases and the all-cover vanishing law.}
\label{claim:GSL14}
Let $e:X'\to X$ be a scheme isomorphism and let $M',M$ be actual
line bundles, with a generated isomorphism $M'\simeq e^*M$.
It induces a linear equivalence of section modules, preserving
specified coefficient-ring actions when $e$ is over the base.
It transports bases and their unique coordinate expressions.
Moreover $\mathrm V(M)$ implies $\mathrm V(M')$, with the exact
all-cover meaning used above.

\begin{proof}
Refine the charts of $M'$ by inverse images of charts of $M$.
On a common chart the generated isomorphism is multiplication by
a specified unit $u$.  Send a coordinate function $s$ to
$u\,e^*s$ (or to $u^{-1}e^*s$, according to the displayed direction
of the isomorphism).  The defining unit compatibility is exactly
the equation ensuring that these functions form a section.
GSL1 glues them.  The inverse scheme map and inverse units give
the inverse section map.  Both maps are linear, commute with the
base structure maps, and are mutual inverses on every chart.
They therefore transport bases and unique coefficients.

For vanishing, start with a cocycle for $M'$ on an arbitrary
subordinate cover $V_a$.  Refine it by the inverse images of
the $M$-charts and transport the cocycle through the inverse units.
The unit compatibility turns its cocycle equation into the one
for $e^*M$.  Transporting across the isomorphism $e$ gives a
cocycle for $M$; solve it using $\mathrm V(M)$ and pull the
solution back.  Multiply by the original units.  This gives a
solution on the refinement for the initial cocycle.

For fixed $a$, its solutions on two refined pieces of $V_a$
agree: their required difference is the diagonal cocycle
$\phi_{aa}$, which is zero by the cocycle equation with three
equal indices and $g_{i(a)i(a)}=1$.  Thus the solutions glue to
one function on each $V_a$.  The coboundary equation holds
locally and hence on every $V_a\cap V_b$.  This proves the
original \texttt{h1Zero\_\allowbreak of\_\allowbreak is\allowbreak Iso} statement, including changes
of both the scheme and the literal line-bundle charts.
\end{proof}
% END GSL14

% BEGIN GSL15
\paragraph{GSL15: the required projective presentation after any affine base change.}
\label{claim:GSL15}
Assume $\mathcal P$.  For every $P/S$, every commutative $R$,
$t:\operatorname{Spec}R\to S$, and every specified Cartesian model
\[
 \begin{array}{ccc}
 X&\xrightarrow{\pi}&A\\
 f\downarrow&&\downarrow\\
 \operatorname{Spec}R&\xrightarrow{t}&S,
 \end{array}
\]
there exist $N\ge1$, $m\ge0$ and $m+1$ actual sections of
$(\pi^*L)^N$ giving the source's \texttt{Very\allowbreak Ample\allowbreak By} presentation
of $f$.  In particular $f$ satisfies \texttt{Is\allowbreak Proj\allowbreak By}.

\begin{proof}
Apply $\mathcal P$ to obtain $P'/\operatorname{Spec}R$.
The two Cartesian models are uniquely isomorphic over
$\operatorname{Spec}R$, and the generated isomorphism in
$\mathcal P$ identifies their actual pulled line bundles.
Apply the \texttt{Rel\allowbreak Ample} field of $P'$ on the whole affine
base $\operatorname{Spec}R$.  It supplies a positive power,
$m+1$ sections whose nonzero loci cover, an affine nonzero
locus for each section, and on each such locus ratios $r_j$
with $s_j=r_js_i$ and a surjection
\[
 R[X_0,\ldots,X_m]\longrightarrow
 \Gamma(D(s_i),\mathcal O),\qquad X_j\longmapsto r_j.
 \tag{GSL.8}
\]
This is the original definition of relative ampleness used here.

Transport all these sections to $X$ by GSL14, also converting
the repeated-factor power to the actual $\operatorname{pw}$
charts.  Multiplication by a transition unit preserves whether
the germ of a section is a unit; consequently the nonzero
loci are the inverse images of the old ones.  They cover and
remain affine under the scheme isomorphism.  The section
isomorphism multiplies numerator and denominator by the
same unit, so the ratios are pulled functions $r_j$.
The corresponding map (GSL.8) is conjugate through ring
isomorphisms to the old surjection.  It remains surjective.
These are every field of \texttt{Very\allowbreak Ample\allowbreak By}, with the same
positive integer $N$.  No restriction on $R$ or $S$ was imposed.
\end{proof}
% END GSL15

% BEGIN GSL16
\paragraph{GSL16: descent of first-cocycle vanishing by a finite average.}
\label{claim:GSL16}
Let $f:Y\to X$, let a finite group $G$ act on $Y$ by maps
$\tau_g$ over $X$, with the composition convention
$\tau_g\tau_h=\tau_{gh}$, and suppose $|G|$ is a unit in
$\Gamma(X,\mathcal O)$.  Assume that on every open $U\subseteq X$
pullback identifies $\Gamma(U,\mathcal O)$ with the functions on
$f^{-1}U$ invariant under every $\tau_g$: each invariant function
has a unique preimage.  If $\mathrm V(f^*M)$ holds, then
$\mathrm V(M)$ holds.

\begin{proof}
First this function-descent premise makes every pullback
$\Gamma(U,\mathcal O)\to\Gamma(f^{-1}U,\mathcal O)$ injective.
Indeed a pulled function is invariant because $f\tau_g=f$,
and uniqueness of its descended preimage proves injectivity.

Pull an arbitrary cocycle $\phi_{ab}$ for $M$ on a subordinate
cover $V_a$ to $Y$.  Vanishing there gives
\[
 f^*\phi_{ab}=\psi_a-f^*g_{i(a)i(b)}\psi_b.
\]
Set $\Psi_a=\sum_{g\in G}\tau_g^*\psi_a$.  Reindexing the sum
by group multiplication proves that $\Psi_a$ is invariant,
so it descends uniquely to a function $\chi_a$ on $V_a$.
Every pulled transition function and every $f^*\phi_{ab}$ is
invariant.  Sum the displayed coboundary identity, and use the
injectivity just proved to descend the resulting equation:
\[
 |G|\phi_{ab}=\chi_a-g_{i(a)i(b)}\chi_b.
\]
Multiply each $\chi_a$ by the restriction of $|G|^{-1}$.
The resulting functions solve the original cocycle.  The argument
uses every original cover without a refinement or a separability
assumption on $f$ beyond the stated invariant-function law.
\end{proof}
% END GSL16

% BEGIN GSL17
\paragraph{GSL17: positive-power vanishing from the original geometric suppliers.}
\label{claim:GSL17}
Assume $\mathcal P$ and the following original inputs.

\begin{enumerate}
 \item \emph{Serre input.}  For every Noetherian $R$, every actual
 $\texttt{Very\allowbreak Ample\allowbreak By}(f,H,s)$ and every actual line bundle $M$,
 there is $t_0\in\mathbb Z$ such that
 $\mathrm V(M\otimes H^t)$ for all integers $t\ge t_0$.
 \item \emph{Multiplication pullback.}  For every $S,P,n$ there
 is an actual generated trivialization of
 $L^{n^2}\otimes([n]^*L)^{-1}$.
 \item \emph{Multiplication geometry.}  For every locally Noetherian
 $S$, $P/S$ and $n\ge1$, the map $[n]$ is finite, flat and
 surjective, and is etale if $n$ is a unit on $S$.
 \item \emph{Torsion-invariant descent.}  For every algebraically
 closed $K$, $P/K$ and $n$ nonzero in $K$, functions on
 $[n]^{-1}U$ invariant under all translations by $A[n](K)$
 descend uniquely to $U$, for every open $U\subseteq A$.
\end{enumerate}
Then for every algebraically closed field $K$, every original
$P/K$ and every $n\ge1$, one has $\mathrm V(L^n)$.

\begin{proof}
By GSL15 choose $N\ge1$ with $H=L^N$ very ample in the printed
sense.  Apply the Serre input to $L^r$ for each
$0\le r<N$.  Choose a nonnegative integer $T$ at least all
the resulting bounds.  Choose a prime $q$ sufficiently large
that $q$ is nonzero in $K$ and $q^2n/N\ge T+1$.
There are arbitrarily large primes, so this choice works also
in positive characteristic by avoiding that characteristic.
Write
\[
 q^2n=r+Nt,\qquad 0\le r<N.
\]
Then $t\ge T$ and the Serre input gives
$\mathrm V(L^r\otimes(L^N)^t)=\mathrm V(L^{q^2n})$.
The multiplication-pullback generator, raised to its $n$th
power and combined with the identity $q^2n=r+Nt$, gives an
actual generated isomorphism
$[q]^*L^n\simeq L^{q^2n}$.  On common charts its unit is the
$n$th power of that generator times the necessary transition
units; equality of the exponents makes the cocycle equation
identical on both sides.  GSL14 therefore gives
$\mathrm V([q]^*L^n)$.

Let $G=A[q](K)$.  It is finite: its points belong to the finite
fiber over the identity of the finite map $[q]$.  It is an
abelian group killed by the prime $q$, hence a finite vector
space over $\mathbb F_q$ and has order $q^a$ for some $a\ge0$.
Because $q\ne0$ in $K$, this order is a unit in $K$ and hence
in $\Gamma(A,\mathcal O)$.  Translation by $G$ is an action
over $[q]$, since $[q](x+y)=[q]x$ for $y\in G$.
The invariant-function input is exactly the descent premise
of GSL16.  Apply that claim with $f=[q]$ and $M=L^n$.
It yields $\mathrm V(L^n)$.  This proves the entire original
positive-power vanishing supplier, uniformly over all
algebraically closed characteristics, from the four displayed
inputs; no characteristic-zero vanishing theorem was inserted.
\end{proof}
% END GSL17

% BEGIN GSL18
\paragraph{GSL18: the exact exponent forced by growth and square dilation.}
\label{claim:GSL18}
Let $h:\mathbb N\to\mathbb N$ be nondecreasing, $h(1)=1$, and
$G\subseteq\mathbb N$ satisfy $0\notin G$ and, for every $m$,
at least one of $m,m+1$ belongs to $G$.  Suppose
\[
 h(n^2j)=h(n^2)h(j)\quad(n\in G,\ j\ge1), \tag{GSL.9}
\]
and there are $N\ge1$, $c,C\in\mathbb N$ and $g\ge0$ such that
\[
 t^g\le c\,h(Nt),\qquad h(Nt)\le Ct^g\quad(t\ge1).
 \tag{GSL.10}
\]
Then $h(j)=j^g$ for every $j\ge1$.

\begin{proof}
Fix $n\in G$ and put $a=n^2\ge1$.  Iterating (GSL.9) gives
$h(a^ki)=h(a)^kh(i)$ for every $i\ge1$.  By monotonicity and
the upper bound in (GSL.10),
\[
 h(a)^k=h(a^k)\le h(Na^k)\le C(a^g)^k.
\]
Since a ratio greater than one has unbounded powers, this
forces $h(a)\le a^g$.  The lower bound in (GSL.10) gives
\[
 (a^g)^k\le c\,h(Na^k)=c\,h(N)h(a)^k.
\]
Here $h(a)\ge1$, so the same argument forces $a^g\le h(a)$.
Thus $h(n^2)=n^{2g}$ for every $n\in G$.

There are arbitrarily large $n\in G$.  Fix $j\ge1$ and for each
such $n$ choose $m_-,m_+\in G$ immediately below and above
$n\sqrt j$.  They can be chosen with
$0\le n\sqrt j-m_-\le2$ and
$0\le m_+-n\sqrt j\le2$, because $G$ meets every pair of
consecutive integers.  For large $n$ both are positive.
Monotonicity and (GSL.9) yield
\[
 m_-^{2g}\le n^{2g}h(j)\le m_+^{2g}.
\]
Divide by $n^{2g}$ and let these $n$ tend to infinity.
Both bounds tend to $(\sqrt j)^{2g}=j^g$.
This proves the assertion, including $g=0$.  The use of real
limits merely shortens the source's equivalent integer estimates.
\end{proof}
% END GSL18

% BEGIN GSL19
\paragraph{GSL19: the original geometric rank is exactly \(n^g\).}
\label{claim:GSL19}
Assume the following source inputs for every original $P/K$ over
an algebraically closed field.

\begin{enumerate}
 \item Its geometric total space is integral (the source obtains this
 from smoothness and geometric connectedness).
 \item $h^0(L)=1$.
 \item For $n\ne0$ in $K$ and $j\ge1$ there is an actual
 coefficient-linear equivalence
 \[
 \Gamma(A,L^{n^2j})\simeq
       \prod_{x\in A[n](K)}\Gamma(A,L^j).
 \tag{GSL.11}
 \]
 \item Every $\Gamma(A,L^j)$ for $j\ge0$ is finite, and for some
 $N,c,C\in\mathbb N$ with $N,c\ge1$ its dimensions obey
 (GSL.10) with the original $g=P.g$.
\end{enumerate}
Then \texttt{Global\allowbreak Sections\allowbreak Rank\allowbreak Alg\allowbreak Closed} holds: for every such
$P$ and every $j\ge1$ the literal \texttt{Sec\allowbreak M} has dimension $j^g$
over $\Gamma(\operatorname{Spec}K,\mathcal O)$.

\begin{proof}
Use GSL1 and the identity-base product isomorphism to identify
the displayed modules with the literal \texttt{Sec\allowbreak M} modules,
preserving the canonical base coefficient ring.  Put
$h(j)=\dim_K\Gamma(A,L^j)$.  The rank-one input gives a nonzero
actual section $s$ of $L$.  Multiplication by $s$ maps
$\Gamma(A,L^j)$ injectively into $\Gamma(A,L^{j+1})$.
Indeed, on an integral scheme a nonzero section of a line bundle
is nonzero in every nonempty trivializing open: otherwise its
zero coordinate, transferred across a nonempty intersection,
would make it zero in every chart.  The coordinate rings of
nonempty opens of an integral scheme are domains, so
$fs=0$ forces every local coordinate of $f$ to vanish.
Thus $h$ is nondecreasing, using the finite-dimensionality input.

Let $G=\{n:(n:K)\ne0\}$.  It omits zero and meets each pair
of consecutive integers, since both cannot vanish in a field.
For $n\in G$, the $j=1$ instance of (GSL.11) implies that
$A[n](K)$ is finite: otherwise its product with the nonzero
one-dimensional space $\Gamma(A,L)$ would contain arbitrarily
large independent families of coordinate vectors, contradicting
the finite dimension of $\Gamma(A,L^{n^2})$.
The same instance gives $h(n^2)=|A[n](K)|$.  The general $j$
instance therefore gives
$h(n^2j)=|A[n](K)|h(j)=h(n^2)h(j)$.
Apply GSL18 with the stated growth bounds.  The resulting
$h(j)=j^g$ is preserved by the canonical coefficient-ring
isomorphism, proving the printed rank law.  The proof uses the
actual section equivalences, not an independently assumed
degree formula for $[n]$.
\end{proof}
% END GSL19

% BEGIN GSL20
\paragraph{GSL20: rank and vanishing on every actual geometric fiber.}
\label{claim:GSL20}
Assume $\mathcal P$, positive-power vanishing on every $P'/K$,
and \texttt{Global\allowbreak Sections\allowbreak Rank\allowbreak Alg\allowbreak Closed}.  Then for every scheme
$S$, original $P/S$, algebraically closed $K$, point
$t:\operatorname{Spec}K\to S$ and $n\ge1$,
\[
 \dim_{\Gamma(\operatorname{Spec}K,\mathcal O)}
    \operatorname{SecM}(P.A,P.L,n,\operatorname{Over.mk}t)
       =n^{P.g}. \tag{GSL.12}
\]
The actual positive power on any specified Cartesian model of this
fiber has $\mathrm V$.  The same assertions apply after any
intermediate base scheme $T\to S$ and any geometric point of $T$.

\begin{proof}
Apply $\mathcal P$ to $t$ and obtain $P'/K$, its Cartesian
map $G:A'\to A$, $P'.g=P.g$, and its actual line-bundle
isomorphism.  Every specified fiber model is uniquely
isomorphic over $K$ to $A'$.  Raise the generated isomorphism
to the $n$th power and apply GSL14.  This transports both
the section module and the all-cover vanishing assertion
from $(L')^n$.  The one-factor \texttt{Lin\allowbreak Comb} chart maps
of GSL1 identify the section module with the exact
\texttt{Sec\allowbreak M} in (GSL.12), so the closed-field rank law gives
the stated integer.

For $T\to S$ and $\sigma:\operatorname{Spec}K\to T$, apply
the same construction to their composite.  The universal
property of the two successive pullbacks gives a Cartesian
map from $A'$ to $T\times_S A$.  Its composite to $A$ is
$G$, so the actual pulled powers and their chart isomorphisms
are precisely those already used.  Neither $S$ nor $T$ has
been assumed Noetherian or reduced.  This proves the exact
full-fiber rank and vanishing inputs of the two section-law
theorems, not only their instances at a chosen affine chart.
\end{proof}
% END GSL20

% BEGIN GSL21
\paragraph{GSL21: the geometric root and the actual translated-power generator.}
\label{claim:GSL21}
Assume the multiplication geometry of GSL17 and the original
square input: for every $S,P$, every $T\to S$, points $x,y\in A(T)$
and point of $T$, the bundle
\[
 t_{x+y}^*L\otimes L\otimes(t_x^*L)^{-1}
                    \otimes(t_y^*L)^{-1}
 \tag{GSL.13}
\]
has a generated trivialization after restriction above some open
neighborhood of that point.  Here the notation is on $T\times_S A$.
Then, for every locally Noetherian $S$, every $P/S$, every
algebraically closed field $K$, $t:\operatorname{Spec}K\to S$,
$q\ge1$ and $c\in A(K)$ above $t$, there exist $d$ with $qd=c$
and an actual generated trivialization of
\[
 L^{q-1}\otimes t_c^*L\otimes(t_d^*L)^{-q}. \tag{GSL.14}
\]
This is the full original \texttt{twisted\_\allowbreak pow\_\allowbreak translate} supplier.

\begin{proof}
The fiber of the finite surjective map $[q]$ above the $K$-point
$c$ is a nonempty finite $K$-scheme.  It has a closed point with
residue field finite over $K$, hence equal to $K$ because $K$ is
algebraically closed.  Its point gives $d$ with $qd=c$.  This
argument does not require $q$ to be invertible in $K$.

On $\operatorname{Spec}K$ any neighborhood of its unique point
is the whole base.  Thus the square input gives an actual
global generator for (GSL.13) for each pair of $K$-points.
Put $M_x=t_x^*L$ on the geometric fiber.  The generator
identifies $M_{x+y}\otimes M_0$ with $M_x\otimes M_y$.
Multiplying these generators successively with $(x,y)=(ad,d)$
proves by induction that
\[
 M_{qd}\otimes M_0^{q-1}\otimes M_d^{-q}
\]
has an actual generator.  The starting case $q=1$ is the
identity generator; the step multiplies the previous generator
by the square generator and cancels the common factors.
All cancellations are products or inverses of the specified
units on a common chart, so they yield a generated
\texttt{Lin\allowbreak Comb} isomorphism, not just a numerical identity.
Since $qd=c$, this is (GSL.14) with the exact maps
$[\mathrm{snd},\mathrm{translPt}(c),\mathrm{translPt}(d)]$
and exponents $[q-1,1,-q]$.
\end{proof}
% END GSL21

% BEGIN GSL22
\paragraph{GSL22: a pure-power model of each twisted geometric fiber.}
\label{claim:GSL22}
Assume $\mathcal P$ and the conclusion of GSL21 at its printed
locally Noetherian scope.  Let $S$ now be \emph{arbitrary},
$P/S$ original, $T\to S$ any base, $c\in A(T)$, $q\ge1$ and
$\sigma:\operatorname{Spec}K\to T$ an algebraically closed point.
Put $X=T\times_S A$ and
$M=L_T^{q-1}\otimes t_c^*L$ on $X$.  There are a polarized
$P'/K$ with $g'=g$ and a Cartesian map
$\pi:A'\to X$ over $\sigma$ for which
\[
 \pi^*M\simeq (L')^q \tag{GSL.15}
\]
is an actual generated isomorphism.  One may equally use the
literal $\operatorname{baseOver}(K)\times_K A'$ model and its
$\operatorname{LPow}$ in (GSL.15).  Thus, under the rank and
vanishing inputs of GSL20, $\pi^*M$ has dimension $q^g$ and
has $\mathrm V$.

\begin{proof}
Apply $\mathcal P$ to the composite $\operatorname{Spec}K\to S$,
obtaining $P'/K$, its multiplication-preserving Cartesian
map $G:A'\to A$, and $G^*L\simeq L'$.  The point $c$ lifts
uniquely to $c'\in A'(K)$.  The resulting Cartesian map
$j:A'\to X$ satisfies, on every test scheme,
\[
 j^*L_T\simeq L',\qquad
 j^*(t_c^*L)\simeq t_{c'}^*L'.
 \tag{GSL.16}
\]
The second equality uses the actual multiplication-preserving
field of $\mathcal P$: $G(x+c')=G(x)+c$.
Take the powers and tensor products of its given units to
obtain these actual generated isomorphisms.

Apply GSL21 to the new base $\operatorname{Spec}K$ and $P'$.
This base is Noetherian regardless of the original $S$.
It gives $d\in A'(K)$ with $qd=c'$ and a generated
isomorphism
$(L')^{q-1}\otimes t_{c'}^*L'\simeq t_d^*((L')^q)$.
Let $\tau=t_{-d}:A'\to A'$.  It is an isomorphism over $K$
with inverse $t_d$.  Set $\pi=j\circ\tau$.
Precomposing a Cartesian square by this isomorphism over the
base keeps it Cartesian.  Pulling the preceding line-bundle
isomorphism by $\tau$ gives (GSL.15), because
$t_d\circ t_{-d}=1$.

For the source's literal \texttt{base\allowbreak Over} product, compose
with its canonical isomorphism to $A'$ and use the one-factor
chart maps of GSL1.  This is the exact source construction
$\pi=\tau\gg j$, including its inverse-translation sign.
GSL14 transports the rank and all-cover vanishing of $(L')^q$
to this actual $\pi^*M$.  In particular no locally Noetherian
assumption on $S$ has entered through the geometric twist input.
\end{proof}
% END GSL22

\subsection{The original global laws, with their exact quantifier order}

The remaining geometric model premise, denoted $\mathcal F$, is the
literal source statement
\[
 \begin{gathered}
 \text{for every commutative }R,\quad
 f:X\to\operatorname{Spec}R\text{ smooth and }\texttt{Is\allowbreak Proj\allowbreak By},\\
 \text{and every actual line bundle }M,\quad
        \operatorname{Nonempty}(\texttt{Perf}(f,M)).
 \end{gathered} \tag{GSL.17}
\]
GSL5--GSL10 prove its complete finite-complex construction once the
actual section cube has the stated Noetherian flat model with finite
cohomology.  The global existence of that geometric model is a retained
supplier.  More precisely, the original source supplies it through
\texttt{exists\_\allowbreak noetherian\_\allowbreak model\_\allowbreak gr\allowbreak Cube} and the dehomogenization
comparison with the actual section cube, followed by
\texttt{exists\_\allowbreak finite\_\allowbreak projective\_\allowbreak complex\_\allowbreak cech\allowbreak LB}.  Its graded
presentation, smooth finite-presentation descent and finite graded
cohomology inputs are not proved by the matrix argument alone.

Write $\mathcal V$ for positive-power $\mathrm V$ over every
algebraically closed field, $\mathcal R$ for
\texttt{Global\allowbreak Sections\allowbreak Rank\allowbreak Alg\allowbreak Closed}, and $\mathcal T$ for the
full locally Noetherian geometric translated-power conclusion in
GSL21.  GSL17, GSL19 and GSL21 respectively give their complete
reductions from the displayed original geometric inputs.

% BEGIN GSL23
\paragraph{GSL23: both original section base-change laws.}
\label{claim:GSL23}
Assume $\mathcal F,\mathcal P,\mathcal V,\mathcal R$.
Then \texttt{Sec\allowbreak Base\allowbreak Change\allowbreak Any\allowbreak St} and \texttt{Sec\allowbreak Base\allowbreak Change\allowbreak St} both
hold.  Explicitly, for every scheme $S$, every original $P/S$,
$n\ge1$ and point $p\in S$, there exist an open $V\ni p$
and actual sections $b_i$, $i\in\operatorname{Fin}(n^g)$,
of the specified $\operatorname{LPow}$ over $V$, such that
for every commutative ring $R'$ and every
$t:\operatorname{Spec}R'\to V$, \emph{every} tuple $b_i'$
satisfying the literal \texttt{Is\allowbreak Pull\allowbreak Sec} equalities is
linearly independent and spans the literal \texttt{Sec\allowbreak M}
over $\Gamma(\operatorname{Spec}R',\mathcal O)$.

\begin{proof}
Choose an affine open neighborhood of $p$, written
$t_0:\operatorname{Spec}R\hookrightarrow S$, and let
$\mathfrak p\subset R$ be its point above $p$.
On the actual pullback $X=\operatorname{Spec}R\times_S A$
let $M=\operatorname{LPow}(A,L,n,\operatorname{Over.mk}t_0)$.
The morphism $X\to\operatorname{Spec}R$ is smooth by base
change.  GSL15 supplies its \texttt{Is\allowbreak Proj\allowbreak By} presentation.
Thus $\mathcal F$ provides an actual $\texttt{Perf}(f,M)$.

Take $K$ to be an algebraic closure of the residue field
$\kappa(\mathfrak p)$ and let $\gamma:R\to K$ be the
canonical map.  Its kernel is $\mathfrak p$.
GSL20 gives $\mathrm V$ on an actual Cartesian model of
the geometric pullback and gives its section dimension
$n^g$.  Apply GSL12 to obtain a scalar $r\notin\mathfrak p$
and exactly $n^g$ actual sections $b_i^0$ over $X$.
The columns representing these sections were chosen over
$R$ and satisfy $D_0B=0$ there.

Let $V$ be the image of $D(r)\subseteq\operatorname{Spec}R$
under the chosen open immersion.  It contains $p$ and
is canonically isomorphic to $\operatorname{Spec}R[r^{-1}]$.
Pull the $b_i^0$ to the specified restriction object over
$V$ to obtain $b_i$.  For any later $t:\operatorname{Spec}R'\to V$,
the corresponding map $R\to R'$ sends $r$ to a unit.
GSL12 therefore gives unique section coordinates in the
pulled columns.  A tuple $b_i'$ satisfying the original
\texttt{Is\allowbreak Pull\allowbreak Sec} conditions consists of exactly those
sections by GSL1 and GSL11.  Its unique-coordinate property
is equivalent to linear independence and full span in
$\operatorname{SecM}$: the map from $R'^{n^g}$ taking
coefficients to the linear combination of $b_i'$ is a
linear bijection.  The canonical $\Gamma\operatorname{Spec}$
ring isomorphism gives the printed coefficient action.

The same $V$ and tuple $b_i$ were selected before all
$R',t,b'$.  This proves \texttt{Sec\allowbreak Base\allowbreak Change\allowbreak Any\allowbreak St} with
arbitrary original $S$ and arbitrary coefficient ring $R'$.
Restricting the original $S$ to locally Noetherian schemes
gives \texttt{Sec\allowbreak Base\allowbreak Change\allowbreak St}.  Thus every premise in the
original \texttt{sec\allowbreak Base\allowbreak Change\_\allowbreak of\_\allowbreak perf} theorem has either
been retained exactly or supplied above: projectivity by
GSL15, isomorphism invariance by GSL14, the universal matrix
theorem by GSL3--GSL4 and the full geometric-fiber rank
statement by GSL20.
\end{proof}
% END GSL23

% BEGIN GSL24
\paragraph{GSL24: the actual twisted basis over every local coefficient ring.}
\label{claim:GSL24}
Assume $\mathcal F,\mathcal P,\mathcal V,\mathcal R,\mathcal T$.
For every arbitrary scheme $S$, original $P/S$, $q\ge1$,
commutative local ring $R$, $t:\operatorname{Spec}R\to S$
and $c\in A(R)$ above $t$, the actual bundle
\[
 M=\operatorname{twistLB}(A,L,q,c)
     =L_R^{q-1}\otimes t_c^*L
\]
has a basis of $q^g$ actual sections $b_i$ over $R$.
For every actual section $s$ there is a unique
$a:\operatorname{Fin}(q^g)\to\Gamma(\operatorname{Spec}R,\mathcal O)$
such that on each literal common chart $j$,
\[
 s_j=\sum_i c_f(a_i)b_{i,j}. \tag{GSL.18}
\]
The same sections remain a basis after every ring map $R\to R'$.

\begin{proof}
Set $X=\operatorname{Spec}R\times_S A$ with its actual projection
$f$.  Smoothness is stable under base change, and GSL15 makes
$f$ \texttt{Is\allowbreak Proj\allowbreak By}.  The tensor formula for $M$ uses the
given transition units, so $M$ is an actual line bundle.
Obtain $F=\texttt{Perf}(f,M)$ from $\mathcal F$.

Let $\kappa$ be the residue field of $R$ and $K$ an algebraic
closure of $\kappa$, with $\gamma:R\to\kappa\hookrightarrow K$.
GSL22 produces a specified Cartesian model of $M_K$ as a
positive power of an original $P'/K$, with dimension $q^g$
and $\mathrm V$.  Apply GSL3--GSL4 to the matrices of $F$
at this geometric model.  The resulting $r$ has nonzero
image in $K$, hence is a unit of $R$, since $\gamma$ factors
through the residue field.  Let $b_i$ be the sections
represented by the fixed columns of $B$.

Over $K$ these sections form a basis.  Their number is
therefore $q^g$ by GSL22, and we reindex the fixed columns
once by $\operatorname{Fin}(q^g)$.  GSL11 and GSL4 give
unique coordinates over $R$, yielding (GSL.18).
Because $r$ is already a unit, its image is a unit under
every ring map $R\to R'$, so the same proof gives the stated
universal basis.  The source's identity-pullback comparison
is removed by the coordinate-preserving chart isomorphism
of GSL1; the resulting $b_i$ belong to the literal
\texttt{twist\allowbreak LB}, and (GSL.18) is its exact \texttt{Is\allowbreak Sec\allowbreak Basis}.
\end{proof}
% END GSL24

% BEGIN GSL25
\paragraph{GSL25: every residue section lifts, giving the full \texttt{Sec\allowbreak Twist\allowbreak St}.}
\label{claim:GSL25}
Under GSL24, let $\overline X=\operatorname{Spec}\kappa\times_S A$
and $\rho_X:\overline X\to X$ be the prescribed map induced by
the residue homomorphism $R\to\kappa$.  For \emph{every} actual
section $\overline s$ of the literal bundle
\[
 \operatorname{LinComb}\!\left(
   L,\ [\,\mathrm{snd}\circ\rho_X,
              \mathrm{translPt}(c)\circ\rho_X\,],\
          [\,q-1,1\,]\right),
 \tag{GSL.19}
\]
where composition has its usual right-to-left meaning,
there is a section $s$ of
$\operatorname{twistLB}(A,L,q,c)$ with
$\operatorname{IsPullSecLC}(s,\overline s)$.
Together with GSL24 this is the complete original
\texttt{Sec\allowbreak Twist\allowbreak St}, with no condition on $S$ and no
Noetherian, reduced or flatness condition on the local $R$.

\begin{proof}
The columns $B$ constructed in GSL24 form a basis after the
residue map, because their denominator was a unit of $R$.
The residue homomorphism is surjective.  Apply GSL13 to
this same matrix and the actual bundle $M$.  It lifts every
$\overline s$, with exactly the chart formula (GSL.7).
Equivalently, pull the basis $b_i$ to obtain $\overline b_i$,
write uniquely $\overline s=\sum_i\overline a_i\overline b_i$,
choose arbitrary lifts $a_i\in R$, and set $s=\sum_i a_i b_i$.
Pullback preserves this sum and its coefficients, giving
$\overline s$ on every chart.

The target bundle (GSL.19) is precisely the original
\texttt{Lin\allowbreak Comb} obtained by composing the two maps with the
whiskered residue map.  GSL1 identifies it with the nested
pullback of $M$ by equality of the local coordinate functions,
so the equality just established is the required
\texttt{Is\allowbreak Pull\allowbreak Sec\allowbreak LC}, not an equality only after an unspecified
line-bundle isomorphism.  No lift of the geometric root $d$
to $R$ is needed: that root was used only to prove geometric
exactness and rank before choosing the matrix $B$.
The original section, point $c$, local ring and residue
section remain unchanged throughout.
\end{proof}
% END GSL25

% BEGIN GSL26
\paragraph{GSL26: the exact conditional installation and the remaining boundary.}
\label{claim:GSL26}
The literal inputs $\mathcal F,\mathcal P,\mathcal V,\mathcal R$
supply both global section base-change laws; adding the original
$\mathcal T$ supplies the full twisted law.  The constructions
above establish this implication on the original arbitrary bases,
all affine target rings and all local coefficient rings, with
the quantified choices and actual section objects of the source.
Equivalently, $\mathcal V,\mathcal R,\mathcal T$ may be supplied
by the exact upstream packages in GSL17, GSL19 and GSL21.
The finite-complex part of $\mathcal F$ is supplied by
GSL5--GSL10 from its original geometric section-cube model.

Consequently these law instances may be installed in the already
proved TDR/NFR/NFL arguments for the \emph{same} supplied
\texttt{Moduli} record and actual normalized family.
This removes no other premise of those arguments and does not
assert existence of that moduli record or of a new common bundle.

\begin{proof}
GSL23 proves each quantifier of \texttt{Sec\allowbreak Base\allowbreak Change\allowbreak Any\allowbreak St};
its restriction proves \texttt{Sec\allowbreak Base\allowbreak Change\allowbreak St}.  GSL24 and
GSL25 prove the two conjuncts of \texttt{Sec\allowbreak Twist\allowbreak St}, with
the same $P,q,R,t,c$ in both.  In GSL23 the open and its
basis precede every affine test map.  In GSL24 the geometric
root is allowed only inside the proof, whereas the resulting
sections lie on the original bundle over $R$; GSL25 lifts
every residue section with that same fixed basis.  These
are precisely the clauses consumed by TDR13 and NFL22.
Substitution of the supplied laws in those proved implications
therefore keeps their original objects and quantifier order.

For the upstream form, GSL15 derives the projectivity premise
from $\mathcal P$, GSL14 supplies the full isomorphism-invariance
premise, and GSL3--GSL4 supply the original universally quantified
matrix premise.  GSL17 supplies $\mathcal V$; GSL19 supplies
$\mathcal R$; GSL20 makes their full original-fiber versions.
GSL21 supplies $\mathcal T$, and GSL22 explains why its
locally Noetherian original scope suffices for an arbitrary
$S$ in the final twisted law.  Finally GSL8--GSL10 give
every field of $\mathcal F$ whenever the geometric
Noetherian model of the actual section cube is provided.

The remaining global suppliers are thus explicitly geometric:
the universal section-cube model, polarization base change,
the Serre and multiplication/descent inputs of GSL17,
the initial rank-one, section-decomposition and growth inputs
of GSL19, and the square input of GSL21.  Their formal proofs
are accepted; their complete human reconstructions are not
claimed by this packet.  The separate supplied \texttt{Moduli}
and inherited HUF/PUN foundations likewise remain.  This is
the strongest complete original-scope conditional section-law
join proved here, and supplies no automatic closure of
\texttt{Is\allowbreak Unf} on unsupplied data or of a selected RH/GRH leaf.
\end{proof}
% END GSL26

\subsection*{Exact source correspondence}

{\small\raggedright
All paths in this paragraph are relative to the pinned repository's
\texttt{src/\allowbreak Development/\allowbreak Abelian\allowbreak Variety\allowbreak Theory/\allowbreak } unless explicitly
stated otherwise.  Write $A$ for
\texttt{Proofs/\allowbreak Abelian\allowbreak Scheme\allowbreak Law\allowbreak Proofs/\allowbreak } and $C$ for
\texttt{Proofs/\allowbreak Coherent\allowbreak Cohomology\allowbreak Proofs/\allowbreak }.  The complete source files,
exact ranges and SHA-256 bindings are provided in this component's
\texttt{source\_\allowbreak manifest.\allowbreak json} and \texttt{provenance.\allowbreak json}.

\begin{itemize}
 \item GSL1, GSL11--GSL12 and GSL23:
 $A$\texttt{Sec\allowbreak Base\allowbreak Change\allowbreak Of\allowbreak Perf.\allowbreak lean}, lines 22--170,
 241--287 and 305--401.  The original complete conditional
 join is \texttt{sec\allowbreak Base\allowbreak Change\_\allowbreak of\_\allowbreak perf}, lines 333--401.
 \item GSL2--GSL4:
 $A$\texttt{Exists\allowbreak Mul\allowbreak Eq\allowbreak Smul\allowbreak One\allowbreak Of\allowbreak Mul\allowbreak Vec\allowbreak Injective.\allowbreak lean},
 lines 9--53, 55--156 and 158--181, respectively.
 \item GSL5:
 $C$\texttt{Exists\allowbreak Matrices\allowbreak Of\allowbreak Projective\allowbreak Complex.\allowbreak lean},
 especially lines 7--52 and 90--170.
 \item GSL6--GSL7:
 $C$\texttt{Exists\allowbreak Finite\allowbreak Projective\allowbreak Complex\allowbreak Quasi\allowbreak Iso.\allowbreak lean},
 lines 14--91, 124--244, 255--348 and 521--624.
 \item GSL8:
 $C$\texttt{Exists\allowbreak Finite\allowbreak Projective\allowbreak Complex\allowbreak Of\allowbreak Noetherian\allowbreak Model.\allowbreak lean},
 lines 25--60.  The geometric Noetherian-model producer retained
 above is $C$\texttt{Dehomogenising\allowbreak Comp\allowbreak Of\allowbreak Saturating.\allowbreak lean},
 lines 44--114; it is not counted as proved here.
 \item GSL9--GSL10:
 $C$\texttt{Three\allowbreak Term\allowbreak Complex\allowbreak Of\allowbreak Projective\allowbreak Cech\allowbreak Model.\allowbreak lean},
 lines 24--156; the common-chart and model inputs also appear in
 $C$\texttt{Very\allowbreak Ample\allowbreak Compact\allowbreak Quasi\allowbreak Separated\allowbreak Affine\allowbreak Charts.\allowbreak lean},
 lines 24--50, and
 $C$\texttt{Exists\allowbreak Finite\allowbreak Projective\allowbreak Cech\allowbreak Complex\allowbreak Line\allowbreak Bundle.\allowbreak lean},
 lines 19--43.  The last geometric producer remains a premise.
 \item GSL13, GSL14 and GSL22--GSL25:
 $A$\texttt{Finrank\allowbreak Of\allowbreak Is\allowbreak Sec\allowbreak Basis.\allowbreak lean}, lines 62--110,
 126--216, 231--310 and 312--421.  In particular
 \texttt{geometric\_\allowbreak fibre} uses the inverse translation,
 and \texttt{sec\allowbreak Twist\_\allowbreak of\_\allowbreak perf} includes every residue lift.
 \item The all-cover isomorphism argument of GSL14 is
 $A$\texttt{First\allowbreak Cohomology\allowbreak Vanishes\allowbreak Of\allowbreak Iso\allowbreak Of\allowbreak Is\allowbreak Gen.\allowbreak lean},
 lines 129--224.  GSL15 is
 $A$\texttt{Very\allowbreak Ample\allowbreak By\allowbreak Of\allowbreak Rel\allowbreak Ample.\allowbreak lean}, lines 29--121.
 \item GSL16--GSL17:
 $A$\texttt{First\allowbreak Cohomology\allowbreak Vanishes\allowbreak Descent.\allowbreak lean},
 lines 65--119 and 123--228.  The retained Serre input is the
 first conjunct in
 $A$\texttt{Serre\allowbreak Vanishing\allowbreak Very\allowbreak Ample\allowbreak Statement.\allowbreak lean}.
 \item GSL18--GSL19:
 $A$\texttt{Finrank\allowbreak Sections\allowbreak Line\allowbreak Bundle\allowbreak Pow\allowbreak Eq\allowbreak Pow.\allowbreak lean},
 lines 13--123 and 125--260.  GSL20 uses the base-change
 section equivalence in
 $A$\texttt{Sec\allowbreak MInvariant\allowbreak Eq\allowbreak Bot\allowbreak Or\allowbreak Univ\allowbreak At\allowbreak Geometric\allowbreak Point.\allowbreak lean},
 lines 222--243 and 304--423, and
 $A$\texttt{Sec\allowbreak Base\allowbreak Change\allowbreak Of\allowbreak Perf.\allowbreak lean}, lines 305--325.
 \item GSL21:
 $A$\texttt{Twisted\allowbreak Pow\allowbreak Translate.\allowbreak lean}, lines 63--121.
 The precise assembly boundary in GSL26 is recorded in
 $A$\texttt{Laws\allowbreak Of\allowbreak H0Eq\allowbreak One\allowbreak Closed\allowbreak Fields.\allowbreak lean}, lines 64--91.
\end{itemize}

The literal contracts are in
\texttt{Definitions/\allowbreak Abelian\allowbreak Scheme\allowbreak Laws/\allowbreak First\allowbreak Cohomology\allowbreak Vanishes\allowbreak Spec\allowbreak Of\allowbreak Is\allowbreak Local\allowbreak Ring.\allowbreak lean},
lines 16--88, and, relative to the repository's \texttt{src/\allowbreak },
\texttt{Reciprocity\allowbreak Theory/\allowbreak Reductive\allowbreak Groups/\allowbreak Arithmetic\allowbreak Geometry/\allowbreak Extension\allowbreak Descent/\allowbreak }:
\texttt{Nonzero\allowbreak Locus.\allowbreak lean}, lines 15--75;
\texttt{Global\allowbreak Sections\allowbreak Rank\allowbreak Alg\allowbreak Closed.\allowbreak lean}, lines 19--41 and 160--166;
and \texttt{Sec\allowbreak Base\allowbreak Change\allowbreak Any\allowbreak St.\allowbreak lean}, lines 16--58.
These give the actual \texttt{Pol\allowbreak AV}, \texttt{Perf}, basis and lifting
definitions used throughout.  No source declaration is counted as a
human proof merely because its formal verification is accepted.
\par}


\section{The original arithmetic entries and their common marked frames}
\label{sec:arithmetic-frame-entries}

\paragraph{The endpoint and its retained premise.}
We prove the two original nineteen-field implications
\(\texttt{gen\_\allowbreak gd}\) and \(\texttt{gen\_\allowbreak host}\), then obtain the
original \(\texttt{Gd\allowbreak Frame\allowbreak Ex}\) and \(\texttt{Host\allowbreak Frame\allowbreak Ex}\) from
\(P:\texttt{GC.\allowbreak Hypotheses}_{2}\).
The latter premise is exactly the cohomological package retained
in LFC22--23; it is not replaced by the assertion that a formal
proof exists.  The coefficient ring in these two entries is the
original coefficient DVR \(O=c.O\).
The Dat modulus is \(\ell N|\operatorname{disc}K|\), with
\(H_b=G_{\mathbb Q}\).  The Host modulus is \(\ell N\), with
\(H_b=\ker\operatorname{dir}(\kappa_B)\).
The Host argument allows \(\ell=3\), and both arguments work at
the valuation specified in the input.

The main simplifications are elementary.  Quadratic Gauss sums
place the imaginary quadratic field in the cyclotomic field of
its discriminant.  The two Dat nondegeneracy fields follow
directly from the original kernel and weight conditions.
A dyadic conductor calculation supplies the Host witness at
every relevant local place, without first changing the place
or assuming an auxiliary Host frame exists.  We also construct
that original auxiliary frame and prove its full
\(\texttt{exists\_\allowbreak sigma}\) consequence.

\paragraph{Notation and foundations.}
Let \(\Omega\) be the specified algebraic closure of \(\mathbb Q\),
and put \(G=\operatorname{Gal}(\Omega/\mathbb Q)\).
For a valuation subring \(B\) above a rational prime \(p\), write
\(D_B,I_B\) for its actual decomposition and inertia subgroups.
We use the usual extension-of-valuations dictionary: \(D_B\)
is the absolute Galois group of the local completion, \(I_B\)
is the kernel of its action on the residue algebraic closure,
and \(D_B/I_B=\widehat{\mathbb Z}\) with arithmetic Frobenius
acting as \(x\mapsto x^p\).
Finite and infinite Galois correspondence, the existence of
these prolongations, elementary integral-basis theory, and
Hensel's lemma for the complete ring \(\mathbb Z_\ell\) are
the classical foundations used below.  The cyclotomic,
character, discriminant and conductor calculations needed
from them are proved explicitly.

For a positive integer \(n\), let
\[
 U_n=(\mathbb Z/n\mathbb Z)^\times,\qquad
 c_n:G\longrightarrow U_n
\]
be the action on \(\mu_n\).  Let
\(\epsilon:G\to\mathbb Z_\ell^\times\) be the inverse-limit
cyclotomic character.  If \(O=c.O\), put
\[
 k=O/\mathfrak m_O,\qquad
 \omega:G\longrightarrow\mathbb F_\ell^\times
                         \longrightarrow k^\times .
\]
The character \(\omega\) is the source's literal
\(\texttt{om\allowbreak O}\).  For a Dirichlet character
\(\psi:U_N\to O^\times\), write
\[
 \operatorname{dir}(\psi)=\psi c_N,\qquad
 \Psi=\epsilon^e\operatorname{dir}(\psi),\qquad
 \overline\Psi=\operatorname{red}\circ\Psi.
                                                        \tag{AFE.1}
\]
The power of \(\epsilon\) in (AFE.1) uses the specified
\(\mathbb Z_\ell\)-algebra map to \(O\).
All finite characters are given the discrete finite-image
topology.  Their global pullbacks are locally constant.

\subsection{Residue fields, character components and cyclotomic inertia}

% BEGIN AFE1
\paragraph{AFE1. The coefficient residue field and torsion reduction.}
\label{claim:AFE1}
For the original coefficient \(c\), the field \(k\) is finite
of characteristic \(\ell\).  Reduction is injective on roots
of unity of order prime to \(\ell\), and kills roots of unity
of \(\ell\)-power order.  The reduction of the specified map
\(\mathbb Z_\ell\to O\) is the usual map
\(\mathbb Z_\ell\to\mathbb F_\ell\to k\).

\begin{proof}
The ring \(O\) is finite over \(\mathbb Z_\ell\) and is a
subring of \(\overline{\mathbb Q}_\ell\).
Every \(x\in O\) is therefore integral over \(\mathbb Z_\ell\).
If \(|x|_\ell>1\), a monic integral equation for \(x\) has
its leading term strictly larger in absolute value than
every other term, which is impossible.  Hence
\(|x|_\ell\le1\).  In particular \(\ell^{-1}\notin O\), so
\(\ell\) is a nonunit and lies in \(\mathfrak m_O\).
The quotient \(O/\ell O\) is a finite-dimensional
\(\mathbb F_\ell\)-space, and \(k\) is a quotient of it.
Thus \(k\) is finite and has characteristic \(\ell\).
For \(z\in\mathbb Z_\ell\), choose its integer residue
representative \(a\).  The difference \(z-a\) belongs to
\(\ell\mathbb Z_\ell\), so its image in \(k\) vanishes.
This proves the exact reduction convention.

More generally, let \(R\) be any local ring with residue
characteristic \(\ell\).  If \(x^n=1\), \(\ell\nmid n\),
and \(\bar x=1\), then
\[
 (x-1)(1+x+\cdots+x^{n-1})=0.
\]
The second factor reduces to \(n\ne0\) and is a unit.
Thus \(x=1\).  Applying this to the ratio of two roots
proves injectivity on prime-to-\(\ell\) torsion.
In a field of characteristic \(\ell\),
\(X^{\ell^a}-1=(X-1)^{\ell^a}\); this proves the
assertion about \(\ell\)-power roots.

We will also use the elementary fact that a finite subgroup
of the multiplicative group of a field is cyclic.  Indeed,
if \(d\) is its exponent, the primary decomposition of a
finite abelian group produces an element of order \(d\).
Every element is a root of \(X^d-1\), so the subgroup has
at most \(d\) elements and hence is generated by that
element.  In particular all the finite character orders
used below are positive.
\end{proof}
% END AFE1

% BEGIN AFE2
\paragraph{AFE2. The exact modulus and image-primary decompositions.}
\label{claim:AFE2}
Write \(N=\prod_p p^{a_p}\).
Let \(\operatorname{pr}_p:U_N\to U_N\) retain the
\(p^{a_p}\)-component in the Chinese remainder decomposition
and set every other component to \(1\).  Then
\[
 \psi_p=\psi\operatorname{pr}_p=\texttt{comp }p\,\psi,\qquad
 \psi=\prod_{p\mid N}\psi_p.                              \tag{AFE.2}
\]
The conductor of \(\psi_p\) is \(p^{b_p}\), with
\(0\le b_p\le a_p\).  In particular the dyadic character
\(\psi_2\) has \(2\)-power order and conductor
\(2^{c_2(\psi)}\).

The source's \(\kappa=\texttt{kap }\ell\,\psi\) and
\(\theta=\texttt{th0 }\ell\,\psi\) are exactly the
\(\ell\)-primary and prime-to-\(\ell\) parts of the character:
\[
 \psi=\theta\kappa,\quad
 \kappa_L=\kappa_\ell,\quad
 \kappa_B=\kappa\kappa_L^{-1},\quad
 \psi=\theta\kappa_B\kappa_L.                             \tag{AFE.3}
\]
For odd \(\ell\), \(\kappa_2=1\), and consequently
\(\psi_2=\theta_2\).

\begin{proof}
Put \(A=p^{a_p}\) and \(C=N/A\).  The integer
\(C^{\varphi(A)}\) is \(1\) modulo \(A\), and is \(0\)
modulo \(C\); when \(A=1\) it is \(0\) modulo \(N\).
Thus the source's idempotent
\(\texttt{idem }p\,N\) has exactly those two components.
The formula
\[
 u\longmapsto \texttt{idem }p\,N\cdot u+
                         (1-\texttt{idem }p\,N)
\]
is precisely \(\operatorname{pr}_p\).  The product of these
projections is \(u\), proving (AFE.2), as well as
idempotence, commutation of different projections, and
compatibility with products and inverses of characters.
A character depending only on the \(p\)-component factors
through the least quotient \(U_{p^{b_p}}\) on which it is
well defined.  This is its conductor: triviality on the
kernel is exactly the condition for factoring through a
surjective group map.  Its conductor therefore divides
\(p^{a_p}\).  The group \(U_{2^{a_2}}\) has \(2\)-power
order, proving the dyadic assertion.

For a finite-order element \(x\), write its order as
\(\ell^a d\), with \(\ell\nmid d\).
The exponent \(d^{\varphi(\ell^a)}\) is \(1\) modulo
\(\ell^a\) and is \(0\) modulo \(d\).  Thus
\[
 x^{\,d^{\varphi(\ell^a)}}
\]
is exactly the \(\ell\)-primary part of \(x\).
This is the literal \(\texttt{l\allowbreak Part}\) formula, including
the case \(a=0\).  Apply it in the finite character group
to obtain \(\kappa\); its complementary factor is
\(\theta=\psi\kappa^{-1}\).
Projection to modulus components commutes with taking
these powers, giving (AFE.3).
Finally \(\kappa_2\) has both \(\ell\)-power and
\(2\)-power order.  For odd \(\ell\) it is trivial.
\end{proof}
% END AFE2

% BEGIN AFE3
\paragraph{AFE3. Cyclotomic inertia and corrected Frobenius at every place.}
\label{claim:AFE3}
For every \(B\) above \(p\), the image of \(I_B\) under
\(c_N\) is precisely \(\operatorname{pr}_p(U_N)\).
If the integer \(m\) is a unit in a valuation subring,
its inertia fixes \(\mu_m\).
For every \(B\) above \(\ell\),
\(\epsilon(I_B)=\mathbb Z_\ell^\times\).
For each \(B\) above \(p\) there is \(\gamma\in D_B\)
whose action on every root of unity of order prime to
\(p\) is \(\zeta\mapsto\zeta^p\).  If \(p=\ell\), it
can be chosen with \(\epsilon(\gamma)=1\).
If \(p\ne\ell\) and \(p\equiv1\pmod\ell\), an ordinary
Frobenius lift already has \(\omega(\gamma)=1\).

\begin{proof}
A root of unity and its inverse are integral at every
valuation, so both have valuation \(1\).
If inertia sends a root \(\zeta\in\mu_m\) to \(\zeta'\),
their residues agree.  The ratio \(x=\zeta'/\zeta\)
satisfies \(x^m=1\) and has residue \(1\).
When \(m\) is a unit the geometric-sum argument in
AFE1, applied to this valuation ring, gives \(x=1\).
Thus inertia fixes all prime-to-\(p\) roots at \(B\).

For \(a\ge1\), the polynomial
\(\Phi_{p^a}(X+1)\) is Eisenstein at \(p\).
Its reduction is \(X^{p^{a-1}(p-1)}\), its constant
term is \(p\), and its leading coefficient is \(1\).
It follows that \(\mathbb Q_p(\mu_{p^a})/\mathbb Q_p\)
has degree \(\varphi(p^a)\) and is totally ramified.
The faithful cyclotomic action therefore identifies its
Galois group with \(U_{p^a}\), and inertia surjects onto
that group.  Together with the preceding prime-to-\(p\)
assertion and the Chinese remainder theorem, this proves
the exact image \(\operatorname{pr}_p(U_N)\).
The argument takes place in the completion belonging to
the given \(B\), so it applies to every prolongation.

The compatible maps \(I_B\to U_{\ell^a}\) are all onto.
For any specified element of \(\mathbb Z_\ell^\times\),
the sets of its lifts at levels \(a\) are nonempty nested
closed subsets of the compact group \(I_B\).
Their intersection is nonempty.  This proves
\(\epsilon(I_B)=\mathbb Z_\ell^\times\).

Choose a lift \(\gamma\in D_B\) of arithmetic Frobenius
in \(D_B/I_B=\widehat{\mathbb Z}\).
For a prime-to-\(p\) root \(\zeta\), the elements
\(\gamma\zeta\) and \(\zeta^p\) have the same residue;
injectivity of reduction on those roots gives equality.
This holds simultaneously for all such roots.
If \(p=\ell\), choose \(\tau\in I_B\) with
\(\epsilon(\tau)=\epsilon(\gamma)^{-1}\) and replace
\(\gamma\) by \(\gamma\tau\).  Inertia acts trivially
on the prime-to-\(p\) roots, so their Frobenius action
is unchanged and the new cyclotomic character is \(1\).
If \(p\ne\ell\), Frobenius acts on \(\mu_\ell\)
by the exponent \(p\), which proves the last assertion.
\end{proof}
% END AFE3

% BEGIN AFE4
\paragraph{AFE4. A quadratic discriminant square root with the exact sign character.}
\label{claim:AFE4}
Let \(K\) be the original imaginary quadratic field,
embedded in \(\Omega\), and let
\(\eta:G\to\{\pm1\}\) have kernel \(G_K\).
Write \(d=\operatorname{disc}K\) and \(D=|d|\).
Then \(D>0\), and there is \(\delta\in K\subset\Omega\)
such that
\[
       \delta^2=-D,\qquad s\delta=\eta(s)\delta
                          \quad(s\in G).                \tag{AFE.4}
\]

\begin{proof}
Write \(K=\mathbb Q(\sqrt m)\) for a negative squarefree
integer \(m\).  Here is the integral-basis calculation
including the possible dyadic factor.
If \(a+b\sqrt m\) is integral, its trace gives
\(2a\in\mathbb Z\), and its norm gives
\(4mb^2\in\mathbb Z\).
Squarefreeness of \(m\), prime by prime in the rational
denominator of \(b\), implies \(2b\in\mathbb Z\).
Set \(r=2a,s=2b\).  Integrality is then equivalent to
\[
                  r,s\in\mathbb Z,\qquad
                  r^2-ms^2\equiv0\pmod4.
\]
If \(m\equiv1\pmod4\), the two integers have the same
parity; otherwise both are even.  Conversely these
conditions make trace and norm integral and hence give
a monic quadratic equation.  It follows that
\[
 \mathcal O_K=
 \begin{cases}
  \mathbb Z[(1+\sqrt m)/2],&m\equiv1\pmod4,\\
  \mathbb Z[\sqrt m],&m\equiv2,3\pmod4.
 \end{cases}
\]
For the resulting integral basis \((1,u)\), the determinant
of the embedding matrix is \(u'-u\), where the prime denotes
the other embedding.  Thus
\[
        d=(u-u')^2=
        \begin{cases}m,&m\equiv1\pmod4,\\4m,&\text{otherwise}.
        \end{cases}
\]
In particular \(d<0\).  Take \(\delta=u-u'\).
It is nonzero and generates \(K\) over \(\mathbb Q\).
An element of \(G\) fixing \(K\) fixes \(\delta\);
one not fixing \(K\) exchanges its two roots
\(\delta,-\delta\).
These alternatives agree exactly with
\(\eta=1,-1\), since \(\ker\eta=G_K\).
This proves (AFE.4).
\end{proof}
% END AFE4

% BEGIN AFE5
\paragraph{AFE5. The discriminant cyclotomic containment and its inertia consequence.}
\label{claim:AFE5}
For the field and character in AFE4,
\[
                  K\subset\mathbb Q(\mu_D).
                                                        \tag{AFE.5}
\]
Consequently an element fixing \(\mu_D\) lies in
\(\ker\eta\), and if \(\ell\nmid D\) then
\(I_B\subset\ker\eta\) for every \(B\) above \(\ell\).

\begin{proof}
We prove the required quadratic containment directly.
For an odd prime \(p\), choose a primitive \(p\)-th
root \(\zeta\), let \(\chi\) be the quadratic character
of \(\mathbb F_p^\times\), extended by \(0\) at \(0\),
and put
\[
       \tau_p=\sum_{a\in\mathbb F_p}\chi(a)\zeta^a,\qquad
       \tau_p^-=\sum_a\chi(a)\zeta^{-a}.
 \]
The cyclicity in AFE1 gives
\(\chi(-1)=(-1)^{(p-1)/2}\).  Changing \(a\) to \(-a\)
shows \(\tau_p^-=\chi(-1)\tau_p\).
In the product \(\tau_p\tau_p^-\), change variables
\(a=tb\) for \(b\ne0\).  The inner sum
\(\sum_{b\ne0}\zeta^{(t-1)b}\) is \(p-1\) for
\(t=1\) and \(-1\) otherwise.  Since
\(\sum_t\chi(t)=0\), this gives
\[
              \tau_p\tau_p^-=p,\qquad
              \tau_p^2=(-1)^{(p-1)/2}p=:p^*.
                                                        \tag{AFE.6}
\]
In particular these roots are nonzero.

For odd squarefree \(m<0\), the product of the \(p^*\)
over \(p\mid |m|\) has absolute value \(|m|\) and is
congruent to \(1\) modulo \(4\).  It equals \(m\) if
\(m\equiv1\pmod4\), and \(-m\) if \(m\equiv3\pmod4\).
The product of the corresponding \(\tau_p\)'s therefore
gives \(\sqrt m\) inside \(\mathbb Q(\mu_{|m|})\) in
the first case; in the second, multiplying by \(i\)
gives it inside \(\mathbb Q(\mu_{4|m|})\).
The case \(m=-1\) uses the empty product and the same
factor \(i\).
For even squarefree \(m=2u<0\), first obtain
\(\sqrt u\) inside \(\mathbb Q(\mu_{4|u|})\) by the
odd calculation, and multiply by
\(\zeta_8+\zeta_8^{-1}\), whose square is \(2\).
This puts \(\sqrt m\) inside
\(\mathbb Q(\mu_{8|u|})\).
By the discriminant calculation in AFE4, the modulus
in each case is exactly \(D\).
This proves (AFE.5), including all dyadic cases.

An automorphism fixing \(\mu_D\) therefore fixes \(K\).
When \(\ell\nmid D\), AFE3 says that every inertia
group above \(\ell\) fixes \(\mu_D\).
The equality \(\ker\eta=G_K\) proves both consequences.
\end{proof}
% END AFE5

% BEGIN AFE6
\paragraph{AFE6. One inertia exponent, and the two Dat nondegeneracy fields.}
\label{claim:AFE6}
For every \(\psi\) there is \(a\ge0\), chosen independently
of the valuation \(B\), such that
\[
       \overline\Psi|_{I_B}=\omega^{e+a}|_{I_B}
                       \quad(B\mid\ell).                 \tag{AFE.7}
\]
If \(\ell\nmid\operatorname{ord}(\psi)\), the full
\(\ell\)-modulus component has the exact form
\[
       \operatorname{dir}(\psi_\ell)=\widetilde\omega^a,
                                                        \tag{AFE.8}
\]
where \(\widetilde\omega\) is the Teichmuller lift of
\(\omega\) through \(\mathbb Z_\ell^\times\).
For the original Dat input, these facts imply
\[
 \ker\omega\cap\ker\overline\Psi\not\subset\ker\eta,\qquad
 \forall B\mid\ell\ \exists s\in I_B:
                  \eta(s)=1,\quad\overline\Psi(s)^2\ne1.
                                                        \tag{AFE.9}
\]

\begin{proof}
By AFE3, restriction to \(I_B\) sees only \(\psi_\ell\).
The kernel of \(U_{\ell^b}\to\mathbb F_\ell^\times\)
is an \(\ell\)-group, and has trivial image in \(k^\times\)
by AFE1.  Thus the reduced character factors through
\(\mathbb F_\ell^\times\).  Every root of
\(X^{\ell-1}-1\) in \(k\) is in its prime field:
the \(\ell-1\) nonzero prime-field elements already
exhaust the roots.  The cyclicity of
\(\mathbb F_\ell^\times\) consequently writes this
character as \(x\mapsto x^a\).
This is a statement about the original component
\(\psi_\ell\), so one exponent works for all \(B\).
The exact cyclotomic reduction in AFE1 proves (AFE.7).

If the order of \(\psi\) is prime to \(\ell\), the wild
kernel also has trivial image before reduction:
its image has both \(\ell\)-power and prime-to-\(\ell\)
order.  Hence \(\psi_\ell\) itself factors through
\(\mathbb F_\ell^\times\).
For each \(u\in\mathbb F_\ell^\times\), Hensel's lemma
applied to \(X^{\ell-1}-1\) gives a unique lift
\([u]\in\mathbb Z_\ell^\times\) of order dividing
\(\ell-1\).  Uniqueness proves \([uv]=[u][v]\).
The roots \(\psi_\ell(u)\) and \([u]^a\) have the same
residue and order prime to \(\ell\), so AFE1 makes
them equal.  This proves (AFE.8).

Now retain the original Dat and all its GD fields.
The order hypothesis gives \(\theta=\psi\), so its
character \(\chi_P\) is precisely
\(\psi\psi_\ell^{-1}\).
The GD kernel hypothesis supplies \(t\in G\) with
\[
      c_\ell(t)=1,\qquad
      \operatorname{dir}(\chi_P)(t)=1,\qquad
      t\notin G_K.
 \]
By (AFE.8), \(\operatorname{dir}(\psi_\ell)(t)=1\).
Thus \(\operatorname{dir}(\psi)(t)=1\) and
\(\overline\Psi(t)=1\), while
\(\eta(t)\ne1\).  This proves the first assertion
of (AFE.9), with the literal \(\omega\)-kernel since
\(\mathbb F_\ell\to k\) is injective.

The GD condition \(\psi_\ell^2=1\), together with
surjectivity of \(I_B\to\mathbb F_\ell^\times\), gives
\(\ell-1\mid2a\).
Its weight condition is \(\ell-1\nmid2e\); hence
\(\ell-1\nmid2(e+a)\).
Choose \(s\in I_B\) whose \(\omega(s)\) generates
\(\mathbb F_\ell^\times\).  Equation (AFE.7) gives
\(\overline\Psi(s)^2\ne1\).
The original coprimality
\(\ell\nmid D\,h_K\) implies \(\ell\nmid D\),
so AFE5 gives \(\eta(s)=1\).
This proves the second assertion at every \(B\).
Only the Dat branch uses this weight condition;
it imposes no restriction on the Host prime.
\end{proof}
% END AFE6

\subsection{The full Dat entry}

% BEGIN AFE7
\paragraph{AFE7. All nineteen original fields of \(\texttt{gen\_\allowbreak gd}\).}
\label{claim:AFE7}
Let the input be exactly the original
\(\texttt{Dat}(\Omega,K,c,e,\psi)\), with its specified
embedding of the imaginary quadratic \(K\), and let
\(B\) be any valuation above \(\ell\).
Then
\[
 \texttt{Gen}\bigl(\ell,\eta,\Psi,\mathrm{cc},B,
                     \ell N|\operatorname{disc}K|,G\bigr)
                                                        \tag{AFE.10}
\]
holds.  In particular this is the full original
\(\texttt{gen\_\allowbreak gd}\), with \(H_b=\top\).

\begin{proof}
Put \(D=|\operatorname{disc}K|\) and \(M=\ell ND\).
We verify the actual nineteen fields.

\emph{The fields \(\texttt{odd},\texttt{fin},
\texttt{M\_\allowbreak ne},\texttt{l\_\allowbreak dvd},\texttt{two\_\allowbreak dvd},
\texttt{cc},\texttt{over\_\allowbreak l}\).}
The prime is odd by GD.  AFE1 gives finiteness of \(k\).
The integers \(\ell,N,D\) are positive, so \(M\ne0\)
and \(\ell\mid M\).  GD says \(N\) is even, hence
\(2\mid M\).  The Dat embedding into \(\mathbb C\)
realizes its specified \(\mathrm{cc}\) as complex
conjugation.  The valuation condition is the given one.

\emph{The field \(\texttt{quad}\).}
Use \(D,\delta\) from AFE4.  Every prime dividing
\(D\) divides \(M\), and (AFE.4) is the exact
required square and transformation law.

\emph{The field \(\texttt{cyc}\).}
An element fixing \(\mu_M\) fixes \(\mu_D\), and
therefore has \(\eta=1\) by AFE5.
It also fixes \(\mu_N\) and \(\mu_\ell\), so
\(\operatorname{dir}(\psi)=1\) and \(\omega=1\).
By (AFE.1) and AFE1, \(\overline\Psi=1\).
Thus the required two \(\texttt{Kills\allowbreak Mu}\) clauses
hold with exponent \(1\).

\emph{The fields \(\eta_\ell,\texttt{even},\Psi_\ell\).}
The first is AFE5 using \(\ell\nmid D\).
Complex conjugation acts by inversion on every root
of unity, so
\[
       \epsilon(\mathrm{cc})=-1,\qquad
       \operatorname{dir}(\psi)(\mathrm{cc})=\psi(-1).
 \]
The original parity condition
\(\psi(-1)=(-1)^e\) gives \(\Psi(\mathrm{cc})=1\)
even before reduction.  The common exponent
\(j=e+a\) in AFE6 is precisely \(\Psi_\ell\).

\emph{The fields \(\texttt{g4},\texttt{g5}\).}
These are (AFE.9).  Its second witness lies in
\(I_B\subset D_B\), so it satisfies the decomposition
condition at the specified \(B\).

\emph{The four \(H_b\) fields.}
The subgroup \(G\) is open, contains the commutator
and every inertia group, and has index
\(1=\ell^0\).  These give
\(\texttt{Hb\_\allowbreak open},\texttt{Hb\_\allowbreak ab},\texttt{Hb\_\allowbreak idx},
\texttt{Hb\_\allowbreak l}\) literally.

\emph{The field \(\texttt{teich}\).}
Let \(B'\) be any valuation not above \(\ell\)
and \(s\in I_{B'}\).  The integer \(\ell\) is
a unit in this valuation ring, so AFE3 shows
that \(s\) fixes every \(\ell\)-power root.
Thus \(\epsilon(s)=1\) and
\(\Psi(s)=\operatorname{dir}(\psi)(s)\).
The latter has order prime to \(\ell\) by GD.
If its reduction is \(1\), AFE1 makes it \(1\).
This also covers a trivial valuation, and proves
the exact quantification over all \(B'\).
All nineteen fields have now been supplied.
\end{proof}
% END AFE7

\subsection{The dyadic Host arguments}

% BEGIN AFE8
\paragraph{AFE8. The dyadic conductor forces the Host kernel obstruction.}
\label{claim:AFE8}
Suppose \(\ell\) is odd and \(c_2(\psi)\ge4\).
There exists \(s\in G\) such that
\[
       \omega(s)=1,\qquad \overline\Psi(s)=1,\qquad
                         \eta_4(s)=-1,                 \tag{AFE.11}
\]
where \(\eta_4\) is the quadratic character of
\(\mathbb Q(i)\).  This proves the full Host
\(\texttt{g4}\) obstruction.

\begin{proof}
The conductor of \(\psi_2\) is \(2^{c_2(\psi)}\)
by AFE2, so \(4\mid N\).
If \(\psi_2\) were trivial on the kernel of
\(U_N\to U_4\), it would factor through that
surjective map and have conductor dividing \(4\).
Therefore there is \(v\in U_N\), \(v\equiv1\pmod4\),
with \(\psi_2(v)\ne1\).
This value has order \(2^{b+1}\) for some \(b\ge0\).
For \(w=v^{2^b}\) we have \(w\equiv1\pmod4\)
and \(\psi_2(w)=-1\): in the characteristic-zero
domain \(O\), the only nonidentity element whose
square is \(1\) is \(-1\).
The value \(\psi_2(-1)\) is either \(1\) or \(-1\).
In the first case put \(u=-1\), and in the second
put \(u=-w\).  In both cases
\[
                    u\equiv-1\pmod4,\qquad
                    \psi_2(u)=1.                       \tag{AFE.12}
\]

Choose any prolongation of the \(2\)-adic valuation.
AFE3 supplies \(s\) in its inertia group with
\(c_N(s)=\operatorname{pr}_2(u)\).
As \(\ell\) is odd, the same inertia fixes
\(\mu_\ell\), so \(\omega(s)=1\).
By (AFE.12),
\(\operatorname{dir}(\psi)(s)=\psi_2(u)=1\),
whence \(\overline\Psi(s)=1\).
Its action on \(\mu_4\) has exponent \(-1\),
so it sends \(i\) to \(-i\) and has
\(\eta_4(s)=-1\).
This proves (AFE.11) for every exponent \(e\).
\end{proof}
% END AFE8

% BEGIN AFE9
\paragraph{AFE9. The dyadic power bound from the actual conductor.}
\label{claim:AFE9}
Let \(p>1\) be odd, \(o\ge1\), and suppose
\[
                  8op^2<2^{c_2(\psi)}.                 \tag{AFE.13}
\]
If the dyadic component of \(u\in U_N\) is the
residue of \(p\), then
\[
                         \psi_2(u)^{4o}\ne1.           \tag{AFE.14}
\]
The argument only requires that the character values
lie in a characteristic-zero domain.

\begin{proof}
We first give the principal-unit calculation.
If \(t\ge2\) and \(b\) is odd, repeated squaring gives
\[
             (1+2^t b)^{2^h}
                    =1+2^{t+h}b_h,\qquad b_h\ \text{odd}.
                                                        \tag{AFE.15}
\]
For the induction step, the new odd coefficient is
\(b_h+2^{t+h-1}b_h^2\).
It follows that, modulo \(2^A\) with \(A>t\),
the element \(1+2^t b\) has exact order \(2^{A-t}\).
Indeed (AFE.15) gives this for powers of \(2\);
an odd additional exponent preserves the valuation
of the difference from \(1\), by the geometric sum.
The group of units congruent to \(1\) modulo \(2^t\)
has exactly \(2^{A-t}\) elements.  Thus this element
generates that entire group.

Write \(o=2^j m\) with \(m\) odd, and
\(p^2-1=2^s b\) with \(b\) odd.
Since \(p\) is odd, \(s\ge3\).
Assume, for contradiction, that (AFE.14) fails.
The value \(\psi_2(u)\) has \(2\)-power order;
cancelling the odd factor \(m\) in its order gives
\[
                         \psi_2(u)^{2^{j+2}}=1.
 \]
By (AFE.15), the dyadic component of
\(u^{2^{j+2}}\) is
\[
         p^{2^{j+2}}=(p^2)^{2^{j+1}}
                    =1+2^{s+j+1}b',\qquad b'\ \text{odd}.
 \]
Put \(t=s+j+1\) and \(c=c_2(\psi)\).
The numerical estimate
\[
        2^t=2\cdot2^j\,2^s
                    <2op^2\le8op^2<2^c
 \]
shows \(t<c\).
If \(2^A\) is the exact dyadic divisor of \(N\),
then \(c\le A\) by AFE2.
The principal-unit calculation shows that the dyadic
component of \(u^{2^{j+2}}\) generates
\(\ker(U_{2^A}\to U_{2^t})\).
Our assumed character equality says that \(\psi_2\)
is trivial on this generator, hence on that kernel.
Since \(\psi_2\) is trivial on all other modulus
components, it factors through \(U_{2^t}\).
Its conductor therefore divides \(2^t\), contrary
to \(c>t\).
This proves (AFE.14) with the actual conductor,
not merely an upper bound on it.
\end{proof}
% END AFE9

% BEGIN AFE10
\paragraph{AFE10. A nonquadratic decomposition witness at every Host place.}
\label{claim:AFE10}
Retain the original
\(\texttt{Host\allowbreak Data}(\ell,e,\psi,o)\).
For every rational prime \(p\in T_b(\ell,\psi)\)
and every valuation \(B\) above \(p\), there exists
\(x\in D_B\) with
\[
 \eta_4(x)=1,\qquad c_\ell(x)=1,\qquad
 \overline\Psi(x)^2\ne1,\qquad
                         \Psi(x)^2-1\in O^\times.
                                                        \tag{AFE.16}
\]
There is no restriction \(\ell>3\).

\begin{proof}
We first spell out which primes occur.
The set \(T_b\) is \(\{\ell\}\cup S_b\), where
\[
 S_b=\{p\mid N:p\ne\ell,\ \kappa_p\ne1\}.
 \]
For \(p\in S_b\), the nontrivial character
\(\kappa_p\) has \(\ell\)-power order, so
\(\ell\mid |U_{p^{a_p}}|=p^{a_p-1}(p-1)\).
Since \(p\ne\ell\), this gives
\(\ell\mid p-1\).  In particular \(p\) is odd and
\(p\equiv1\pmod\ell\).
The other possibility is \(p=\ell\), which is
also odd by HostData.

By AFE2, \(\psi_2=\theta_2\).
The HostData order field states that each odd
modulus component \(\theta_r\) has order dividing
\(2o\).  Applying (AFE.2) gives the character identity
\[
                         \theta^{4o}=\psi_2^{4o}.
                                                        \tag{AFE.17}
\]
Choose \(\gamma\in D_B\) by AFE3.  If \(p=\ell\),
use the corrected Frobenius with \(\epsilon(\gamma)=1\);
if \(p\ne\ell\), use \(p\equiv1\pmod\ell\).
In both cases \(c_\ell(\gamma)=1\), and for
\(u=c_N(\gamma)\) the dyadic component is \(p\).
The field \(\texttt{R}\) in HostData supplies (AFE.13).

Suppose that \(\overline{\psi(u)}^{\,4}=1\).
The character \(\kappa\) reduces to \(1\) by AFE1,
so \(\overline{\theta(u)^4}=1\).
The value \(\theta(u)^4\) has order prime to \(\ell\);
AFE1 gives \(\theta(u)^4=1\).
Equation (AFE.17) would then give
\(\psi_2(u)^{4o}=1\), contradicting AFE9.
Consequently
\[
             \overline{\psi(u)}^{\,4}\ne1.
 \]
Set \(x=\gamma^2\).
It still belongs to \(D_B\);
the sign character gives \(\eta_4(x)=1\), and
\(c_\ell(x)=1\).
Using \(\omega(\gamma)=1\) in (AFE.1) yields
\[
            \overline\Psi(x)^2
                     =\overline{\psi(u)}^{\,4}\ne1.
 \]
In a local ring an element is a unit precisely when
its residue is nonzero.  Therefore
\(\Psi(x)^2-1\) is a unit.  The construction was
performed in the specified \(D_B\); no change of
valuation or globally chosen substitute was made.
\end{proof}
% END AFE10

% BEGIN AFE11
\paragraph{AFE11. The literal Host subgroup and the away-from-\(\ell\) implication.}
\label{claim:AFE11}
For odd \(\ell\), define
\[
                   H_b=\ker\operatorname{dir}(\kappa_B).
                                                        \tag{AFE.18}
\]
Then \(H_b\) is open, contains the commutator of \(G\),
has index a power of \(\ell\), and contains every
inertia group above \(\ell\).
For every valuation \(B'\) not above \(\ell\),
\[
 s\in I_{B'}\cap H_b,\quad\overline\Psi(s)=1
                         \ \Longrightarrow\ \Psi(s)=1.
                                                        \tag{AFE.19}
\]

\begin{proof}
The character \(\kappa_B\) has finite \(\ell\)-power
order and factors through \(U_N\).
Its global pullback is locally constant with finite
abelian \(\ell\)-power image.
Its kernel is therefore open, contains the commutator,
and has index equal to the cardinality of that image,
which is \(\ell^a\) for some \(a\ge0\).
The \(\ell\)-modulus component of \(\kappa_B\) is
\(\kappa_\ell\kappa_L^{-1}=1\).
By AFE3 its pullback is trivial on every \(I_B\)
above \(\ell\), proving the inertia containment.

For \(s\in I_{B'}\) as in (AFE.19), AFE3 gives
\(\epsilon(s)=1\) and
\(\operatorname{dir}(\kappa_L)(s)=1\):
the first follows from all \(\ell\)-power roots,
and the second depends on only a finite level of them.
Membership in \(H_b\) kills \(\kappa_B\).
Equation (AFE.3) therefore gives
\[
                         \Psi(s)=\operatorname{dir}(\theta)(s).
 \]
The right side has order prime to \(\ell\).
Its residue being \(1\), AFE1 makes it \(1\).
The argument uses only that \(\ell\) is a unit in
the valuation ring, so it covers every \(B'\) in
the original field's quantifier.
\end{proof}
% END AFE11

% BEGIN AFE12
\paragraph{AFE12. All nineteen original fields of \(\texttt{gen\_\allowbreak host}\).}
\label{claim:AFE12}
For the original HostData, any specified complex
conjugation \(\sigma\), and any specified \(B\) above
\(\ell\), the subgroup (AFE.18) satisfies
\[
        \texttt{Gen}(\ell,\eta_4,\Psi,\sigma,B,\ell N,H_b).
                                                        \tag{AFE.20}
\]
This proves the full original existential
\(\texttt{gen\_\allowbreak host}\), including \(\ell=3\).

\begin{proof}
Set \(M=\ell N\) and keep precisely the subgroup
(AFE.18).  We again check all nineteen fields.

\emph{The first seven fields.}
HostData gives \(\ell\ne2\) and even \(N\).
AFE1 gives finite \(k\).
As \(\ell\) is prime and \(N>0\),
\(M\ne0\), \(\ell\mid M\), and \(2\mid M\).
The supplied \(\sigma\) and \(B\) give
\(\texttt{cc}\) and \(\texttt{over\_\allowbreak l}\).

\emph{The field \(\texttt{quad}\).}
Choose \(i\in\Omega\) with \(i^2=-1\), and use the
positive integer \(D=1\).
The prime-divisor condition is vacuous.
The action on \(\mu_4\) is \(i\mapsto i\) or
\(i\mapsto-i\); by the definition of the character
\(\chi_4:U_4\to\{\pm1\}\), this is
\(s(i)=\eta_4(s)i\).

\emph{The field \(\texttt{cyc}\).}
Use exponent \(2\).
Since \(M\) is even, \(4\mid M^2\);
also \(N\mid M^2\) and \(\ell\mid M^2\).
Fixing \(\mu_{M^2}\) therefore kills \(\eta_4\),
\(\operatorname{dir}(\psi)\), and \(\omega\).
Equation (AFE.1) gives the required two
\(\texttt{Kills\allowbreak Mu}\) conclusions.

\emph{The fields \(\eta_\ell,\texttt{even},\Psi_\ell\).}
As \(\ell\) is odd, AFE3 says that inertia at
\(\ell\) fixes \(\mu_4\), proving \(\eta_\ell\).
Complex conjugation acts by inversion on roots;
the original parity condition again gives
\(\Psi(\sigma)=(-1)^e\psi(-1)=1\).
AFE6 supplies a single \(j=e+a\) before all
valuations, which is exactly \(\Psi_\ell\).

\emph{The fields \(\texttt{g4},\texttt{g5}\).}
HostData's \(c_2(\psi)\ge4\) and AFE8 prove
\(\texttt{g4}\).
The prime \(\ell\) belongs to \(T_b\), so AFE10
at the specified \(B\) proves \(\texttt{g5}\).

\emph{The four \(H_b\) fields and \(\texttt{teich}\).}
They are precisely AFE11.
Thus no \(H_b\) condition has been replaced by the
corresponding easier condition for \(\top\).
The conductor and Frobenius arguments require only
that \(\ell\) be odd, so all these fields apply when
\(\ell=3\).
\end{proof}
% END AFE12

\subsection{The original auxiliary Host frame and its local witnesses}

% BEGIN AFE13
\paragraph{AFE13. Constructing the full original \(\texttt{Is\allowbreak Host\allowbreak Frame}\).}
\label{claim:AFE13}
Under the original HostData there exist the finite
set \(S\), actual prolongations \(A_v\), and the
original frame \(F\) satisfying all ten
\(\texttt{Is\allowbreak Host\allowbreak Frame}\) fields.
This is the full \(\texttt{exists\_\allowbreak host\allowbreak Frame}\);
it does not require \(P\).

\begin{proof}
Identify finite places of \(\mathbb Q\) with rational
prime ideals, and let
\[
       S=\{v:v\mid\ell N\}.
 \]
This is finite because \(\ell N\ne0\).
Choose a prolongation \(A_v\) of each \(v\in S\).
Let \(R_S\) be the closed normal subgroup of \(G\)
generated by all inertia groups at places outside
\(S\), and set
\[
       \Gamma=G/R_S=G_{\mathbb Q,S},\qquad
       \pi:G\twoheadrightarrow\Gamma.
 \]
This is exactly the source's \(\texttt{GS}\) and
\(\texttt{to\allowbreak GS}\), not a different quotient.
An inertia group outside \(S\) fixes \(\mu_N\),
\(\mu_4\), and every \(\ell\)-power root by AFE3:
for \(\mu_4\), use \(2\mid N\).
The kernels of \(c_N,c_4\) are open and closed,
and the kernel of \(\epsilon\) is an intersection
of such closed kernels.  Normality and closedness
therefore show that all three characters kill \(R_S\).
They descend to characters
\[
       \mu_N:\Gamma\to U_N,\quad
       \mu_4:\Gamma\to U_4,\quad
       \epsilon_S:\Gamma\to\mathbb Z_\ell^\times
 \]
whose pullbacks along \(\pi\) are the original ones.

Choose an embedding \(\Omega\hookrightarrow\mathbb C\);
its image is the field of algebraic numbers and is
stable under complex conjugation.
Conjugation therefore pulls back to
\(\sigma\in G\) with \(\texttt{Is\allowbreak CC}\,\sigma\).
Define the actual frame by
\[
 \begin{gathered}
 F.\pi=\pi,\qquad F.\eta=\chi_4\mu_4,\qquad
 F.\Psi=\epsilon_S^e\,\psi\mu_N,\qquad F.c=\pi(\sigma),\\
 F.D_v=D_{A_v},\qquad F.I_v=I_{A_v},\qquad
 F.v_\ell=(\ell),\\
 F.S_b=\{v\in S:\exists p\in S_b(\ell,\psi),\ v=(p)\}.
 \end{gathered}                                         \tag{AFE.21}
\]
The distinguished place is in \(S\), and each prime
in \(S_b(\ell,\psi)\) divides \(N\), so these choices
have the required types.
In the source's frame the groups \(F.D_v,F.I_v\)
are subgroups of the original \(G\); their maps to
\(\Gamma\) are induced by \(\pi\), just as in (AFE.21).

The ten fields can now be read off without additional
choices: membership of \(S\) is exactly divisibility
by \(\ell N\); \(A_v\) is over \(v\);
the quotient map is literally \(\texttt{to\allowbreak GS}\);
the two character pullbacks are \(\eta_4,\Psi\);
\(\sigma\) gives the complex-conjugation field;
the decomposition and inertia groups are the actual
ones; \(\ell\) belongs to the distinguished prime;
and the displayed \(S_b\) is its exact required
prime-divisibility condition.
These are respectively
\(\texttt{mem},\texttt{is\allowbreak Over},\pi,\eta,\Psi,
\texttt{cc},D,I,\texttt{vl},\texttt{Sb}\).
\end{proof}
% END AFE13

% BEGIN AFE14
\paragraph{AFE14. The full original \(\texttt{exists\_\allowbreak sigma}\) on any Host frame.}
\label{claim:AFE14}
Let \(F,S,(A_v)\) be any supplied original
\(\texttt{Is\allowbreak Host\allowbreak Frame}\), not necessarily the one
chosen in AFE13.
For every \(v\in F.T_b=\{F.v_\ell\}\cup F.S_b\)
there is \(x\in F.D_v\) with
\[
          F.\eta(F.\pi x)=1,\qquad
          F.\Psi(F.\pi x)^2-1\in O^\times,\qquad
                              c_\ell(x)=1.              \tag{AFE.22}
\]

\begin{proof}
If \(v=F.v_\ell\), the field \(\texttt{vl}\)
identifies its rational prime with \(\ell\).
If \(v\in F.S_b\), the field \(\texttt{Sb}\)
supplies \(p\in S_b(\ell,\psi)\) with \(v=(p)\).
In either case \(v=(p)\) for \(p\in T_b(\ell,\psi)\).
The field \(\texttt{is\allowbreak Over}\) puts \(A_v\) above
that exact \(p\), so AFE10 gives \(x\in D_{A_v}\)
with (AFE.16).
Use the literal equalities
\[
        F.D_v=D_{A_v},\quad
        F.\eta\circ F.\pi=\eta_4,\quad
        F.\Psi\circ F.\pi=\Psi
 \]
to obtain all three assertions (AFE.22).
AFE1 supplies the residue characteristic even when
it is not separately listed in this source theorem.
Thus this argument proves the full supplied-frame
statement and is compatible with every allowed choice
of the local prolongations.
\end{proof}
% END AFE14

\subsection{Joining the entries to the completed common frame}

% BEGIN AFE15
\paragraph{AFE15. The full original \(\boldsymbol{P\Rightarrow\texttt{Gd\allowbreak Frame\allowbreak Ex}}\).}
\label{claim:AFE15}
Assume the original \(P:\texttt{GC.\allowbreak Hypotheses}_{2}\).
For every original Dat input and every specified
\(B\) above \(\ell\), there are one CM field \(L\),
one \(g_1\in G\), one frame \(\Phi\), and one \(C_0\)
such that, with \(M=\ell N|\operatorname{disc}K|\),
\[
 \operatorname{GalLink}(\eta,\Psi,\mathrm{cc},g_1,B,M,L,\Phi),
 \quad \Phi.\texttt{Ok}(C_0),\quad \Phi.\texttt{Big}(C_0),
 \quad \texttt{Seed\allowbreak Clause}(c,\eta,\Psi,\mathrm{cc},L).
                                                        \tag{AFE.23}
\]
This is the original \(\texttt{Has\allowbreak Frame}\) conclusion
and hence the full \(\texttt{Gd\allowbreak Frame\allowbreak Ex}\).

\begin{proof}
AFE7 supplies the complete original Gen with \(H_b=\top\).
Apply the inherited human character construction CAF7
to that input first.  It chooses a positive modulus
\(Q\) and a seed-producing function before choosing
the eventual field \(L\) and auxiliary prime \(w\).
Apply LFC22 with this \(Q\), the same Gen and the
original \(P\).  It supplies one
\[
               L,\ g_1,\ \Phi,\ C_0,\ w
 \]
with actual Linked data, all twenty-eight Ok fields,
Big for the same \(C_0\), the original Frobenius
witness at \(w\), \(w\equiv1\pmod Q\), and
\(\operatorname{upl}(L)\subset H^{\mathrm{th}}(w)\).
Let \(k_0\) denote its GalLink.
Its abelian and fixed-field clauses and those
conditions at \(w\) are exactly the field conditions
required by CAF7's previously chosen function.

We verify the remaining marked-element condition
on this very \(L,g_1\).
The GalLink field \(\texttt{g0}\) gives
\(g_1\in\operatorname{upl}(L)\) and
\(k_0.p(g_1)=\Phi.g_0\).
Its sign relation and \(\Phi.g_0\in\ker\Phi.\delta\)
give \(\eta(g_1)=1\).
The fixed-field equality in GalLink consequently
gives \(g_1\in G_L\).
Its character relation gives
\[
             \overline\Psi(g_1)=\Phi.\Psi(\Phi.g_0).
 \]
The actual Ok field \(\texttt{base\allowbreak Subgroup\_\allowbreak spec}\)
says that the square of the right side is not \(1\).
Thus \(g_1\) supplies exactly the nonquadratic marker
needed by CAF11 to install the seed.
Apply CAF11, equivalently the final installation in
LFC23, to the already chosen \(L,w\).
It returns the actual SeedClause on this same \(L\);
it chooses neither a new field nor a new marker.
Together with the GalLink, Ok and Big already
obtained, this is (AFE.23).
Every quantified original Dat and every specified
\(B\) were arbitrary, so the full \(\texttt{Gd\allowbreak Frame\allowbreak Ex}\)
follows.
\end{proof}
% END AFE15

% BEGIN AFE16
\paragraph{AFE16. The full original \(\boldsymbol{P\Rightarrow\texttt{Host\allowbreak Frame\allowbreak Ex}}\).}
\label{claim:AFE16}
Assume the same original \(P\).
For every original HostData, every specified complex
conjugation \(\sigma\), and every specified \(B\)
above \(\ell\), there are one CM field \(L\), one
\(g_1\), one \(\Phi\), and one \(C_0\) satisfying
\[
 \operatorname{GalLink}(\eta_4,\Psi,\sigma,g_1,B,\ell N,L,\Phi),
 \quad\Phi.\texttt{Ok}(C_0),\quad\Phi.\texttt{Big}(C_0),
 \quad\texttt{Seed\allowbreak Clause}(c,\eta_4,\Psi,\sigma,L).
                                                        \tag{AFE.24}
\]
This is the full original \(\texttt{Host\allowbreak Frame\allowbreak Ex}\),
with \(\ell=3\) included.

\begin{proof}
AFE12 supplies full Gen with the actual subgroup
\(H_b=\ker\operatorname{dir}(\kappa_B)\) and
the exact modulus \(M=\ell N\).
Use CAF7 to choose its \(Q\) and character data
before the field and prime.
LFC22, at this \(Q\) and this Gen, supplies
the actual linked \(L,g_1,\Phi,C_0,w\), all Ok
fields, Big, the Frobenius witness,
\(w\equiv1\pmod Q\), and the required subgroup
inclusion.  The character construction is being
used at the original coefficient DVR \(O=c.O\);
we do not extend its seed assertion to a general
local coefficient ring.

For completeness the marking check is unchanged
in content but uses the Host character:
\(\texttt{g0}\) gives \(g_1\in\operatorname{upl}(L)\)
and \(k_0.p(g_1)=\Phi.g_0\);
\(\texttt{del}\) gives \(\eta_4(g_1)=1\);
\(\texttt{gal}\) gives \(g_1\in G_L\);
and \(\texttt{psi}\), followed by
\(\texttt{base\allowbreak Subgroup\_\allowbreak spec}\), gives
\(\overline\Psi(g_1)^2\ne1\).
CAF11 therefore installs its previously selected
character seed on this same \(L,w\).
Keeping the GalLink, Ok and Big from the same
construction proves (AFE.24).
No step invokes the Dat weight condition or
requires \(\ell>3\).
\end{proof}
% END AFE16

\subsection{The original marking and the bounded finite-point transport}

The remaining claims make a further, separate advance.
The representation-theoretic input \(\texttt{Fin\allowbreak Pt}\)
is retained, but its transport to \(\texttt{Fin\allowbreak End}\)
is proved, including the actual normalized representation
and its local characters.  In this subsection \(O\) may
be any commutative local ring of residue characteristic
\(\ell\), as in the original transport theorem.
No coefficient-DVR hypothesis is inserted into that
general statement.

% BEGIN AFE17
\paragraph{AFE17. The same \(g_1\) satisfies all original marking fields.}
\label{claim:AFE17}
The \(g_1\) in AFE15 satisfies
\(\texttt{Marks}(\eta,\Psi,\mathrm{cc},g_1)\).
The \(g_1\) in AFE16 satisfies
\(\texttt{Marks}(\eta_4,\Psi,\sigma,g_1)\).
Thus the frame and its marker already have the
original downstream marking, not merely a residual
version of it.

\begin{proof}
The five fields of Marks are
\[
 \sigma^2=1,\qquad \eta(\sigma)=-1,\qquad
 \Psi(\sigma)=1,\qquad \eta(g_1)=1,\qquad
                         \Psi(g_1)^2-1\in O^\times.
                                                        \tag{AFE.25}
\]
Complex conjugation squares to the identity on
\(\Omega\).
In the Dat case it acts nontrivially on the imaginary
quadratic field, so its sign under \(\eta\) is \(-1\).
In the Host case it sends \(i\) to \(-i\), so its
sign under \(\eta_4\) is \(-1\).
The original parity condition gives
\(\Psi(\sigma)=1\) before reduction, as proved in
AFE7 and AFE12.
The common-marker calculation in AFE15 and AFE16
already gives
\(\eta(g_1)=1\) and \(\overline\Psi(g_1)^2\ne1\).
The latter says that \(\Psi(g_1)^2-1\) has nonzero
residue, hence is a unit of the local ring.
These are exactly all five conditions in (AFE.25).
\end{proof}
% END AFE17

\paragraph{The literal ordinary-point and normalized-representation data.}
Let \(w_0=1,w_1=0,w_2=-1\), and let
\(J\) be the \(3\)-by-\(3\) antidiagonal matrix
whose antidiagonal entries are \(1\).
The group \(\texttt{GG}\,3\,A\) is
\(\mathrm{GL}_3(A)\rtimes\{\pm1\}\), where the sign
\(-1\) acts by
\[
                         X\longmapsto J(X^{-1})^tJ.
 \]
For \(r:G\to\texttt{GG}\,3\,A\), put \(R(g)=(r(g)).\mathrm{mat}\).
An original \(\texttt{Ord\allowbreak Pt}\) consists of the following
five conditions: its sign is \(\eta\), and
\[
 R(g)\bmod\mathfrak m_A
   =\operatorname{diag}\bigl(\Psi(g)^{w_i}\bmod\mathfrak m_A\bigr);
 \]
its matrix entries are locally constant modulo every
\(\mathfrak m_A^n\); it has the original normalization
\[
 R(\sigma)=1,\quad
 \bigl(R(g_1)-R(\sigma g_1^{-1}\sigma)\bigr)_{01}
 =\bigl(R(g_1)-R(\sigma g_1^{-1}\sigma)\bigr)_{10}=0;
                                                        \tag{AFE.26}
\]
away from \(\ell\) it kills inertia elements on which
both \(\eta\) and the full \(\Psi\) are \(1\);
and at the specified \(D_B\) there is \(h\equiv1
\pmod{\mathfrak m_A}\) making the matrices upper
triangular on \(\ker\eta\), with inertia diagonals
exactly the images of \(\Psi^{w_i}\).
The last is the literal condition \(\texttt{Cond\allowbreak T}\).

An \(\texttt{NRep}(\Phi,3,A)\) consists of a matrix
representation \(\rho:\Phi.\Gamma_0\to\mathrm{GL}_3(A)\)
and local characters \(\psi_{v,i}:\Phi.G_v\to A^\times\).
It requires
\[
 \begin{gathered}
 \rho(\Phi.\star g)=J\rho(g)^tJ,\\
 \rho(\Phi.g_0)-\rho(\Phi.\star\Phi.g_0)
                         \ \text{is diagonal},\\
 \prod_{i=0}^{2}
   \bigl(\rho(\Phi.\operatorname{loc}_{0,v}(g_i))
                         -\psi_{v,i}(g_i)1\bigr)=0
                       \quad\text{for every }(g_0,g_1,g_2).
 \end{gathered}                                         \tag{AFE.27}
\]
The product is in the displayed order.
Its condition \(\texttt{Is\allowbreak Lift\allowbreak A}\) means adic continuity
of all these entries and characters, residual diagonal
entries and local characters prescribed by
\(\Phi.\Psi^{w_i}\), and zero residual off-diagonal
entries.  These are the source's exact data, rather
than a semisimplified substitute.

% BEGIN AFE18
\paragraph{AFE18. Transporting the actual triangular characters.}
\label{claim:AFE18}
Let \(r:G\to\texttt{GG}\,3\,A\) have sign
\((r(g)).\mathrm{right}=\eta(g)\), the residual
diagonal condition, adic continuity and
\(\texttt{Cond\allowbreak T}\) at subgroups \(D,I\).
Let \(\varphi:H\to G\) be continuous, and suppose
\[
 \eta(\gamma)=1,\quad \eta(\varphi(s))=1,\quad
                         \gamma^{-1}\varphi(s)\gamma\in D
                           \quad(s\in H).
 \]
Then there are characters \(\chi_i:H\to A^\times\)
with the ordered product condition for \(R\varphi\),
adic continuity, residual values
\(\Psi(\varphi(s))^{w_i}\), and
\[
 \chi_i(s)=\alpha(\Psi(\varphi(s))^{w_i})
 \quad\text{if }\gamma^{-1}\varphi(s)\gamma\in I,
                                                        \tag{AFE.28}
\]
where \(\alpha:O\to A\) is the given local coefficient map.
Here \(A\) need only be a commutative local ring;
no flatness or completeness is used.

\begin{proof}
Choose the actual \(h\equiv1\pmod{\mathfrak m_A}\)
from CondT.  Put
\[
 d(s)=\gamma^{-1}\varphi(s)\gamma,\qquad
 T(s)=h^{-1}R(d(s))h,\qquad \chi_i(s)=T(s)_{ii}.
 \]
All relevant signs are \(1\), so \(T\) is a matrix
representation.  Its matrices are upper triangular.
For such matrices the diagonal of a product is the
product of the diagonals: in the sum
\(\sum_j T(s)_{ij}T(t)_{ji}\), terms \(j<i\)
vanish in the first factor and terms \(j>i\)
in the second.  Consequently each \(\chi_i\)
is a homomorphism; evaluating at \(s^{-1}\)
shows its values are units.

Let \(E_j\) be the span of the first \(j\) coordinate
vectors, with \(E_0=0\).
Upper triangular matrices preserve this filtration.
The factor \(T(g_i)-\chi_i(g_i)1\) maps
\(E_{i+1}\) to \(E_i\).
Thus the rightmost factor for \(i=2\) maps \(A^3\)
to \(E_2\), the middle one maps \(E_2\) to \(E_1\),
and the leftmost maps \(E_1\) to \(0\).
Their ordered product is zero.
Since \(\eta(\gamma)=1\),
\[
 R(\varphi(s))=P\,T(s)\,P^{-1},\qquad P=R(\gamma)h.
 \]
Conjugating the product proves the required version
for \(R\varphi\).

Modulo \(\mathfrak m_A\), both \(h\) and \(h^{-1}\)
are \(1\).
The residual condition therefore gives
\(\bar\chi_i(s)=
 \overline{\alpha(\Psi(d(s))^{w_i})}\).
Since \(\Psi\) is a character,
\(\Psi(d(s))=\Psi(\varphi(s))\).
The exact inertia diagonals of CondT give the same
equality before reduction whenever \(d(s)\in I\).
This proves (AFE.28) with the original character,
not only its residue.

Finally
\(T(s)=P^{-1}R(\varphi(s))P\).
After reduction modulo each \(\mathfrak m_A^n\),
this is a fixed matrix expression in the locally
constant entries of \(R\varphi\).
Its diagonal entries are locally constant as well.
This proves precisely the required adic continuity,
without making any topological assumption on
coefficient changes beyond those stated.
\end{proof}
% END AFE18

% BEGIN AFE19
\paragraph{AFE19. The full original ordinary-point-to-\(\texttt{NRep}\) descent.}
\label{claim:AFE19}
Let \(O\) be a commutative local ring with residue
characteristic \(\ell\).
Take an original
\[
 k_0:\operatorname{GalLink}(\eta,\Psi,\sigma,g_1,B,M,L,\Phi),
 \quad \Phi.\texttt{Ok}(C_0),\quad \ell\mid M,\quad\sigma^2=1.
 \]
Let \(A\) be any Noetherian local domain, with a local
\(O\)-algebra map \(\alpha\).
For every original \(\texttt{Ord\allowbreak Pt}\) \(r\), there is
an actual \(\texttt{NRep}(\Phi,3,A)\) satisfying
\(\texttt{Is\allowbreak Lift\allowbreak A}\) and
\[
 \begin{aligned}
 \rho(k_0.p_0(s))&=R(s) &&(s\in G_L),\\
 \psi_{v,i}(t)&=\alpha(\Psi(k_0.j_v(t))^{w_i})
                                      &&(t\in\Phi.I_v).
 \end{aligned}                                         \tag{AFE.29}
\]
The inclusion \(G_L\to G\) is understood in the first
line.  This is the full original
\(\texttt{exists\_\allowbreak nrep3}\), with no extra arithmetic
or representation-theoretic premise.

\begin{proof}
Put \(U=\operatorname{upl}(L)\) and \(p=k_0.p:U\to\Phi.\Gamma\).
The subgroup \(U\) is closed in the profinite \(G\):
it is an intersection of stabilizers of algebraic
elements, and each such stabilizer is open and closed
because the orbit is finite.
Thus \(U\) is compact.
The exact GalLink kernel is the closure of the normal
subgroup generated by those \(t\in U\) for which
\[
 \exists B'\nmid\ell:\quad
 t\in I_{B'},\quad \eta(t)=1,\quad\overline\Psi(t)=1.
                                                        \tag{AFE.30}
\]
Its Teichmuller field gives \(\Psi(t)=1\) for each
generator in (AFE.30), and the away field of OrdPt
then gives \(r(t)=1\).
Hence \(r\) kills their normal subgroup.

We justify passage to its closure.
For the Noetherian local ring \(A\),
\(\bigcap_n\mathfrak m_A^n=0\).
One standard algebraic proof applies Artin--Rees
to the ideal \(I=\bigcap_n\mathfrak m_A^n\subset A\):
for a sufficiently large integer \(c\),
\(I\cap\mathfrak m_A^{c+1}\subset\mathfrak m_A I\).
Since \(I\) lies in every power, \(I=\mathfrak m_A I\);
it is finitely generated, so Nakayama gives \(I=0\).
Therefore
\[
 \{u\in U:R(u)=1\}=
 \bigcap_n\{u\in U:R(u)\equiv1\pmod{\mathfrak m_A^n}\}
 \]
is closed: each set on the right is a fiber of a
locally constant map.  It contains the normal
subgroup just considered, hence its closure.
For \(u\in\ker p\), GalLink's sign relation also gives
\(\eta(u)=\Phi.\delta(pu)=1\).
Both the matrix and the sign of \(r(u)\) are therefore
trivial.  Thus \(\ker p\subset\ker(r|_U)\).

Surjectivity of \(p\) gives a unique homomorphism
\[
                  r':\Phi.\Gamma\to\texttt{GG}\,3\,A,
                         \qquad r'(pu)=r(u).
 \]
A continuous surjection from compact \(U\) to the
Hausdorff profinite \(\Phi.\Gamma\) is a quotient map:
it is closed, so a subset of the target is open
exactly when its inverse image is open.
Applying this to the fibers of each matrix reduction
shows that \(r'\) is adically continuous.
The sign equality \((r'(x)).\mathrm{right}=\Phi.\delta(x)\)
follows by surjectivity from GalLink's sign relation.
Hence
\[
             \rho(g)=(r'(g)).\mathrm{mat}
                         \quad(g\in\Phi.\Gamma_0)
 \]
is a matrix representation.

The GalLink fields at \(\sigma,g_1\) give
\(r'(\Phi.c)=r(\sigma)\) and
\(r'(\Phi.g_0)=r(g_1)\).
Ok gives \(\Phi.\delta(\Phi.c)=-1\) and \(\Phi.c^2=1\);
normalization gives \(R(\sigma)=1\).
In the displayed semidirect product,
\[
 (1,-1)(X,1)^{-1}(1,-1)=(JX^tJ,1).
 \]
Since \(\Phi.\star g=\Phi.c\,g^{-1}\Phi.c\),
this proves the polarization in (AFE.27).
For any matrix \(X=(x_{ij})\),
\[
 X-JX^tJ=
 \begin{pmatrix}
 x_{00}-x_{22}&x_{01}-x_{12}&0\\
 x_{10}-x_{21}&0&x_{12}-x_{01}\\
 0&x_{21}-x_{10}&x_{22}-x_{00}
 \end{pmatrix}.
 \]
Its two entries \(01,10\) control all off-diagonal
entries.  At \(X=R(g_1)\), (AFE.26) makes both zero,
so the normalization field of NRep follows.

For each \(v\), the original
\(\texttt{pl\_\allowbreak conj}\) field provides \(\gamma_v\)
with \(\eta(\gamma_v)=1\) and
\(k_0.\mathrm{pl}_v=\gamma_v B\).
Set \(\varphi=k_0.j_v\), viewed as a map to \(G\).
The fields \(\texttt{j\_\allowbreak range}\) and \(\texttt{j\_\allowbreak ine}\)
identify its image and inertia with the actual
decomposition and inertia at this place.
Conjugation by \(\gamma_v^{-1}\) therefore carries
them into \(D_B,I_B\).
Furthermore GalLink's sign and localization equalities
give
\[
           \eta(k_0.j_v(t))
             =\Phi.\delta(\Phi.\operatorname{loc}_v(t))=1,
 \]
since this localization has its specified lift to
\(\Phi.\Gamma_0\).
Apply AFE18 to obtain \(\psi_{v,i}\).
GalLink's localization equality and the definition of
\(r'\) give
\[
          \rho(\Phi.\operatorname{loc}_{0,v}(t))
                           =R(k_0.j_v(t)).
 \]
Thus AFE18 gives the ordered product field in
(AFE.27), continuity, and the second line of (AFE.29).

It remains to check the literal residual fields.
The local coefficient map induces
\(k\to k_A=A/\mathfrak m_A\), characterized by
\(\operatorname{red}_A\alpha=
 (k\to k_A)\operatorname{red}_O\).
For a lift \(u\in U\) of \(g\in\Phi.\Gamma_0\),
OrdPt gives the residual diagonal
\(\alpha(\Psi(u)^{w_i})\), and GalLink gives
\(\Phi.\Psi(pu)=\overline\Psi(u)\).
These are exactly the prescribed diagonal entries;
the off-diagonal entries vanish by the same OrdPt
condition.  AFE18 gives the corresponding local
character identities.  All weights are the literal
\(\texttt{w}\,3\,i=1-i=w_i\).
Together with the adic continuity already proved,
these are all fields of IsLiftA.
Finally \(k_0.p_0\) is the restriction of \(p\)
to \(G_L\), so \(r'(pu)=r(u)\) gives the first
line of (AFE.29).
\end{proof}
% END AFE19

% BEGIN AFE20
\paragraph{AFE20. The full original \(\texttt{Fin\allowbreak Pt}\Rightarrow\texttt{Fin\allowbreak End}\) transport.}
\label{claim:AFE20}
At the exact GalLink, Ok, residue-characteristic,
\(\ell\mid M\), and \(\sigma^2=1\) premises of
AFE19, the original \(\texttt{Fin\allowbreak Pt}(O,\Phi,3)\)
implies
\[
                    \texttt{Fin\allowbreak End}(\ell,\eta,\Psi,\sigma,g_1,B).
                                                        \tag{AFE.31}
\]
The finite-index subgroup can be chosen as
\(G_L\) before every coefficient ring and every
ordinary point.

\begin{proof}
Choose \(U_0=G_L\subset G\) once.
Its index is \([L:\mathbb Q]\), a positive finite
integer.  This choice depends only on the given
frame and field.
Now take any original \(\texttt{Is\allowbreak CD}(O,A)\)
coefficient algebra and its local map
\(\alpha:O\to A\), with
\(\alpha(\mathfrak m_O)=0\), and any ordinary point
\(r:G\to\texttt{GG}\,3\,A\).
Here IsCD means a complete DVR domain \(A\), a
local coefficient map, and finite residue field;
in particular \(A\) is a Noetherian local domain.
There is a unique ring map \(f:k\to A\) with
\[
                         f\operatorname{red}_O=\alpha.
 \]
Give \(A\) this \(k\)-algebra structure.
It forms the required scalar tower \(O\to k\to A\).

AFE19 supplies \(\rho\) with IsLiftA and (AFE.29).
Suppose \(t\in\Phi.I_v\) and
\(\Phi.\Psi(\Phi.\operatorname{loc}_v(t))=1\).
GalLink's character relation makes this
\(\overline\Psi(k_0.j_v(t))=1\).
Because \(\alpha\) factors through \(k\), the exact
inertia formula in (AFE.29) gives
\[
                    \psi_{v,i}(t)
                       =f\bigl(\overline\Psi(k_0.j_v(t))^{w_i}\bigr)
                       =1.
 \]
This is precisely the extra inertia condition in
the original FinPt, including the negative weight.
FinPt therefore gives, for every \(g\in\Phi.\Gamma_0\),
\[
       \operatorname{charpoly}(\rho(g))
         =\prod_{i=0}^{2}
                 \left(X-f(\Phi.\Psi(g)^{w_i})\right).
                                                        \tag{AFE.32}
\]

Take any \(s\in U_0\), retaining the original
\(\eta(s)=1\) condition of FinEnd.
Represent it as an element of \(G_L\) and apply
(AFE.32) to \(k_0.p_0(s)\).
By (AFE.29) the matrix on the left is \(R(s)\).
The character relation and \(f\operatorname{red}_O=\alpha\)
identify the right side with
\[
                  \prod_{i=0}^{2}(X-\alpha(\Psi(s)^{w_i})).
 \]
The matrix of the original residual reference
representation \(r_0(\eta,\Psi)(s)\) is
\(\operatorname{diag}(\Psi(s),1,\Psi(s)^{-1})\).
After applying \(\alpha\), its characteristic
polynomial is exactly this product.
This proves the original FinEnd equality.
Since \(U_0\) was chosen before \(A,r,s\), its
quantifier order is the required one.
The proof has used the supplied FinPt and has not
asserted its separate patching-theoretic producer.
\end{proof}
% END AFE20

% BEGIN AFE21
\paragraph{AFE21. The common marked frame now carries the original finite-end implication.}
\label{claim:AFE21}
Under \(P\), each original Dat input, and each original
HostData input with its specified \(\sigma,B\), admits
the same output \(L,g_1,\Phi,C_0\) as in AFE15 or
AFE16, together with its full original Marks and
the implication
\[
          \texttt{Fin\allowbreak Pt}(c.O,\Phi,3)
              \ \Longrightarrow\
          \texttt{Fin\allowbreak End}(\ell,\eta,\Psi,\sigma,g_1,B).
                                                        \tag{AFE.33}
\]
The \(\eta,\sigma,M,H_b\) in the two cases retain
their respective original meanings.

\begin{proof}
Choose the output only once, by AFE15 or AFE16.
AFE17 supplies Marks on its \(g_1\), including
\(\sigma^2=1\).
The actual modulus is divisible by \(\ell\) in
both cases.  The GalLink and all Ok fields are
those already chosen, and AFE1 supplies the
residue characteristic.
Apply AFE20 to this same data.
Its subgroup is exactly \(G_L\), independent of
the later coefficient algebra and ordinary point.
This proves (AFE.33) while preserving \(Q\) before
\(L,w\), the same \(g_1\), and the same frame.
No simultaneous choice involving an automorphic
representation has been added.
\end{proof}
% END AFE21

\paragraph{Boundary after these joins.}
The original cohomological premise \(P\) remains in
AFE15--16.  The resulting SeedClause supplies the
specified arithmetic characters; it does not supply
the separate rank-two and rank-three automorphic
seed representations.  Those original producers,
and the later \(\texttt{Fin\allowbreak Pt}\), \(\texttt{Aut\allowbreak Pt}\),
or their distinct Alt versions, remain explicit
inputs at their later applications.  The two Gen
entries, the full HasFrame conclusions at \(P\),
the actual auxiliary Host-frame witness statements,
the normalized-representation transport, and
\(\texttt{Fin\allowbreak Pt}\Rightarrow\texttt{Fin\allowbreak End}\)
are complete at their original premises.
The full common \(\texttt{Gd\allowbreak End}\) and
\(\texttt{Host\allowbreak End}\) still require their remaining
finite-point and automorphic suppliers.


\section{The actual standard and dual Bessel quotients}
\label{sec:bessel-quotient-suppliers-v52}

This chapter supplies the two original \texttt{Quot\allowbreak PStatement} and
\texttt{Quot\allowbreak MStatement}.  The spaces at the end are images of the actual
tensor coefficient maps, and the kernel being removed is the kernel of
all of the prescribed periods.  The intermediate statements retain the
broader source scopes: arbitrary generated-module tuples, arbitrary
cyclic lowest-weight modules, and intertwining maps into inner-product
spaces which need not be complete.

We use the reviewed DC, BSP, BST, BSE and BNV chapters.  Coordinates are
$0,\ldots,n-1$, $G_\epsilon=U(J_\epsilon)$ and
$K_\epsilon=G_\epsilon\cap U(n)$.  Write $p=p(\epsilon)$, $q=q(\epsilon)$,
$q_0=\min(p,q)$,
\[
 \rho_i=(n-1)/2-i,\qquad
 \Lambda(\lambda)=\lambda+\rho_n-\rho_c,\qquad
 c(\lambda)=\|\lambda\|^2-\|\rho\|^2.
 \tag{BQS.1}
\]
Positive roots are $\alpha_{ab}=e_a-e_b$ for $a<b$; compact roots have
$\epsilon_a=\epsilon_b$.  The \emph{integral positive-root cone} is
\[
 Q_+=\left\{\sum_{a<b}m_{ab}\alpha_{ab}:m_{ab}\in\mathbb Z_{\ge0}\right\}.
 \tag{BQS.2}
\]
For a weight $\gamma$ put $P_j(\gamma)=\sum_{i<j}\gamma_i$ and
$\operatorname{ht}(\gamma)=\sum_i(n-1-i)\gamma_i$.
Every $\gamma\in Q_+$ has integral coordinates, nonnegative integral
partial sums, total sum zero, and
$\operatorname{ht}(\gamma)=\sum_{a<b}(b-a)m_{ab}$.
Consequently $Q_+\cap(-Q_+)=\{0\}$, and a nonempty sum of positive
roots has strictly positive height.

An \emph{original unitary datum} below means every datum quantified in
\texttt{LowestKTypeEuler\allowbreak{}TensorStandard}: a complete complex Hilbert space
$H$, an actual continuous unitary representation $\sigma$ of $G_\epsilon$,
a derived invariant domain $D\subset H$ with matrix-unit operators
$E_{ab}$, and a finite-dimensional compact irreducible
$V\subset D$ of highest weight $\Lambda(\lambda)$, having Casimir
$c(\lambda)$ and stable under the actual compact group.  The parameter is
strictly decreasing and
$\lambda_i-(n-1)/2\in\mathbb Z$.  Also $p,q\ge1$ and
$\operatorname{GapLap}(\epsilon,\lambda,h)$ holds: the noncompact gaps
are at least $h$ and $h>h_{\rm need}(\epsilon)$, where exactly
\[
\begin{split}
h_{\rm need}(\epsilon)=\max\bigl\{&
2+\max(3n/2,1+npq),\
16nq_0+8\sqrt{2q_0}+2,\\
&8\sqrt{2q_0}\,[4\max(2nq_0,2(q-1)n^2)+1]\bigr\}.
\end{split}
\tag{BQS.3}
\]
In particular this datum is not assumed to come with a chosen coefficient
identity.  BSP2 supplies that identity at this exact gap when it is used.
Put $S=U(\mathfrak{gl}_n)V$.  The notation \emph{compact highest vector
of weight $\nu$} means $E_{aa}v=\nu_av$ and $E_{ab}v=0$ for
$a<b$ of equal sign; a zero vector is allowed in this predicate.

For any complex $\mathfrak{gl}_n$-module $M$, use the literal tensor
operators
\[
\begin{aligned}
 (A^+_{ab}f)_d&=E_{ab}f_d+\delta_{da}f_b,&
 (\Omega^+f)_d&=\sum_cE_{cd}f_c,\\
 (A^-_{ab}f)_d&=E_{ab}f_d-\delta_{db}f_a,&
 (\Omega^-f)_d&=-\sum_cE_{dc}f_c,\\
 \theta_s^+&=\lambda_s-(n-1)/2,&
 \theta_s^-&=-\lambda_s-(n-1)/2,\\
 X_s^\pm&=S^n\cap\ker(\Omega^\pm-\theta_s^\pm).
\end{aligned}
\tag{BQS.4}
\]
These are \texttt{tens\allowbreak V}, \texttt{tens\allowbreak Vd}, \texttt{ee\allowbreak V},
\texttt{ee\allowbreak Vd} and \texttt{Xs}; products of operators act from right to
left.

% BEGIN BQS1
\paragraph{BQS1. Tensor generation, commuting Euler operators and finite compact hulls.}
\label{claim:BQS1}
For every complex module $M$, every matrix-unit representation $E$ and
every subspace $V$, the module $S^n$, $S=U(\mathfrak{gl}_n)V$, is generated
by $V^n$ under either $A^+$ or $A^-$.  Both tensor actions satisfy the
matrix-unit relations and commute with the corresponding $\Omega^\pm$.
If $V$ is finite-dimensional and stable under the compact matrix units,
every $f\in X_s^\pm$ is contained in a finite-dimensional compact-stable
subspace of $X_s^\pm$.  No Hilbert completion is required.

\begin{proof}
Write $[v]_c$ for the tuple with $v$ in coordinate $c$ and zero elsewhere.
The identities
\[
 A^+_{ab}[v]_c=[E_{ab}v]_c+\delta_{bc}[v]_a,\qquad
 A^-_{ab}[v]_c=[E_{ab}v]_c-\delta_{ac}[v]_b
 \tag{BQS.5}
\]
show inductively that the tensor module generated by $V^n$ contains
$[E_{a_1b_1}\cdots E_{a_mb_m}v]_c$ for every word and $v\in V$.
The induction subtracts the second term of (BQS.5), whose word length
is smaller.  The reverse inclusion follows since both terms preserve
$S^n$.  This proves the full original \texttt{Gen\allowbreak Statement}.

The tensor relations follow either by substituting (BQS.5) or by adding
the commuting actions on the two factors $M\otimes\mathbb C^n$ and
$M\otimes(\mathbb C^n)^*$.  To check Euler commutation in the standard
case, the $d$-component of $[\Omega^+,A^+_{ab}]f$ is
\[
 \sum_c[E_{cd},E_{ab}]f_c+E_{ad}f_b
       -\delta_{da}\sum_cE_{cb}f_c=0.
 \tag{BQS.6}
\]
Indeed $[E_{cd},E_{ab}]=\delta_{da}E_{cb}-\delta_{bc}E_{ad}$.
For the dual case the same substitution gives
$-\sum_c[E_{dc},E_{ab}]f_c+E_{db}f_a
-\delta_{db}\sum_cE_{ac}f_c=0$.
Thus each Euler eigenspace is a full tensor submodule.

Let $F_m$ be the span of all words of length at most $m$ applied to $V$.
It is finite-dimensional.  A compact matrix unit preserves it: move that
unit past each word letter using the commutator relation, and then use
compact stability of $V$.  Every tuple in $S^n$ lies in some $F_m^n$.
The intersection $F_m^n\cap X_s^\pm$ is compact-stable, because
$F_m^n$ is stable under the compact tensor action and $\Omega^\pm$
commutes with that action.  It is the required finite hull.
\end{proof}
% END BQS1

% BEGIN BQS2
\paragraph{BQS2. The adjugate identities give the whole Euler polynomial.}
\label{claim:BQS2}
Suppose $S$ is an invariant subspace of a complex matrix-unit module and
the literal ordered Capelli operator has scalar value
\[
 \mathcal C_E(t)|_S=p(t)I,\qquad
 p(t)=\prod_i(t+\theta_i),\qquad t\in\mathbb C.
 \tag{BQS.7}
\]
Then on every tuple in $S^n$,
\[
 \prod_i(\Omega^+-\theta_i)=0,\qquad
 \prod_i\bigl(\Omega^-+\theta_i+n-1\bigr)=0.
 \tag{BQS.8}
\]
This assertion needs neither distinct roots nor an inner product.

\begin{proof}
For $n=0$, $S^0$ is the zero vector space, so its identity is zero.
Assume $n\ge1$.  BST8 gives, with exactly its ordered complementary
matrices $U(t),V(t)$,
\[
 (\Omega^++tI)U(t)=p(t)I,\qquad
 \bigl(-\Omega^-+(t-n+1)I\bigr)V(t)=p(t)I
 \quad\hbox{on }S^n.
 \tag{BQS.9}
\]
Here the matrices of operators act on column tuples.  Each minor
preserves $S$, so these are genuine identities of endomorphism-valued
polynomials on $S^n$.  Their scalar validity for every $t$ implies
coefficient validity: apply an arbitrary linear functional to their
value on an arbitrary vector and use the scalar polynomial identity.
The diagonal complementary minor has leading term $t^{n-1}$; every
off-diagonal complementary minor has degree at most $n-2$.  Thus both
$U$ and $V$ are monic of degree $n-1$ as operator polynomials.

For clarity, the elementary substitution argument in a noncommutative
endomorphism algebra is as follows.  If
$(tI-B)(t^{n-1}I+Q_1t^{n-2}+\cdots+Q_{n-1})=p(t)I$,
comparison from the highest coefficient downward gives
$Q_j=BQ_{j-1}+p_jI$, where
$p(t)=t^n+p_1t^{n-1}+\cdots+p_n$ and $Q_0=I$.
The constant coefficient then gives
$B^n+p_1B^{n-1}+\cdots+p_nI=0$.
This proves $p(B)=0$ without substituting into polynomials with
uncontrolled noncommuting coefficients.
Apply it first to $B=-\Omega^+$ and then to
$B=\Omega^-+(n-1)I$ in (BQS.9).
Removing the harmless scalar $(-1)^n$ in the first case gives (BQS.8).
\end{proof}
% END BQS2

% BEGIN BQS3
\paragraph{BQS3. Both original lowest-type Euler statements and their projectors.}
\label{claim:BQS3}
For every original unitary datum and every $f\in S^n$,
\[
 \prod_i(\Omega^+-\theta_i^+)f=0,\qquad
 \prod_i(\Omega^--\theta_i^-)f=0.
 \tag{BQS.10}
\]
These are the full original \texttt{LowestKTypeEuler\allowbreak{}TensorStandard}
and \texttt{LowestKTypeEuler\allowbreak{}TensorDual}.  For either sign the operators
\[
 \Pi_s^\pm=
 \prod_{t\ne s}\frac{\Omega^\pm-\theta_t^\pm}
                       {\theta_s^\pm-\theta_t^\pm}
 \quad\hbox{on }S^n
 \tag{BQS.11}
\]
are commuting idempotents with pairwise zero products, sum $I$, and
image $X_s^\pm$.  They commute with every corresponding tensor matrix
unit, not only the compact ones.

\begin{proof}
The gap (BQS.3) permits the full cyclic rigidity assertion of BSP2.
Choose its actual compact isometry and coefficient identity for this
datum.  BSP3, applied to that identity, gives
$p(t)=\prod_i(t+\lambda_i-(n-1)/2)$ on all of $S$.
BQS2 yields both assertions in (BQS.10), because
$-\theta_i^+-(n-1)=\theta_i^-$.
Thus no coefficient identity has been added to the hypotheses of
either original Euler statement.

The numbers $\theta_i^\pm$ are distinct by strict decrease of $\lambda$.
The scalar Lagrange polynomials $l_s$ satisfy
$l_s(\theta_t)=\delta_{st}$; the polynomials
$\sum_s l_s-1$, $l_s^2-l_s$, $l_sl_t$ for $s\ne t$, and
$(z-\theta_s)l_s(z)$ are divisible by
$\prod_i(z-\theta_i)$.  Evaluate these scalar polynomial identities
at $\Omega^\pm$ and use (BQS.10).
They give the stated projections and sum.  An eigenvector of eigenvalue
$\theta_s$ is fixed by $\Pi_s$, proving equality of its image with
$X_s^\pm$.  Commutation follows from BQS1.
\end{proof}
% END BQS3

% BEGIN BQS4
\paragraph{BQS4. The whole polynomial identity for a cyclic lowest-weight module.}
\label{claim:BQS4}
Let $M$ be any complex matrix-unit module generated by $\phi$, and suppose
\[
 E_{ab}\phi=0\quad(b<a),\qquad
 E_{aa}\phi=(\lambda_a+\rho_a)\phi.
 \tag{BQS.12}
\]
No regularity, integrality, positivity or strict decrease of $\lambda$
is assumed.  Both products in (BQS.10) vanish on every tuple of $M^n$.
This is the full original \texttt{Verma\allowbreak Product\allowbreak Identity}.

\begin{proof}
Set $v_t=[\phi]_t$.  Directly from (BQS.4), (BQS.5) and (BQS.12),
\[
\begin{aligned}
 (\Omega^+-\theta_t^+)v_t
    &=\sum_{d>t}A^+_{td}v_d,\\
 (\Omega^--\theta_t^-)v_t
    &=-\sum_{d<t}A^-_{dt}v_d .
\end{aligned}
\tag{BQS.13}
\]
For example the $t$-coordinate in the first identity is
$(\lambda_t+\rho_t-\theta_t^+)\phi=(n-1-t)\phi$;
the right side has exactly $n-1-t$ such contributions.  The other
coordinates are $E_{td}\phi$ for $d>t$, and zero for $d<t$.
In the second identity the $t$-coordinate is $t\phi$, and the same
coordinate check gives the displayed minus sign.

For the standard action let $F_j$ be the tensor submodule generated by
$v_t$, $t\ge j$.  For the dual action let $F_j$ be generated by
$v_{n-1-t}$, $t\ge j$.  In either case BQS1 gives $F_0=M^n$ and
$F_n=0$.  The $j$th quotient is cyclic, generated by the corresponding
$v_t$, and (BQS.13) says that $\Omega$ has scalar value $\theta_t$
on it.  BQS1 allows this equality to propagate to the whole quotient.
Hence $(\Omega-\theta_t)F_j\subset F_{j+1}$.
The product of the $n$ factors sends $F_0$ into $F_n=0$.
All factors are polynomials in one operator, so their order is
immaterial.  This proves both signs, including repeated roots.
\end{proof}
% END BQS4

% BEGIN BQS5
\paragraph{BQS5. The full cyclic-module cone and extremal-line theorem.}
\label{claim:BQS5}
Under (BQS.12), now assume $\lambda$ strictly decreasing.  Every
nonzero simultaneous diagonal weight vector $f$ of weight $\nu$ in
$\ker(\Omega^\pm-\theta_s^\pm)\subset M^n$ satisfies
\[
              \nu-(\lambda+\rho\pm e_s)\in Q_+.
 \tag{BQS.14}
\]
The simultaneous weight space of weight $\lambda+\rho\pm e_s$ in
that eigenspace has dimension at most one.  This is the whole original
\texttt{Verma\allowbreak Cone\allowbreak Statement}; it does not require $f$ to be highest
for any compact group.

\begin{proof}
Use the flag from BQS4, writing its ordered eigenvalues as
$\vartheta_0,\ldots,\vartheta_{n-1}$ and its $j$th generator as $v$.
Here $j=s$ in the standard case and $j=n-1-s$ in the dual case.
All $\vartheta_i$ are distinct.  If $\Omega f=\vartheta_j f$,
applying the first $j$ flag factors shows
$\prod_{i<j}(\vartheta_j-\vartheta_i)f\in F_j$.
Its nonzero scalar coefficient implies $f\in F_j$.

Modulo $F_{j+1}$ every lower matrix unit kills $v$:
in the standard case $A^+_{ab}v_s=\delta_{bs}v_a$ for $a>b$,
and $a>s$ puts this in the next flag term; in the dual case
$A^-_{ab}v_s=-\delta_{as}v_b$ and $b<s$ does the same.
The diagonal weight of $v$ is $b=\lambda+\rho\pm e_s$.
Commuting lower units to the right and diagonals past upper units
therefore shows that $F_j/F_{j+1}$ is spanned by positive-root
words applied to $v$.  This is only PBW spanning; no freeness or
irreducibility of the cyclic module is used.

Put $R=\prod_{i>j}(\Omega-\vartheta_i)$ and
$a=\prod_{i>j}(\vartheta_j-\vartheta_i)\ne0$.
The flag argument gives $R F_{j+1}=0$.  Thus every such eigenvector
$f$ satisfies
\[
           a f=R y
 \quad\hbox{for a finite sum $y$ of positive-root words on $v$}.
 \tag{BQS.15}
\]
The operator $R$ commutes with all tensor matrix units.  Each word of
length $m$ on $v$ has diagonal weight $b+\beta$, where
$\beta\in Q_+$ and $\operatorname{ht}(\beta)\ge m$.
Its image under $R$ has the same weight or is zero.

Distinct simultaneous weight spaces are linearly independent, even in
an arbitrary module: for any finite set of distinct weights choose a
linear combination of the diagonal operators separating them, and use
its scalar Lagrange polynomials.  Apply this observation to the finite
sum (BQS.15).  A nonzero weight vector can therefore have only a weight
$b+\beta$, proving (BQS.14).
For weight exactly $b$, every positive-length word has a different
weight by the strict positivity of height.  Its component in (BQS.15)
is consequently a scalar multiple of $Rv$.
The extremal weight space is contained in this one line, which may
itself be zero.  This proves the asserted dimension bound for both signs.
\end{proof}
% END BQS5

% BEGIN BQS6
\paragraph{BQS6. The actual principal-minor cyclic module.}
\label{claim:BQS6}
Let $\mu=\lambda+\rho\in\mathbb Z^n$, and let $\Phi_\mu$ be the
principal-minor function of DC2.  Write
\[
 \Omega_{\rm hol}=\{z:\Delta_j(z)\ne0\ (1\le j\le n)\},\quad
 \Omega_P=\{z:z^{-1}\in\Omega_{\rm hol}\},\quad
 \Phi_P(z)=\Phi_\mu(z^{-1}).
 \tag{BQS.16}
\]
On holomorphic functions on $\Omega_P$, extended by zero outside it,
define
$\mathcal E_{ab}F(z)=dF_z(zE_{ab})$ on $\Omega_P$ and zero outside.
The cyclic module $M_\mu=U(\mathfrak{gl}_n)\Phi_P$ is a literal module
of such functions.  Its generator satisfies (BQS.12), and it is
generated by that vector.  These are all fields of the original
\texttt{k2v\allowbreak M} model.

\begin{proof}
All leading minors are polynomial, and the final one is the determinant.
Thus $\Omega_{\rm hol}$ and $\Omega_P$ are open subsets of invertible
matrices.  Integer powers of their nonzero minors are holomorphic.
For a holomorphic $F$, $z\mapsto dF_z(zX)$ is holomorphic there;
extension by zero is only a function outside this open set and is not
claimed to be smooth across its boundary.
Commuting the two vector fields $z\mapsto zX,z\mapsto zY$ gives
the field $z\mapsto z[X,Y]$, so the $\mathcal E_{ab}$ obey the
matrix-unit relations, also outside $\Omega_P$ where they all vanish.

For $a>b$, $(z\exp(tE_{ab}))^{-1}=\exp(-tE_{ab})z^{-1}$.
Triangular multiplication in DC2 leaves every leading minor unchanged
under this lower unipotent multiplication.  Its derivative is zero.
For $E_{aa}$ the same identity and DC2's diagonal character give
$\Phi_P(z\exp(tE_{aa}))=e^{t\mu_a}\Phi_P(z)$ locally.
Thus $\mathcal E_{aa}\Phi_P=\mu_a\Phi_P$ and lower units kill the
generator.  Restricting the operators to their generated submodule
preserves all identities and makes generation tautological.
This is exactly the source's zero-extended \texttt{MMT}, with no
unproved global analytic continuation attached to it.
\end{proof}
% END BQS6

\subsection{The compact integral transfers the cone to the actual datum}

% BEGIN BQS7
\paragraph{BQS7. Full transfer of finite derivative relations through the compact average.}
\label{claim:BQS7}
Let $\lambda$ be Harish--Chandra and let $T$ be any original
compact type of weight $\Lambda(\lambda)$,
$w_T$ its unit highest vector, and $\mu=\lambda+\rho$.
For arbitrary finitely many lists $\mathbf Z_i$ of complex matrices
and arbitrary $v_i\in T$, suppose
\[
 \sum_i\langle T(k)w_T,v_i\rangle
       R^{\rm hol}_{\operatorname{Ad}(k^{-1})\mathbf Z_i}
                         \Phi_P(z)=0
 \quad(k\in K_\epsilon,\ z\in\Omega_P).
 \tag{BQS.17}
\]
Then, for every $g\in G_\epsilon$,
\[
                  \sum_i R_{\mathbf Z_i}\Psi_{T,\mu}(g)v_i=0.
 \tag{BQS.18}
\]
Here every list keeps its order, $R^{\rm hol}_ZF(z)=dF_z(zZ)$, and
$R_Z$ is the original complexified tangential right derivative.
Empty lists and zero summands are allowed.  This is the full original
\texttt{Transfer\allowbreak Statement}.

\begin{proof}
Use probability Haar on $K$; replacing it by the source's Haar changes
the coefficient $\dim T$ into $\dim T/\mu_K(K)$ and changes no identity.
DC2 gives a common open neighborhood $\mathcal U$ of the positive cone
on which every $k^{-1}zk$ is in $\Omega_{\rm hol}$ and
\[
 P(z)=(\dim T)\int_K\Phi_\mu(k^{-1}zk)
                 T(k)P_{w_T}T(k)^{-1}\,dk.
 \tag{BQS.19}
\]
All derivatives of this identity are obtained under the integral.
Indeed for a fixed $z\in\mathcal U$ choose a smaller relatively compact
complex neighborhood whose closure remains in $\mathcal U$.
The conjugated minors have a uniform positive lower bound in modulus
on its product with compact $K$.  Every fixed derivative is a continuous
rational expression with these nonzero denominators, so it is uniformly
bounded there.  Differentiation under this finite measure integral follows.
This argument works for every finite word, without one word-length bound
being imposed on the statement.

Fix $g$.  DC3 supplies $k_0\in K$ and a neighborhood on which
$\Psi(g')=T(k_0)P((g')^{-1}k_0)$.  Put $z_0=g^{-1}k_0$.
If $L_ZF(z)=dF_z(Zz)$, the inverse rule gives
\[
 R_{\mathbf Z}\Psi(g)
       =T(k_0)L_{-\mathbf Z}P(z_0),
 \qquad
 R^{\rm hol}_{\mathbf Z}\Phi_P(y)
       =L_{-\mathbf Z}\Phi_\mu(y^{-1}).
 \tag{BQS.20}
\]
Induction proves these formulas with the displayed order: at each
step $(g'\exp(tZ))^{-1}=\exp(-tZ)(g')^{-1}$.
Complex linearity gives the same result for every complex matrix $Z$.

Differentiate (BQS.19), apply it to each $v_i$, and sum.
The rank-one factor satisfies
\[
 T(k)P_{w_T}T(k)^{-1}v_i
             =\langle T(k)w_T,v_i\rangle T(k)w_T.
 \tag{BQS.21}
\]
At this $k$ use (BQS.20) with
$y=(k^{-1}z_0k)^{-1}\in\Omega_P$.
The scalar multiplying $T(k)w_T$ is exactly the sum in (BQS.17).
It is zero at every $k$.  The integral and hence its left multiplication
by $T(k_0)$ are zero.  This proves (BQS.18).
\end{proof}
% END BQS7

% BEGIN BQS8
\paragraph{BQS8. Formal word space, compact action and highest-vector lifting.}
\label{claim:BQS8}
Let $D$ be any complex inner-product module, possibly incomplete, on
which $E$ satisfies the matrix-unit relations and
$E_{ab}^*=E_{ba}$ for compact pairs.  Let $V$ be finite-dimensional
and compact-stable.  Form the algebraic direct sum
\[
 \mathcal W=\bigoplus_{m\ge0}
             \mathfrak{gl}_n^{\otimes m}\otimes V\otimes\mathbb C^n .
 \tag{BQS.22}
\]
Its elementary vectors are denoted $[\mathbf Z,v,c]$; expansion in
matrix units identifies it with the original finitely supported word
space \texttt{KS}.  Define
\[
\begin{aligned}
 \mathcal T[\mathbf Z,v,c]&=[E_{\mathbf Z}v]_c,\\
 \mathcal B[\mathbf Z,v,c]&=\sum_d[(E_{cd},\mathbf Z),v,d],\\
 \mathcal A_X[\mathbf Z,v,c]&=
 [\operatorname{ad}(X)\mathbf Z,v,c]+[\mathbf Z,E_Xv,c]
                 +[\mathbf Z,v,Xe_c]\quad(X\in\mathfrak k_{\mathbb C}).
\end{aligned}
\tag{BQS.23}
\]
The adjoint term means the sum over the list's individual positions.
Then $\mathcal T$ maps onto $S^n$,
$\mathcal T\mathcal B=\Omega^+\mathcal T$,
$\mathcal T\mathcal A_{ab}=A^+_{ab}\mathcal T$ for compact pairs, and
$[\mathcal A_{ab},\mathcal B]=0$.
Every compact highest vector $f\in S^n$ of weight $\nu$ has a compact
highest lift $u\in\mathcal W$ of that same weight, with $\mathcal T u=f$.
This includes the full original \texttt{lift\_\allowbreak highest\allowbreak Weight\allowbreak Vector}.

\begin{proof}
Surjectivity follows by expanding each component of $f$ as a finite sum
of words on $V$.  Prepending the indicated matrix unit proves
$\mathcal T\mathcal B=\Omega^+\mathcal T$.
For compact $X$, the Leibniz commutator identity
$[E_X,E_{\mathbf Z}]=E_{\operatorname{ad}(X)\mathbf Z}$
proves the second intertwining identity.
The $\mathcal A$ are the sum of the adjoint, $V$, and standard tensor
actions, hence satisfy all compact matrix-unit relations.
The element $\sum_{c,d}E_{cd}\otimes E_{dc}$ is invariant under
simultaneous conjugation.  Expanding its commutator with a matrix unit
cancels the two Kronecker sums; this is exactly
$[\mathcal A_{ab},\mathcal B]=0$ on (BQS.23).

Give each matrix-letter space its Hilbert--Schmidt inner product, give
$V$ its inherited inner product, and take the tensor/direct-sum inner
products in (BQS.22).  On compact pairs,
$(\operatorname{ad}E_{ab})^*=\operatorname{ad}E_{ba}$,
as follows from $\operatorname{tr}(X^*Y)$ and cyclicity of trace.
The same adjoint identity holds on the other two factors, and therefore
$\mathcal A_{ab}^*=\mathcal A_{ba}$.
Words of length at most $m$ form a finite-dimensional compact-stable
space $\mathcal W_{\le m}$.

Choose any preimage of $f$ in some $\mathcal W_{\le m}$.
The kernel of $\mathcal T$ on that space is compact-stable.  Its
orthogonal complement is also compact-stable by the adjoint identities.
The restriction of $\mathcal T$ to this complement is an isomorphism
onto $\mathcal T(\mathcal W_{\le m})$.
Lift $f$ by its inverse.  The intertwining identities and injectivity
show that the lift has the same diagonal eigenvalues and is killed
by every compact upper unit.  All orthogonal complements used here
are finite-dimensional; completeness of $D$ was never invoked.
\end{proof}
% END BQS8

% BEGIN BQS9
\paragraph{BQS9. The full trace-kernel implication for the actual coefficient.}
\label{claim:BQS9}
Let $\lambda$ be Harish--Chandra, and let $(H,\sigma,D,E,V)$
have the complete unitary, derived and compact irreducibility
properties of an original datum.  Suppose an original
$T,\mu$ and linear isomorphism $\iota:V\to T$ satisfy the exact coefficient
identity
\[
 \langle v',\sigma(g)v\rangle
       =\langle\iota v',\Psi_{T,\mu}(g)\iota v\rangle
       \quad(g\in G_\epsilon,\ v,v'\in V).
 \tag{BQS.24}
\]
For this assertion no gap or Casimir hypothesis is additionally required.
Put $v_0=\iota^{-1}w_T$, $\ell(v)=\langle v_0,v\rangle$, and
\[
 \Theta[\mathbf Z,v,c]
       =[\ell(v)\mathcal E_{\mathbf Z}\Phi_P]_c
                   \quad\hbox{in }M_\mu^n.
 \tag{BQS.25}
\]
For \emph{every} compact-stable subspace $W\subset\mathcal W$,
\[
                    \Theta W=0\quad\Longrightarrow\quad\mathcal T W=0.
 \tag{BQS.26}
\]
This proves the full original \texttt{lowest\allowbreak KType\_\allowbreak trace\_\allowbreak integral},
and hence \texttt{lowest\allowbreak KType\_\allowbreak trace}.

\begin{proof}
At $g=1$, (BQS.24) and $\Psi(1)=1$ show that $\iota$ is an isometry.
The finite-dimensional compact-stable space $V$ is stable under the
actual compact group.  To see this directly, solve the orbit differential
equation for each compact real Lie direction on $V$; its matrix
exponential agrees with the actual $\sigma$-orbit by uniqueness.
Every unitary block matrix is the exponential of a skew-Hermitian block
matrix, so this covers $K$.
Putting $g=k$ in (BQS.24), using $\Psi(k)=T(k)$ and all $v'\in V$,
now gives $\iota\sigma(k)|_V=T(k)\iota$.

The actual compact action on the word space is
\[
 \mathcal U(k)=
  \bigoplus_m\operatorname{Ad}(k)^{\otimes m}
                    \otimes\sigma(k)|_V\otimes k.
 \tag{BQS.27}
\]
Its derivative is $\mathcal A$.  It preserves word length.
For each $x\in W$, choose $m$ with $x\in\mathcal W_{\le m}$.
The intersection $W\cap\mathcal W_{\le m}$ is finite-dimensional and
$\mathcal A$-stable.  Matrix exponentials and the preceding logarithms
therefore show $\mathcal U(k)x\in W$ for every $k$.
Consequently $\Theta(\mathcal U(k^{-1})x)=0$.

Write $x=\sum_i[\mathbf Z_i,v_i,c_i]$.
Evaluation of the last zero function at $z\in\Omega_P$ gives, in
coordinate $d$,
\[
 \sum_i(k^{-1})_{dc_i}\,
   \langle T(k)w_T,\iota v_i\rangle
   R^{\rm hol}_{\operatorname{Ad}(k^{-1})\mathbf Z_i}
                                 \Phi_P(z)=0.
 \tag{BQS.28}
\]
Here unitarity moved $\sigma(k^{-1})$ from the second slot to the
first slot of $\ell$; the isometry $\iota$ then gave the displayed
coefficient.  Multiplication of these coordinate equations by $k$
isolates each fixed index $c$ and gives the hypothesis of BQS7 for
the summands with $c_i=c$.
It follows that
$\sum_{c_i=c}R_{\mathbf Z_i}\Psi(g)\iota v_i=0$
for all $g$.

Differentiate (BQS.24) through these finite words.  The derived
domain is invariant, so induction gives
\[
 \left\langle v',\sigma(g)
       \sum_{c_i=c}E_{\mathbf Z_i}v_i\right\rangle=0
             \quad(g\in G_\epsilon,\ v'\in V).
 \tag{BQS.29}
\]
Let $C$ be the closed span of all $\sigma(g)V$.  It contains every
word on $V$: real derived actions are limits of difference quotients
in this closed invariant subspace, and complex matrix units are
complex linear combinations of real directions.
The vector in (BQS.29) lies in $C$ and is orthogonal to every
$\sigma(g^{-1})V$, hence to $C$.  It is zero.
This holds for every coordinate $c$, proving $\mathcal T x=0$.
The argument applies to each $x\in W$ without a finite-dimensional
assumption on $W$ itself.
\end{proof}
% END BQS9

% BEGIN BQS10
\paragraph{BQS10. The universal cone-and-line transfer.}
\label{claim:BQS10}
Here is the full abstract statement used by
\texttt{lowest\allowbreak KType\_\allowbreak tensor\_\allowbreak lift}.  Let $\mathcal S$ be any complex
vector space with compact matrix-unit operators $\mathcal A$ and an
operator $\mathcal B$ commuting with them.  Let
$T:\mathcal S\to M^n$ and $\Theta:\mathcal S\to M_0^n$ be linear,
where $M$ is any complex vector space with an arbitrary operator
family $E$, and $M_0$ is a cyclic matrix-unit module satisfying (BQS.12), with strictly
decreasing $\lambda$.  Let $P\subset M$ be a subspace.  Assume
\[
\begin{aligned}
 T\mathcal B&=\Omega^+_E T,&
 \Theta\mathcal B&=\Omega^+_{\mathcal E}\Theta,\\
 \Theta\mathcal A_{aa}
  &=A^+_{\mathcal E,aa}\Theta+
           \bigl(\Lambda(\lambda)_a-\lambda_a-\rho_a\bigr)\Theta.
\end{aligned}
\tag{BQS.30}
\]
Assume also that every compact highest vector in $P^n$ has a
same-weight highest lift through $T$, and that for every compact-stable
$W\subset\mathcal S$, $\Theta W=0$ implies $TW=0$.
Then every nonzero compact highest vector $f$ of weight $\nu$ in
$P^n\cap\ker(\Omega^+_E-\theta_s^+)$ satisfies
\[
                      \nu-\Lambda(\lambda+e_s)\in Q_+,
 \tag{BQS.31}
\]
and its highest-vector space of weight $\Lambda(\lambda+e_s)$
has dimension at most one.

\begin{proof}
Put $q_s(z)=\prod_{t\ne s}(z-\theta_t^+)$ and
$a_s=q_s(\theta_s^+)\ne0$.
If $Tu=f$ is a highest lift, $u'=q_s(\mathcal B)u$ is still highest,
because $\mathcal B$ commutes with the compact action, and
$Tu'=a_sf$.
BQS4 and the second intertwining in (BQS.30) show that
$\Theta u'$ lies in the selected Euler eigenspace in $M_0^n$.
The same is true of every compact lowering word on $u'$.
If that word lowers the weight by $\delta$, a sum of positive compact
roots, the diagonal relation in (BQS.30) gives its $\Theta$-image
weight
\[
                  \nu-\delta+\lambda+\rho-\Lambda(\lambda).
 \tag{BQS.32}
\]
This identity follows by induction on the word using the compact
matrix-unit commutator; it does not assert that $\Theta$ intertwines
the off-diagonal compact operators.

If $\nu-\Lambda(\lambda+e_s)\notin Q_+$, BQS5 forces every such
image to vanish.  Otherwise a nonzero image would imply
$\nu-\Lambda(\lambda+e_s)-\delta\in Q_+$, and adding $\delta$
would contradict the assumption.  Compact PBW expresses the compact
module generated by the highest vector $u'$ as the span of these
lowering words: commute upper units to the right, where they kill
$u'$, and absorb diagonal units into weights.
This span is compact-stable.  The assumed trace-kernel implication
therefore gives $Tu'=0$, contrary to $a_sf\ne0$.
This proves (BQS.31).

Now take two highest vectors of weight $\Lambda(\lambda+e_s)$,
lift them to $u_1,u_2$, and apply $q_s(\mathcal B)$.
Their two $\Theta$-images have weight $\lambda+\rho+e_s$ in the
selected cyclic-model eigenspace.  BQS5 gives scalars $c_1,c_2$,
not both zero, whose combination has zero image.
Every nonempty compact lowering word on the corresponding highest
combination has model weight $\lambda+\rho+e_s-\delta$ with
$\delta\ne0$ a positive-root sum.  BQS5 and positivity of height
force these images to vanish too.
Apply the trace-kernel implication to its compact lowering span.
The resulting equality
$a_s(c_1f_1+c_2f_2)=0$ gives the required linear dependence.
No finite-dimensionality of $\mathcal S$ or ambient bracket assumptions
beyond the stated compact relations have entered this argument.
\end{proof}
% END BQS10

% BEGIN BQS11
\paragraph{BQS11. The full original standard eigenspace cone and line.}
\label{claim:BQS11}
For every original unitary datum and every $s$, every nonzero compact
highest vector $f\in X_s^+$ of weight $\nu$ has
$\nu-\Lambda(\lambda+e_s)\in Q_+$.
Any two compact highest vectors of weight $\Lambda(\lambda+e_s)$
in $X_s^+$ are linearly dependent.  These are both clauses of the
original \texttt{Eigenspace\allowbreak Weight\allowbreak Cone\allowbreak Standard}.

\begin{proof}
Use BSP2 to identify this datum's lowest type isometrically with $T$
and obtain (BQS.24), retaining the given $\lambda,h,\sigma,D,E,V$.
The vector $v_0=\iota^{-1}w_T$ is a compact highest vector of weight
$\Lambda(\lambda)$: differentiate compact intertwining, first in real
directions and then complexify.
Thus $\ell(E_{aa}v)=\Lambda(\lambda)_a\ell(v)$, using the
compact adjoint identity and reality of the weight.

Use BQS6 for the cyclic model, BQS8 for the formal word space, and
BQS9 for its trace-kernel implication.
The two Euler intertwinings in (BQS.30) follow by prepending a letter.
On $[\mathbf Z,v,c]$, the diagonal commutator in the first term of
$\mathcal A_{aa}$ gives
$\mathcal E_{aa}\mathcal E_{\mathbf Z}\Phi_P
 -(\lambda_a+\rho_a)\mathcal E_{\mathbf Z}\Phi_P$.
The $V$ term gives $\Lambda(\lambda)_a\ell(v)$, and the coordinate
term gives the standard tensor summand.
Their sum is exactly the last equality of (BQS.30).
Highest lifts are BQS8.  Therefore BQS10 applies with $P=S$ and
gives both original conclusions.
\end{proof}
% END BQS11

% BEGIN BQS12
\paragraph{BQS12. Actual reversed duality and the whole dual cone statement.}
\label{claim:BQS12}
Both conclusions of BQS11 hold for $X_s^-$ with
$\Lambda(\lambda-e_s)$ in place of $\Lambda(\lambda+e_s)$.
This proves the full original \texttt{Eigenspace\allowbreak Weight\allowbreak Cone\allowbreak Dual}.
The transport below also preserves arbitrary infinitesimally unitary
target modules, without completing them.

\begin{proof}
Let $\bar a=n-1-a$ and let $j_a\in\{1,-1\}$ be the $a$th diagonal
entry of $J_\epsilon$.  On the same space $D$ set
\[
 \epsilon'_a=\epsilon_{\bar a},\quad
 \lambda'_a=-\lambda_{\bar a},\quad
 E'_{ab}=-j_{\bar a}j_{\bar b}E_{\bar b,\bar a},\quad
 (Uf)_a=j_{\bar a}f_{\bar a}.
 \tag{BQS.33}
\]
This complex-linear signed reversal $U$ is invertible.
The actual new group representation is
$\sigma'(g)=\sigma(R\overline gR)$, where $R e_a=e_{\bar a}$.
The map $g\mapsto R\overline gR$ preserves products, sends
$G_{\epsilon'}$ and $K_{\epsilon'}$ to their original counterparts,
and preserves real one-parameter subgroups.
For a real Lie matrix $X$ the relation
$X^*J_{\epsilon'}+J_{\epsilon'}X=0$ gives
$\overline{X_{ba}}=-j_{\bar a}j_{\bar b}X_{ab}$.
It follows that the complexified derived operators of $\sigma'$ are
exactly $E'$ in (BQS.33).

The matrix-unit relations, Casimir and compact adjoint identities are
preserved.  For example
$(E'_{ab})^*=j_{\bar a}j_{\bar b}E'_{ba}$, the required adjoint sign
for $\epsilon'$, and summing $E'_{ab}E'_{ba}$ cancels both signs and
reindexes the old Casimir.
A highest vector of weight $\nu$ becomes highest of weight
$-\nu\circ\bar{\ }$: if $a<b$, then $\bar b<\bar a$, so
$E'_{ab}$ is a scalar multiple of an original upper unit.
Compact irreducibility and compact stability are unchanged.
The half-root sums satisfy
$\rho_c(\epsilon')=-\rho_c(\epsilon)\circ\bar{\ }$ and likewise for
$\rho_n$, by reversing each root.
Thus
$\Lambda_{\epsilon'}(\lambda')=-\Lambda_\epsilon(\lambda)\circ\bar{\ }$,
$c(\lambda')=c(\lambda)$, and the same original gap and $h_{\rm need}$
hold.  Integrality is preserved because
$-\lambda_{\bar a}-(n-1)/2$ differs from
$-(\lambda_{\bar a}-(n-1)/2)$ by the integer $-(n-1)$.

Direct substitution in (BQS.4) gives
\[
 \Omega^+_{E'}U=U\Omega^-_E,\qquad
 A^+_{E',ab}U
     =U\bigl(-j_{\bar a}j_{\bar b}A^-_{E,\bar b,\bar a}\bigr).
 \tag{BQS.34}
\]
For example in the Euler identity the two coefficients
$j_{\bar c}$ from $U$ cancel $j_{\bar c}^2=1$, leaving
$-j_{\bar d}\sum_cE_{\bar d,c}f_c$.
Generated modules also agree, since (BQS.33) is an invertible
relabeling and rescaling of all matrix units.
Therefore $U$ sends $X_s^-$ to the standard eigenspace for
$(\epsilon',\lambda',\bar s)$, and preserves highest-vector relations
with the just described reversed weights.

Apply BQS11 there.  Its shifted lowest weight is
$-\Lambda_\epsilon(\lambda-e_s)\circ\bar{\ }$.
The map $\beta\mapsto-\beta\circ\bar{\ }$ preserves $Q_+$, because
it sends $\alpha_{ab}$ to $\alpha_{\bar b,\bar a}$.
This proves the cone conclusion; invertibility and linearity of $U$
give the extremal-line conclusion.
The same formula (BQS.33) for operators on any inner-product target
preserves its matrix-unit and adjoint identities by the same checks.
No use of target completeness occurs.
\end{proof}
% END BQS12

\subsection{The lowest type in every nonzero unitary image}

% BEGIN BQS13
\paragraph{BQS13. One highest vector controls every target intertwiner.}
\label{claim:BQS13}
More generally, let $Y$ be any complex inner-product space with a
family of compact operators $A_{ab}$ satisfying the compact adjoint
identities.  Let $X\subset Y$ have finite compact-stable hulls inside
$X$, and assume its highest-vector space of a prescribed real weight
$\nu$ has dimension at most one.
There exists one $y_s\in X$, chosen before any target, such that for
every complex vector space $N$, every family of operators $B_{ab}$
on $N$ and every linear map $\varphi:Y\to N$ intertwining the compact
operators on $X$, every highest vector $w$ of weight $\nu$ in
$\varphi(X)$ is a scalar multiple of $\varphi(y_s)$.
Neither matrix-unit relations on $A$ nor a norm, completeness,
bracket or irreducibility hypothesis on the target are needed.
In particular this holds for either sign with $Y=D^n$,
$X=X_s^\pm$ and $\nu=\Lambda(\lambda\pm e_s)$ for every original
unitary datum.  It supplies the full abstract \texttt{line} and
its original \texttt{line\_\allowbreak V} application.

\begin{proof}
If $w=\varphi(x)$, enclose $x$ in a finite compact-stable
subspace $F\subset X$ as assumed.
The kernel of $\varphi|_F$ is compact-stable, and its orthogonal
complement is compact-stable.  Inverting $\varphi$ on that complement
lifts $w$ to a highest vector of the same weight in $X$.
This proof uses only the intertwining equations on $F$.

If the source highest-vector space is zero, take $y_s=0$; the lifting
argument says that every indicated $w$ is zero.
Otherwise choose any nonzero $y_s$ in it.  The assumed dimension
bound says that this source highest-vector space is one-dimensional.  Every lift is
therefore a scalar multiple of $y_s$, proving the assertion uniformly
for all targets and all maps.
For the stated applications, both tensor actions obey the compact
adjoint identities for the sum inner product on $D^n$.
BQS1 supplies the finite hulls and BQS11 or BQS12 supplies the
one-dimensional source highest space.  Thus the conclusion has
the required original quantifier order, with the same vector for
every target and map.
\end{proof}
% END BQS13

% BEGIN BQS14
\paragraph{BQS14. The full Schmid--Casimir comparison without target completeness.}
\label{claim:BQS14}
Let $N$ be any complex inner-product space and $B$ a matrix-unit
representation with
$B_{ab}^*=j_aj_bB_{ba}$ in the sense of the inner-product identities.
Let $W\subset N$ be a compact irreducible type of highest weight $\nu$.
Put
$\eta=\nu+\rho_c-\rho_n$, and suppose $\eta$ is regular integral,
its noncompact gaps are at least $b$, $p,q\ge1$, and
$h_{\rm need}(\epsilon)<b+2$.
On $\mathfrak p_{\mathbb C}\otimes W$ let
\[
 A_W=\sum_{\epsilon_a=\epsilon_b}
           \operatorname{ad}(E_{ab})\otimes B_{ba}|_W,\qquad
 m(X\otimes w)=B_Xw.
 \tag{BQS.35}
\]
Let $W_-$ be the sum of the negative eigenspaces of this self-adjoint
operator, exactly the original \texttt{Wm}.
If $m(W_-)=0$ and the Casimir has scalar value $c\in\mathbb C$
on $W$, then
\[
                               c=c(\eta).
 \tag{BQS.36}
\]
This is the whole original \texttt{Schmid\allowbreak Casimir\allowbreak Statement}.

\begin{proof}
The first two bounds in (BQS.3) give
\[
 b>npq+1,\qquad b>3n/2,\qquad
 b>16nq_0+8\sqrt{2q_0}.
 \tag{BQS.37}
\]
DC1--DC11 therefore construct the actual unitary model of parameter
$\eta$.  Let $W_0$ be its lowest compact type.  It has highest weight
$\nu$, Casimir $c(\eta)$, and the same negative first-layer relation
by DC8.  Choose an isometric compact intertwiner $i:W\to W_0$ and
put $\Delta=\{(w,iw):w\in W\}$ inside the inner-product direct sum
of $N$ and the model's smooth module.
This direct sum need not be complete.

The compact coupling operator (BQS.35) is natural for compact
intertwiners: $1\otimes i$ intertwines every summand.  Hence the
negative spectral subspace for $\Delta$ is the diagonal copy of the
negative spectral subspace for $W$, also identified with that for
$W_0$.  Multiplication kills it in each coordinate, and therefore
kills it in the direct sum.

Apply the purely algebraic finite-layer argument of DC9 to $\Delta$.
For precision, it requires only this negative first-layer relation
and $b>npq+1$: moving a negative tensor slot to the last position
changes a word by smaller-length terms; a minimal-norm lift of a new
compact type in degree $m$ consequently has every slice in the positive
coupling space.  The compact Casimir identity then gives
\[
 c_K(\tau)\ge c_K(\nu)+2m(b-n+1)\quad(m\ge1).
 \tag{BQS.38}
\]
All lifts and complements in this argument are finite-dimensional.
Since $p,q\ge1$ implies $pq\ge n-1$, the factor $b-n+1$ is positive.
Thus the original compact type $\nu$ occurs in the generated module
only in degree zero and has its one highest line.

Take a nonzero highest vector $(w,iw)$ of $\Delta$.
The total Casimir commutes with the compact action, so its image is
another highest vector of weight $\nu$ in the generated module.
By the preceding uniqueness it is a scalar multiple of $(w,iw)$.
But its two coordinates are $cw$ and $c(\eta)iw$.
As both coordinates are nonzero, the two scalars agree.
This proves (BQS.36) using the same finite-layer comparison as the
source's diagonal-model argument; it assumes neither a completion nor
an actual noncompact group action on $N$.
\end{proof}
% END BQS14

% BEGIN BQS15
\paragraph{BQS15. The exact two-unit gap margin for a candidate compact type.}
\label{claim:BQS15}
Let $\lambda,h$ satisfy the original datum hypotheses, fix $s,t$, and
suppose $\nu$ is compact dominant with
\[
 \nu-\Lambda(\lambda+e_s)\in Q_+,\qquad
 P_j(\nu)\le P_j(\Lambda(\lambda)+e_t)\quad(0\le j\le n).
 \tag{BQS.39}
\]
Then $\eta=\nu+\rho_c-\rho_n$ is regular integral, has noncompact
gaps at least $h-2$, and
\[
 h-2>npq+1,\qquad h_{\rm need}< (h-2)+2,\qquad
 \eta=(\lambda+e_s)+
                 \bigl(\nu-\Lambda(\lambda+e_s)\bigr).
 \tag{BQS.40}
\]

\begin{proof}
Put $\gamma=\nu-\Lambda(\lambda)$ and
$\beta=\nu-\Lambda(\lambda+e_s)\in Q_+$.
Then $\gamma=e_s+\beta$.  The root-cone lower bounds and (BQS.39)
give
\[
        0\le P_j(\gamma)\le1,\qquad
        -1\le\gamma_i=P_{i+1}(\gamma)-P_i(\gamma)\le1.
 \tag{BQS.41}
\]
In particular $|\gamma_a-\gamma_b|\le2$, and every $\gamma_i$ is an
integer.  As $\eta=\lambda+\gamma$, integrality and the noncompact
gap at least $h-2$ follow.

To prove strict decrease it suffices to check adjacent indices $a,a+1$.
For opposite signs,
$\eta_a-\eta_{a+1}\ge h-2>0$.
For equal signs, the half-root counts give
$(\rho_c)_a-(\rho_c)_{a+1}=1$ and
$(\rho_n)_a-(\rho_n)_{a+1}=0$.
Compact dominance gives
$\eta_a-\eta_{a+1}=\nu_a-\nu_{a+1}+1\ge1$.
Finally $h_{\rm need}\ge npq+3$, by (BQS.3), so $h-2>npq+1$.
The other identities in (BQS.40) are the original strict inequality
$h>h_{\rm need}$ and the definition of $\Lambda$.
The margin has not been enlarged.
\end{proof}
% END BQS15

% BEGIN BQS16
\paragraph{BQS16. Every nonzero standard tensor image contains the shifted lowest type.}
\label{claim:BQS16}
Let an original unitary datum and an index $s$ be given, and assume only
that $\lambda+e_s$ is strictly decreasing.
Let $N$ be any complex inner-product space, possibly incomplete, carrying
matrix-unit operators $B$ with the infinitesimal unitary adjoint signs
for $\epsilon$.
For every linear map $\varphi:D^n\to N$ satisfying
\[
       \varphi(A^+_{ab}f)=B_{ab}\varphi(f)
                  \quad(a,b,\ f\in S^n)
 \tag{BQS.42}
\]
and every $x\in X_s^+$ with $\varphi(x)\ne0$, there exists a subspace
$W\subset\varphi(X_s^+)$ which is compact irreducible of highest weight
$\Lambda(\lambda+e_s)$ and on which the Casimir is $c(\lambda+e_s)$.
This is the full original \texttt{Quotient\allowbreak Highest\allowbreak Weight\allowbreak Tensor\allowbreak Standard}.
It assumes neither an actual group representation on $N$ nor a chosen
coefficient identity for the source datum.

\begin{proof}
Let $Y=\varphi(X_s^+)$ and let $\Pi=\Pi_s^+$ from BQS3.
BQS1 and (BQS.42) show that $Y$ is Lie-stable and compact-finite.
If $\varphi\Pi$ vanished on $V^n$, its intertwining and the tensor
generation in BQS1 would force it to vanish on all of $S^n$.
This contradicts $\Pi x=x$ and $\varphi(x)\ne0$.
Thus the finite compact-stable image $\varphi\Pi(V^n)$ is nonzero.

Finite-dimensional compact-unitary modules are orthogonal sums of
compact irreducibles.  The elementary argument suffices here:
the commuting self-adjoint diagonal operators have simultaneous real
weight spaces; choose a maximal weight for the compact positive order
to obtain a highest vector; its generated irreducible summand can be
split off orthogonally and the procedure repeated.
Choose such a highest vector in $\varphi\Pi(V^n)$, of weight $\nu_0$.
Lift it through $\varphi\Pi|_{V^n}$ using the finite-dimensional
kernel-complement argument of BQS13.
If coordinate $t$ of its nonzero highest lift is nonzero, that coordinate
is a weight vector in $V$ of weight $\nu_0-e_t$.
Every weight of the compact highest-weight module $V$ differs from
$\Lambda(\lambda)$ by the negative of a sum of positive compact roots:
compact PBW moves all upper units to its highest vector.
Consequently
\[
              P_j(\nu_0)\le P_j(\Lambda(\lambda)+e_t)
                         \quad\hbox{for every }j .
 \tag{BQS.43}
\]

Every nonzero highest vector in $Y$ lifts to one in $X_s^+$ by the
same finite-hull argument.  BQS11 therefore gives
$\nu-\Lambda(\lambda+e_s)\in Q_+$ for every compact highest weight
occurring in $Y$.
Among those satisfying the upper bound (BQS.43), choose one $\nu$
minimizing the nonnegative integer
$\operatorname{ht}(\nu-\Lambda(\lambda+e_s))$.
This collection is nonempty by $\nu_0$.
Let $W\subset Y$ be a compact irreducible subspace of highest weight
$\nu$, contained in a finite compact hull.
Its weight is compact dominant; for a compact simple root this follows
also from
$\|B_{ba}w\|^2=(\nu_a-\nu_b)\|w\|^2\ge0$ on a highest vector.

BQS15 applies, so
$\eta=\nu+\rho_c-\rho_n$ is regular integral and has the exact gap
$h-2$ needed in BQS14.
We claim that the negative first noncompact layer of $W$ is killed
by multiplication.
The finite tensor calculation of DC8 applies to $W$ with parameter
$\eta$.  Its negative coupling space is a sum of compact types
\[
                   \nu-\alpha_{ab},
             \qquad a<b,\quad\epsilon_a\ne\epsilon_b.
 \tag{BQS.44}
\]
To recall why the sign is exact, the compact standard/dual tensor rule
in the two sign blocks gives types $\nu+e_r-e_t$.
Their coupling eigenvalues are
$\eta_r+(\rho_n)_r-\eta_t-(\rho_n)_t+1-n/2$.
The gap $h-2>npq+1$ separates their signs: a negative eigenvalue
requires $r>t$, giving (BQS.44).

If multiplication had a nonzero image on this negative space, its
finite-dimensional compact-unitary image would contain a highest
vector of one of the weights (BQS.44).
It belongs to $Y$, since $Y$ is Lie-stable.
Its partial sums are no larger than those of $\nu$, hence still
satisfy (BQS.43).  Its root-cone difference from
$\Lambda(\lambda+e_s)$ remains in $Q_+$ by the preceding highest-lift
argument, but its height has decreased by $b-a>0$.
This contradicts minimality.  Thus the required Schmid relation holds.

The original Casimir is scalar on all of $S$ by centrality and generation.
BST10's exact tensor identity gives on $X_s^+$
\[
 C_{A^+}=c(\lambda)+n+2\theta_s^+
         =c(\lambda)+2\lambda_s+1=c(\lambda+e_s).
 \tag{BQS.45}
\]
Intertwining transfers this scalar to $Y$, and in particular to $W$.
BQS14 now identifies it with $c(\eta)$.
Writing $\beta=\nu-\Lambda(\lambda+e_s)\in Q_+$, BQS15 gives
$\eta=\lambda+e_s+\beta$, and hence
\[
 0=\|\lambda+e_s+\beta\|^2-\|\lambda+e_s\|^2
   =2\langle\lambda+e_s,\beta\rangle+\|\beta\|^2.
 \tag{BQS.46}
\]
The first term is nonnegative: it is a sum of nonnegative integer
multiples of $(\lambda+e_s)_a-(\lambda+e_s)_b$, $a<b$.
The second is nonnegative and vanishes only when $\beta=0$.
Thus $\nu=\Lambda(\lambda+e_s)$, proving both required properties
of $W$.  Every use of orthogonal decomposition was in a finite-dimensional
subspace of $N$, so target completeness was not assumed.
\end{proof}
% END BQS16

% BEGIN BQS17
\paragraph{BQS17. The full original dual unitary-image theorem.}
\label{claim:BQS17}
Under the same original datum, suppose $\lambda-e_s$ is strictly
decreasing.  For every infinitesimally unitary matrix-unit module
$(N,B)$, every $\varphi:D^n\to N$ intertwining $A^-$ on all of $S^n$,
and every $x\in X_s^-$ with $\varphi(x)\ne0$, there exists
$W\subset\varphi(X_s^-)$ compact irreducible of highest weight
$\Lambda(\lambda-e_s)$ and Casimir $c(\lambda-e_s)$.
This is the entire original \texttt{Quotient\allowbreak Highest\allowbreak Weight\allowbreak Tensor\allowbreak Dual},
at the same possibly incomplete target scope.

\begin{proof}
Apply the actual reversal BQS12 to the source and set
$B'_{ab}=-j_{\bar a}j_{\bar b}B_{\bar b,\bar a}$ on the same target
$N$.  Its matrix-unit and infinitesimal unitary identities were checked
there.  The map
$\varphi'=\varphi U^{-1}$ intertwines $A^+_{E'}$ with $B'$ on
the whole transformed generated module, by (BQS.34).
It is nonzero on $Ux$.
The relevant standard shift is
\[
            \lambda'+e_{\bar s}
                     =-(\lambda-e_s)\circ\bar{\ },
 \tag{BQS.47}
\]
which is strictly decreasing precisely when $\lambda-e_s$ is.
BQS16 gives the required compact irreducible space $W$ in the same
image $\varphi(X_s^-)$, but expressed using $B'$.
Reversing again returns its highest weight to
$\Lambda_\epsilon(\lambda-e_s)$.
The Casimir is unchanged by the signed transpose reversal and
$c(-(\lambda-e_s)\circ\bar{\ })=c(\lambda-e_s)$.
This proves the original dual conclusion with all source and target
quantifiers preserved.
\end{proof}
% END BQS17

% BEGIN BQS18
\paragraph{BQS18. The original closure and compact-finite assertions.}
\label{claim:BQS18}
Let $\sigma$ be an actual unitary representation of $G_\epsilon$ on
an inner-product space $H$, with derived invariant domain $(D,E)$.
If $W\subset D$ is finite dimensional, stable under the compact
matrix units, and has a real scalar Casimir, then
\[
 C=\overline{U(\mathfrak g_{\mathbb C})W}^{\,H}
                 \quad\hbox{is invariant under every }\sigma(g),
                 \quad g\in G_\epsilon .
 \tag{BQS.48}
\]
Completeness of $H$ is unnecessary for this assertion.
If $H$ is complete and, additionally, $W$ is compact irreducible
of weight $\Lambda(\eta)$, with Casimir $c(\eta)$, and
\[
 \eta\text{ is strictly decreasing},\qquad
 \exists b_0>npq+1:\quad
 |\eta_a-\eta_b|\ge b_0\quad(\epsilon_a\ne\epsilon_b),
 \tag{BQS.49}
\]
then every
finite-dimensional actual $K_\epsilon$-stable subspace of $C$
lies in $U(\mathfrak g_{\mathbb C})W$.
This second assertion is the full original \texttt{data\_\allowbreak kfin\_\allowbreak le}
under \texttt{Std\allowbreak Gap}; it assumes no integrality or coefficient identity.

\begin{proof}
Write $F_j$ for the finite span of words of at most $j$ noncompact
matrix units on $W$.  Commuting a compact unit to the right shows
that every $F_j$ is compact stable and that their union is the
generated module.  A finite compact-stable subspace is stable under
the actual compact group.  Indeed, for a compact real Lie algebra
element $X$, the finite-dimensional curve $e^{tE(X)}w$ and
$\sigma(e^{tX})w$ solve the same differential equation with the
same initial value; differentiating
$\sigma(e^{-tX})e^{tE(X)}w$ makes it constant.
Every unitary matrix in each sign block has a skew-Hermitian logarithm,
so this proves stability for all of $K_\epsilon$.
The same argument applies to $F_j$.

Let $c$ be the scalar Casimir.  Centrality extends it from $W$
to every $F_j$.  The finite compact tensor decomposition of
$\mathfrak p_{\mathbb C}^{\otimes j}\otimes W$ bounds its compact
highest weights by $A+\sqrt2j$ for a fixed $A$.
The energy identity
$L=2C_K-C_G$, obtained by summing the squares of a real orthonormal
compact and noncompact basis with the infinitesimal unitary adjoint
signs, therefore bounds the norm of $E(X)$ on $F_j$ by
$C_X(A'+j)$ for fixed $A'\ge1$ and $C_X$.
Iterating, for $v\in F_k$,
\[
 \|E(X)^m v\|
 \le C_X^m\prod_{r=0}^{m-1}(A'+k+r)\,\|v\|.
 \tag{BQS.50}
\]
For $|t|C_X<1$ the Taylor series converges.
The integral Taylor remainder for
$t\mapsto\sigma(e^{tX})v$ has norm at most
$|t|^m\|E(X)^m v\|/m!$; unitarity gives this bound at every
base point.  Its successive majorants have ratio tending to
$|t|C_X<1$.  The remainder tends to zero, so the series has
the actual orbit as its sum, also if $H$ is incomplete.
All its partial sums lie in the generated module, and thus
$\sigma(e^{tX})v\in C$ for these $t$.
Continuity of the bounded operator $\sigma(e^{tX})$ extends
this inclusion from the dense generated module to $C$.
Subdividing $t=1$ proves stability under $e^X$.
Polar decomposition $g=k e^X$, $X\in\mathfrak p$, together with
the compact stability already proved gives (BQS.48).
Applying it to $g^{-1}$ gives equality.

For the second assertion write $b_0$ for the real bound in
(BQS.49).  If one sign block is empty, all matrix units are compact,
the generated module is $W$, and $C=W$ is already finite dimensional.
Otherwise $b_0>npq+1$ implies $b_0-n+1>0$.
The finite Casimir argument in DC8 supplies the negative first-layer
Schmid relation at exactly this bound.  DC9 then gives, on the
orthogonal new part of $F_j$, $j\ge1$,
\[
 c_K(\nu)\ \ge\ c_K(\Lambda(\eta))
                         +2j(b_0-n+1)
 \tag{BQS.51}
\]
for each nonzero compact highest vector of weight $\nu$.
Neither argument uses integrality of $\eta$: the compact module
$W$ is already supplied and its finite-dimensional unitary
decomposition is the only compact input.
Consequently every vector of a fixed compact type $\nu$ in the
generated module lies in some one finite layer $F_{J(\nu)}$.
To see this for arbitrary word length, decompose a vector
orthogonally through the nested finite compact-stable layers.
Each nonzero new component has the same type and its layer index
is bounded by (BQS.51).

The actual compact character projector onto $\nu$ is bounded on
$H$: it is the integral over probability Haar of the continuous
unitary operators multiplied by the finite-dimensional character
and its dimension.  On every finite compact hull, finite
Schur orthogonality identifies its range with the $\nu$-isotypic
subspace.  Its image on the dense generated module is therefore
contained in $F_{J(\nu)}$, a closed finite-dimensional subspace.
Continuity gives the same inclusion on $C$.
Any finite-dimensional actual compact-stable subspace of $C$ is
the direct sum of finitely many such compact types, so it lies
in the generated module.  This proves the stated full weak-gap
assertion without inserting the stronger constructive gap.
\end{proof}
% END BQS18

% BEGIN BQS19
\paragraph{BQS19. The full universal-image closure lemma.}
\label{claim:BQS19}
Let $Y$ be any complex vector space, $A_{ab}\in\operatorname{End}Y$
any family of linear operators, and $G,X\subset Y$ any subspaces.
No matrix-unit relations on $A$, stability of $G$ or $X$, or inclusion
$X\subset G$ is assumed here.
Fix a real parameter $\eta$ and write
$\mathcal D=\operatorname{smD}(r_{\rm reg})$ for the smooth domain
in the actual right regular $L^2(G_\epsilon)$ representation, with
derived operators $B$.
Call $\psi:Y\to\mathcal D$ admissible when
\[
               \psi(A_{ab}y)=B_{ab}\psi(y)
                            \quad(y\in G,\ a,b).
 \tag{BQS.52}
\]
Assume the following two universal properties.
For every admissible $\psi$, every $x\in X$ with $\psi(x)\ne0$
gives a compact irreducible $W_\psi\subset\psi(X)$ of weight
$\Lambda(\eta)$ and Casimir $c(\eta)$.
There is one specified $y_s\in Y$ such that, for every admissible
$\psi$, every highest vector of weight $\Lambda(\eta)$ in $\psi(X)$
is a scalar multiple of $\psi(y_s)$.
Given an admissible $\varphi$ and a compact irreducible
$W\subset\mathcal D$ of this weight and Casimir with
$\varphi(y_s)\in W$, one has
\[
       \varphi(X)\subset
       \overline{U(\mathfrak g_{\mathbb C})W}^{\,L^2}.
 \tag{BQS.53}
\]
This is the original \texttt{image\_\allowbreak le\_\allowbreak clos\allowbreak Gen}, without a gap
assumption and without an assumption that $W\subset\varphi(X)$.

\begin{proof}
Let $C$ be the closure on the right.  BQS18 makes it invariant
under the actual right regular representation, and unitarity and
invariance under inverses make $C^\perp$ invariant as well.
Its bounded orthogonal projection $P$ commutes with that representation.
If $v$ is smooth, its orbit is locally the restriction of a smooth
ambient map; composing that map with $P$ proves that $Pv$ is smooth.
Differentiating $P r(g)v=r(g)Pv$ along each real one-parameter
subgroup and complexifying gives $PB_{ab}v=B_{ab}Pv$.
Thus $\psi=P\varphi$, viewed as a map to $\mathcal D$, is admissible.
As $\varphi(y_s)\in W\subset C$, one has $\psi(y_s)=0$.

If $\psi(x)\ne0$ for some $x\in X$, the universal image property
gives a nonzero compact irreducible $W_\psi\subset\psi(X)$.
Take its nonzero highest vector $w$.
The universal line property says $w=a\psi(y_s)=0$, a contradiction.
Therefore $P\varphi(X)=0$.  Since $C$ is closed, $\ker P=C$,
which is exactly (BQS.53).
Completeness has been used only for the actual $L^2$ orthogonal
projection; the source $Y$ remains arbitrary.
\end{proof}
% END BQS19

% BEGIN BQS20
\paragraph{BQS20. The entire original generic quotient theorem.}
\label{claim:BQS20}
Use the arbitrary source $Y,A,G,X$, parameter $\eta$, target
$\mathcal D$ and admissibility condition (BQS.52) of BQS19.
Assume now $X\subset G$ and \texttt{Std\allowbreak Gap} (BQS.49) for $\eta$.
Retain the universal nonzero-image property of BQS19, and assume
there exists \emph{one} $y_s\in X$ with its universal highest-line
property for every admissible map.
Suppose every $x\in X$ lies in a finite-dimensional compact
matrix-unit-stable subspace $S_x\subset X$.
Let $Z\subset X$ be stable under every $A_{ab}$, and let $\varphi$
be admissible and satisfy the exact saturation condition
\[
 y\in X,\quad z\in Z,\quad \varphi(y)=\varphi(z)
                       \quad\Longrightarrow\quad y\in Z.
 \tag{BQS.54}
\]
For every $x\in X\setminus Z$ there exist subspaces $Y'\subset X$
and $W\subset\mathcal D$ such that
\[
 \varphi(Y')=W,\qquad Y'\cap Z=0,\qquad
 W\text{ is compact irreducible of weight }\Lambda(\eta),
 \qquad C_B|_W=c(\eta),
 \tag{BQS.55}
\]
and $W$ is stable under every element of the actual $K_\epsilon$.
No relations on $A$, complete norm on $Y$, global intertwining
outside $G$, or full Lie stability of $X$ has been added.

\begin{proof}
If $\varphi(x)=0=\varphi(0)$, (BQS.54) would give $x\in Z$.
Thus the universal image property gives a compact irreducible
$W\subset\varphi(X)$ of the required weight and Casimir.
Its nonzero highest vector has the form $a\varphi(y_s)$ by the
universal line property.  Here $a\ne0$, so $\varphi(y_s)\in W$.
BQS19 gives $\varphi(X)\subset C=\overline{U(\mathfrak g)W}$.

The image $\varphi(S_x)$ is finite dimensional and compact
matrix-unit-stable: $S_x\subset X\subset G$ permits (BQS.52)
on every vector of $S_x$.  The finite-dimensional group
argument in BQS18 makes it stable under the actual compact
group.  The weak-gap conclusion of BQS18 therefore implies
\[
                       \varphi(x)\in U(\mathfrak g)W.
 \tag{BQS.56}
\]
Put $N_0=\varphi(Z)$.  For $z\in Z$, the inclusion
$z\in X\subset G$ and (BQS.52) give
$B_{ab}\varphi(z)=\varphi(A_{ab}z)\in N_0$.
Hence $W\cap N_0$ is compact stable, so compact irreducibility
makes it either zero or $W$.
The second alternative would imply
$U(\mathfrak g)W\subset N_0$.
By (BQS.56), $\varphi(x)=\varphi(z)$ for some $z\in Z$,
contradicting (BQS.54).  Thus $W\cap N_0=0$.

Let $Y_0=X\cap\varphi^{-1}(W)$.
Since $W\subset\varphi(X)$, the restriction
$\varphi:Y_0\to W$ is surjective.
Choose an algebraic complement $Y'$ to its kernel in $Y_0$.
Then $\varphi|_{Y'}$ is an isomorphism onto $W$.
If $y\in Y'\cap Z$, its image lies in $W\cap N_0=0$,
and injectivity gives $y=0$.
This proves all of (BQS.55).  Finally $W$ is finite dimensional
and stable under the compact derived operators in the actual
smooth right regular representation; the logarithm and orbit
argument in BQS18 proves its actual $K_\epsilon$ stability.
The complement is an algebraic choice and imposes no unwanted
continuity requirement on the original source.
\end{proof}
% END BQS20

\subsection{The original smooth maps and the actual period kernel}

For the last three claims let $n=2r+k+1$, $p=k+r$, $q=r+1$.
The original step hypotheses are
\[
 \max(3n/2,1+npq)<h-2,\qquad
 4\max(2nq_0,2rn^2)<h/(8\sqrt{2q_0})-1.
 \tag{BQS.57}
\]
These are exactly \texttt{Step\allowbreak Gap}; BSP1 proves $q_0\ge1$
from these inequalities, including the original zero-division convention.
Let $\lambda$ be Harish--Chandra with noncompact gaps at least $h$.
We retain every original unitary datum with this parameter and, in the
last two claims, the supplied $T,\mu,\iota$ with
$\mu=\lambda+\rho$ and the actual coefficient identity (BSP.4).

The Bessel data are arbitrary original data:
$u\in\mathbb R^k$, a middle vector $z$ with $T_z=-z^*H_{pk}z>0$,
$c\in(\mathbb C^\times)^r$, a complete complex Hilbert space $V_\pi$,
a representation $\pi$ satisfying the full
$\texttt{Is\allowbreak HCType\allowbreak W}(k,r,z,u,\pi)$, and ambient matrix measures
$\mu_W,\mu_N$ satisfying the original weak
$\texttt{Is\allowbreak Haar\allowbreak On}(H_z,\mu_W)$ and
$\texttt{Is\allowbreak Haar\allowbreak On}(N,\mu_N)$.
Here $H_z=\texttt{UWset}(k,r,z)$, $N=\texttt{Nset}(k,r)$,
and $\psi_{z,c}$ is the actual character in (BNV.1).
Let $P$ be any invertible Cayley matrix
$P^*J_\epsilon P=H_{pk}$.  Euclidean unitarity is not required.
Put $x_1=x_J(P,z)=P\,x_{\rm One}(z)$, with the prescribed Bessel
vector $x_{\rm One}$, and fix any
$x_0\in G_\epsilon$, $w\in V$, and $\xi\in V_\pi$.
For a tuple $f\in D^n$, the exact two maps on ambient matrices are
\[
 \begin{split}
 L^+(f)(g)&=\sum_c
  (x_1^*J_\epsilon ge_c)\,
       \langle w,\sigma(x_0^{-1}g)f_c\rangle,\\
 L^-(f)(g)&=\sum_c
  (e_c^t g^{-1}x_1)\,
       \langle w,\sigma(x_0^{-1}g)f_c\rangle .
 \end{split}
 \tag{BQS.58}
\]
Off the group, the inverse has the same total matrix-inverse convention
as the source.  These are its \texttt{t\allowbreak Map} and \texttt{t\allowbreak Map\allowbreak D},
not equivalence classes of functions.
For either scalar function $F$ the actual vector-valued period is
\[
 \mathcal P_F(g)=
 \int F(P h' n' gP^{-1})\,
       \overline{\psi_{z,c}(n')}\,\pi(h'^{-1})\xi\,
             d(\mu_W\otimes\mu_N)(h',n'),\qquad g\in G_H=U(H_{pk}).
 \tag{BQS.59}
\]
The integral is on the two ambient matrix spaces; the weak Haar
predicates give support on $H_z$ and $N$ almost everywhere.
Every choice in this paragraph, including $\xi$, remains fixed
through the quotient construction.

% BEGIN BQS21
\paragraph{BQS21. The smooth lift and the exact simultaneous period kernel.}
\label{claim:BQS21}
First, in arbitrary rank and signature, let $Y$ be any complex vector
space, let $A_{ab}$ be any endomorphisms, let $G\subset Y$ be stable
under them, and let $L:Y\to(\operatorname{Mat}_n(\mathbb C)\to\mathbb C)$
be linear.
Suppose for every $f\in G$ that $Lf$ is smooth on $G_\epsilon$,
all its derivative words are in $L^2$, and
\[
                   E_{ab}(Lf)=L(A_{ab}f)\quad\hbox{on }G_\epsilon.
 \tag{BQS.60}
\]
Then there is a linear map $\Phi:Y\to\mathcal D$ such that on $G$
\[
        \Phi(f)=\operatorname{toL2}(Lf),\qquad
        \Phi(A_{ab}f)=B_{ab}\Phi(f).
 \tag{BQS.61}
\]
This is the entire original \texttt{exists\_\allowbreak lin}.

For the original data (BQS.57)--(BQS.59), apply this to
$Y=D^n$, $G=S^n$, $A=A^\pm$, $L=L^\pm$.
For each $s$, the subspace
\[
 Z_s^\pm=\{y\in X_s^\pm:
                     \mathcal P_{L^\pm(y)}(g)=0
                           \text{ for every }g\in G_H\}
 \tag{BQS.62}
\]
is stable under every $A^\pm_{ab}$ and satisfies the exact
saturation condition
\[
 y\in X_s^\pm,\quad z'\in Z_s^\pm,\quad \Phi(y)=\Phi(z')
                      \quad\Longrightarrow\quad y\in Z_s^\pm .
 \tag{BQS.63}
\]
If the supplied tuple $x\in X_s^\pm$ has a nonzero period
at some $g\in G_H$, then $x\notin Z_s^\pm$.

\begin{proof}
On $G$, every $\operatorname{toL2}(Lf)$ is an actual $L^2$ class,
so the operation is linear there; the source's value zero outside
the $L^2$ domain never enters.  These classes are smooth vectors.
Here are the analytic details of that last assertion.
For a real Lie algebra element $X$, all derivatives $R_X^mLf$
are finite complex linear combinations of the given word derivatives
and belong to $L^2$.
Right translations are strongly continuous on $L^2$:
it holds first for continuous compactly supported functions by
uniform continuity on a compact neighborhood and then for all
$L^2$ functions by density and unitarity.
The pointwise fundamental theorem of calculus for $Lf$ therefore
becomes the Bochner identity
\[
 r(e^{tX})[Lf]-[Lf]
                =\int_0^t r(e^{sX})[R_XLf]\,ds
                   \quad\hbox{in }L^2 .
 \tag{BQS.64}
\]
The integrand is strongly continuous and has constant $L^2$ norm;
on finite intervals Fubini and Cauchy--Schwarz identify its
Bochner integral with the pointwise integral.
Dividing by $t$ gives the derived vector $[R_XLf]$.
Repeating for every mixed derivative, and using coordinates which
are finite products of exponentials of a real Lie basis, makes
the orbit map $C^\infty$.  A local smooth retraction from ambient
matrices to the Lie group supplies the exact local ambient smooth
extensions used by \texttt{sm\allowbreak D}, as in BSP5.
Thus the map $f\mapsto[Lf]$ on $G$ lands in $\mathcal D$ and
(BQS.60) gives its intertwining.
Choose an algebraic projection $q:Y\to G$ with $q|_G=1$ and set
$\Phi(f)=[L(qf)]$.  This proves (BQS.61) on all the originally
required vectors; no continuity or relations on $A$ were assumed.

For the actual maps, BQS1 gives invariance of $S^n$ and of
$X_s^\pm$, and BST5 gives (BQS.60), smoothness, and all word
estimates with the unchanged exponent
$1-h/(8\sqrt{2q_0})$.
Thus the generic construction applies.
BSP5 and BNV1 supply absolute Bochner integrability of every
period in (BQS.59), for every such word and every $g\in G_H$.
Consequently period addition and scalar multiplication are genuine
integral identities.  They make (BQS.62) a linear subspace.

The exact differentiation identity of BSP5 identifies the period
of $E_{ab}(L^\pm y)$ with the corresponding complexified right
derivative, in the Cayley-transported direction, of
$g\mapsto\mathcal P_{L^\pm y}(g)$.
If the latter vanishes everywhere on $G_H$, each real tangential
derivative vanishes, and then so does its complexification.
Together with (BQS.60) and invariance of $X_s^\pm$, this proves
$A^\pm_{ab}Z_s^\pm\subset Z_s^\pm$.
There is no integration by parts against a different character or
a different measure in this step.

Finally, $\Phi(y)=\Phi(z')$ says that the two continuous functions
$L^\pm y$ and $L^\pm z'$ have the same $L^2$ class on
$G_\epsilon$.  Haar measure has full support, so a nonzero value
of their continuous difference would give a nonempty open set
of positive measure on which its absolute value is bounded below.
Their difference is therefore zero at every point of $G_\epsilon$.
For almost every $(h',n')$ in the given ambient product measure,
$h'\in H_z$, $n'\in N$, and hence
$P h'n'gP^{-1}\in G_\epsilon$ for every $g\in G_H$.
The two actual integrands in (BQS.59) agree almost everywhere.
Thus their periods agree for every $g$, proving (BQS.63).
The final assertion is precisely the negation of the defining
condition in (BQS.62).
\end{proof}
% END BQS21

% BEGIN BQS22
\paragraph{BQS22. Full original \texttt{Quot\allowbreak PStatement} and \texttt{Quot\allowbreak MStatement}.}
\label{claim:BQS22}
Take every original datum in (BQS.57)--(BQS.59), every index $s$,
and an arbitrary tuple $x\in X_s^\pm$.
Assume $\lambda'=\lambda\pm e_s$ is Harish--Chandra, its
noncompact gaps are at least the \emph{same} $h$, and
$\mathcal P_{L^\pm(x)}(g)\ne0$ for some $g\in G_H$.
There exist a subspace $W_f$ of complex-valued functions on
ambient matrices and a subspace $W'\subset\mathcal D$ with all
of the following original properties:
\[
 \begin{array}{ll}
 {\rm (a)}&
  \text{Every }F\in W_f\text{ equals }L^\pm(y)
                        \text{ for some }y\in X_s^\pm;\\
 {\rm (b)}&
  v\in W'\ \Longleftrightarrow\
       \exists F\in W_f:\ v=\operatorname{toL2}(F)
                         \quad(v\in\mathcal D);\\
 {\rm (c)}&
  W'\text{ is compact irreducible of highest weight }
                               \Lambda(\lambda');\\
 {\rm (d)}& C_B v=c(\lambda')v\quad(v\in W');\\
 {\rm (e)}&
  r_{\rm reg}(k')W'=W'\quad\text{for every actual }k'\in K_\epsilon;\\
 {\rm (f)}&
  F\in W_f,\quad \mathcal P_F(g)=0\text{ for every }g\in G_H
               \quad\Longrightarrow\quad F|_{G_\epsilon}=0.
 \end{array}
 \tag{BQS.65}
\]
The sign $+$ gives the full original \texttt{Quot\allowbreak PStatement};
the sign $-$ gives the full original \texttt{Quot\allowbreak MStatement}.
The tuples are arbitrary elements of the original Euler eigenspaces,
not only the tuples produced by a complementary minor.

\begin{proof}
BSP1 turns (BQS.57) into $p,q\ge1$ and
$h>h_{\rm need}(\epsilon)$, so the supplied unitary datum is in
the exact scope of BQS1--BQS17.
Also $h>npq+1$; the hypotheses on $\lambda'$ therefore give
the weak \texttt{Std\allowbreak Gap} (BQS.49) used by the generic quotient
theorem.  The rank-zero impossibility is supplied by BSP1,
not by an extra restriction on the quantified $k,r$.

Use BQS20 with
\[
 Y=D^n,\quad G=S^n,\quad X=X_s^\pm,\quad
 A=A^\pm,\quad Z=Z_s^\pm,\quad \eta=\lambda',
 \tag{BQS.66}
\]
and with the smooth lift $\Phi$ supplied by BQS21.
Here $X\subset G$, and BQS1 supplies every finite compact hull.
BQS21 supplies stability and saturation of $Z$, admissibility
of $\Phi$, and $x\notin Z$.
For \emph{every} admissible $\psi:Y\to\mathcal D$, the target
$\mathcal D$ has its actual infinitesimally unitary matrix-unit
action.  It need not be complete in its inherited norm.
BQS16, respectively BQS17, therefore supplies the universal
nonzero-image property required by BQS20 at exactly this target.
BQS13 supplies one $y_s\in X$ before the choice of $\psi$ and
the required universal highest-line property.
Thus every hypothesis of BQS20 holds with the same $G,X,Z,\Phi$.

Let $Y'\subset X$ and $W$ be its conclusion, and define
\[
                         W_f=L^\pm(Y'),\qquad W'=W.
 \tag{BQS.67}
\]
Property (a) is immediate.
If $v\in W'=\Phi(Y')$, choose $y\in Y'$ with $v=\Phi(y)$;
BQS21 gives $v=\operatorname{toL2}(L^\pm y)$.
Conversely, if $F=L^\pm y$ with $y\in Y'$ and
$v=\operatorname{toL2}(F)$, the same identity gives $v=\Phi(y)\in W'$.
This proves both directions of the exact image statement (b).
Properties (c), (d), and (e) are exactly the three representation
conclusions of BQS20.

For (f), write $F=L^\pm y$ with $y\in Y'$.
If every prescribed period vanishes, (BQS.62) says $y\in Z$.
Since $Y'\cap Z=0$, one has $y=0$ and hence $F=0$ as a function
on all ambient matrices, which in particular gives the required
pointwise vanishing on $G_\epsilon$.
All measures, character, Cayley matrix, $\xi$, $w$, and $x_0$
in this proof were the original supplied ones.
The same proof with the already established dual image and
line theorems proves every field for the minus sign.
\end{proof}
% END BQS22

% BEGIN BQS23
\paragraph{BQS23. Full one-weight propagation and the original endomorphism nonvanishing.}
\label{claim:BQS23}
The complete original \texttt{Step\allowbreak Statement} now holds.
Consequently the complete original \texttt{Bessel\allowbreak NVGap} holds:
for every even $k\ge2$, every $r\ge0$, all original
$z,c,u,v,b$ satisfying the Harish--Chandra, interlacing,
half-integrality and \texttt{Gap\allowbreak Big} conditions of (BNV.1)--(BNV.3),
every representation $\sigma$ satisfying
$\texttt{Is\allowbreak DSHpk}(k,r,\epsilon^{48},\lambda,\sigma)$,
every $\pi$ satisfying $\texttt{Is\allowbreak HCType\allowbreak W}(k,r,z,u,\pi)$,
every original lowest-type unit vector $w$, and every original
weak Haar pair, there exist $x_0,x\in G_H$ such that
\[
 \int
 \langle\sigma(x_0)w,\sigma(h'n')\sigma(x)w\rangle\,
    \overline{\psi_{z,c}(n')}\,\pi(h')^{-1}\,
                  d(\mu_W\otimes\mu_N)(h',n')\ne0
       \quad\text{in }\operatorname{End}_{\mathbb C}(V_\pi).
 \tag{BQS.68}
\]
The original \texttt{Concrete\allowbreak Setting\allowbreak Exists} also follows on the
same supplied \texttt{Forms0} scope as BWP24.

\begin{proof}
BQS22 discharges both exact quotient premises retained in BSE9.
For clarity, the one-weight argument uses the already proved
EvalP and EvalM to turn a nonzero source coefficient period
into a nonzero period of the respective complementary-minor
tuple.  BST's complementary Euler identity places that tuple
in $X_s^+$ or $X_s^-$.
BQS22 produces its actual function quotient with all six fields
(BQS.65), while BST5 supplies the inherited smoothness and
unchanged decay exponent for every function in it.
The full \texttt{Fin\allowbreak Statement} proved in BSP, with the full
\texttt{Int\allowbreak SStatement} of BNV1, therefore produces the shifted
original $\texttt{Bessel\allowbreak NZ}$.
The exact evaluation factors in BSE7--BSE8 give precisely the
original exclusions $u_j\ne\lambda_s\pm\tfrac12$.
Both signs are included, so this proves the whole
\texttt{Step\allowbreak Statement}, as assembled in BST11 and BSE9.

BNV23 now has its last premise supplied.
Explicitly, BNV7 chooses an admissible real gap before the finite
walk; BNV18 gives a nonzero period at a sufficiently far deformed
tail parameter; BNV6 walks back using the two step directions
just proved.  BNV21 transports the resulting existential
nonvanishing to every original datum, and BNV22 returns it to
the literal endomorphism-valued integral (BQS.68).
That return includes the independently existential Cayley
witnesses, arbitrary lowest unit vector, arbitrary original
nondegenerate character and Haar normalization, and the
actual order $h'n'$ of the two group factors.
Thus (BQS.68) is the prescribed operator integral, not merely
the existence of an abstract branching functional.

Finally BWP24 consumes exactly this original
\texttt{Bessel\allowbreak NVGap}, constructs and normalizes the original
wave packets on the same adelic data, and joins the reviewed
ASG and AMD setting and measure producers.
It gives \texttt{Concrete\allowbreak Setting\allowbreak Exists} for each supplied
\texttt{Forms0} datum.  That premise is retained exactly:
this invocation does not construct the ambient form record.
\end{proof}
% END BQS23

\paragraph{Scope of this completion.}
This chapter completes the original Euler polynomials, K2 cone and
extremal line, K3 image types, both actual smooth quotients,
both one-weight propagations, and the original Bessel nonvanishing
theorem.  Its supplied-form setting consequence is conditional only
on the inherited, explicit \texttt{Forms0} input.
The ambient form supplier, orbital/fiber comparison, Siegel--Weil
comparison, local uniqueness and the downstream global spectral
identifications remain separate producers.
In particular this chapter does not close the whole selected
\texttt{GLOBAL\_\allowbreak SPECTRAL} obligation or assert the unrestricted
RH or Dirichlet-GRH endpoints.


\section{Selmer signatures on the actual representation family}
\label{sec:selmer-signature-v052}

\status{SSC1--SSC16 prove the full original Selmer-signature comparison,
including integral trace conjugacy, arbitrary local injective changes of
coefficient DVR, the original crystalline local conditions, infinite
lengths and the tame-conductor term. SSC17 removes the second of the two
remaining suppliers in DFD15. The coherent Wach functor and comparison
remain an explicit input to the family setup; the actual common
\texttt{Concl\allowbreak F} producer is not claimed here.}

Fix the original $D:\texttt{D7}$ and write $G=G_E$, $O_i$,
$K_i=\operatorname{Frac}(O_i)$, $\ell_i$, $X_i$ and
$Z_i^{\mathrm{good}}$ as in DFD. For a coefficient domain $O$ put
\[
 A_O^r=\operatorname{Frac}(O)^r/O^r.
\]
It has the discrete topology. An integral matrix representation acts
on this quotient by the same matrices. We use continuous cochains;
\(H^1(G,A)\) always means cocycles modulo coboundaries in that sense.
For a DVR $O$ and natural number $\ell$, let
\[
 n_O(M)=
 \begin{cases}
 \infty,&\operatorname{length}_O M=\infty,\\
 \displaystyle\frac{\operatorname{length}_O M}
 {\operatorname{length}_O(O/\ell O)},
       &\operatorname{length}_O M<\infty.
 \end{cases}                                                    \tag{SSC.1}
\]
The finite denominator in this notation uses the source's
\(\mathrm{toNat}\) convention if a module has infinite length, and
division by zero is zero in $\mathbb Q$. At every actual location
$\ell_i\ne0$ is a nonunit and the denominator is positive and finite.
The extended nonnegative version \texttt{len\allowbreak N} uses division in
$\mathbb R_{\ge0}\cup\{\infty\}$; we will distinguish these conventions
when comparing them.

For a signature $s$, $T_s$ denotes the source's \texttt{lat\allowbreak T}:
the representation $\rho_s$ twisted by the inverse of its specified
first character. For an injective coefficient map $f$ and character
$\psi$, write $T_{s,f,\psi}=f(T_s)\psi^{-1}$. The local condition
\texttt{ext\allowbreak Cond} is defined on cocycles: take the coefficient-ring
span of pushes of the original allowed cocycles, intersect with the
cocycles for $T_{s,f,\psi}$, and then pass to cohomology. Keeping this
order is necessary when $\psi$ need not be trivial locally.
The Selmer module \texttt{sel\allowbreak T} uses this condition above $\ell$,
no condition at the specified bad and auxiliary places, and the
unramified condition elsewhere. All places, embeddings, characters,
and bad sets below are those supplied by $D$.

\subsection{Integral representations and the local conditions}

% BEGIN SSC1
\paragraph{SSC1: trace equality gives integral conjugacy.}
\label{claim:SSC1}
Let $O$ be any commutative local ring, $G$ any group, and
$\rho,\rho':G\to\operatorname{GL}_a(O)$. If the reduction of
$\rho$ is absolutely irreducible and the two traces agree on $G$,
then some $U\in\operatorname{GL}_a(O)$ satisfies
\[
                    \rho'(g)U=U\rho(g)\quad(g\in G).             \tag{SSC.2}
\]
The case $a=0$ is immediate. For $a>0$, absolute irreducibility
and the finite-dimensional density theorem imply that the residual
matrices span $M_a(k)$, where $k$ is the residue field. One may
check the use of this standard theorem after an algebraic closure:
Burnside's theorem gives the full matrix algebra there, and dimension
descends under field extension. Select $a^2$ original matrices whose
reductions form a basis. Their coordinate determinant is a unit, so
they are already a basis of $M_a(O)$. Thus the group-algebra map
$\phi:O[G]\to M_a(O)$ induced by $\rho$ is onto.

Here is a direct local-ring argument for the remaining assertion.
Write $\phi'$ for the other group-algebra map and choose an
$O$-linear right inverse $s$ to $\phi$. Set $\theta=\phi's$ and
\[
 T(X)_{ij}=\operatorname{tr}\bigl(X\theta(E_{ji})\bigr).
\]
For every $x\in O[G]$, trace equality applied to $xs(E_{ji})$
gives $T(\phi'(x))_{ij}=\phi(x)_{ij}$. Consequently
$T\theta=1$. The square matrix of $T$ has a right inverse,
so its determinant is a unit and $T$ is bijective. It follows that
$\theta\phi=\phi'$. Surjectivity of $\phi$ now implies that
$\theta$ is a unital algebra homomorphism and preserves traces.

Fix an index $o$. Since $\operatorname{tr}\theta(E_{oo})=1$,
some diagonal entry $u=\theta(E_{oo})_{yy}$ is a unit: all
nonunits of a local ring lie in its maximal ideal. Define
\[
 U_{xi}=\theta(E_{io})_{xy},\qquad
 V_{jx}=\theta(E_{oj})_{yx}.
\]
The matrix-unit multiplication rules give $VU=uI$. Taking
determinants proves that $U$ is invertible. Those same rules,
applied to $XE_{io}=\sum_k X_{ki}E_{ko}$, give
$\theta(X)U=UX$. Substitute $X=\rho(g)$ to obtain (SSC.2).
This proves the original DVR conjugacy theorem and, in fact, its
stronger local-ring version, without choosing a new integral lattice.
% END SSC1

% BEGIN SSC2
\paragraph{SSC2: the actual rank-$r$ summand has the triple's trace.}
\label{claim:SSC2}
Retain the original \texttt{Cheb} premise. Fix $i,\pi$, an
\texttt{Is\allowbreak Rho\allowbreak Pt} representation $\rho_i(\pi)$, one list $S_g$
satisfying \texttt{Has\allowbreak Param\allowbreak X}, and a signature $s$ satisfying
\texttt{Is\allowbreak Triple} for that same list. Then
\[
 \operatorname{tr}\rho_i(\pi)(g)
   =f_i\bigl(\operatorname{tr}\rho_s(g)\bigr),\qquad g\in G,
 \quad f_i=D.\mathrm{ntw}.\mathrm{twist}_2(i).                  \tag{SSC.3}
\]
Indeed the actual rational splitting in \texttt{Is\allowbreak Rho\allowbreak Pt} gives
\[
 \operatorname{tr}R_\pi(g)
    =\operatorname{tr}\rho_i(\pi)(g)
       +\sum_{c=1}^{t}\chi_{i,c}(g).                           \tag{SSC.4}
\]
This follows by conjugating to the two diagonal blocks; it makes
no integral splitting assertion. At every permitted split Frobenius,
the Hecke polynomial identifies the trace on the left with the
normalized parameter of $S_g$ together with its specified character
block. The triple identifies the trace of $\rho_s$ with the parameter
of $S_g$ alone. In the source normalization the character block is
\[
 \sum_{c=0}^{t-1}\iota_{QC}(\bar\chi_{c+1}(\sigma_v))
   =\operatorname{sumPart}(\mu\eta,t;v)
          (Nv)^{(1-n_V)/2}.                                   \tag{SSC.5}
\]
To check this equality, insert
$\bar\chi_{c+1}=\bar\chi_1\bar\varepsilon^{-c}$ in each term.
The stated Frobenius values are $\bar\varepsilon(\sigma_v)=Nv$
and, after the fixed complex embedding,
$\bar\chi_1(\sigma_v)=\operatorname{sumEig}(\mu\eta;v)
(Nv)^{(t-1)/2+(1-n_V)/2}$. Add the real exponents in each
of the $t$ summands. The norm $Nv$ is a positive integer,
so the ordinary real-power addition law is valid. The same
normalizing factor occurs in both parameter identities; (SSC.5)
cancels the character summand of (SSC.4), giving (SSC.3) on the
permitted Frobenius set.

Both traces in (SSC.3) are continuous in the original $\ell$-adic
coefficient field. In particular $f_i$ is continuous: the inverse
image of the target maximal ideal is a nonzero prime of the source
DVR, hence its maximal ideal, and powers of that ideal map into
the corresponding powers in the target. The original embeddings
into the $\ell$-adic coefficient field are continuous. DFD3, with
the same finite avoidance and split-place requirements, extends the
identity to $G$. No continuity of an embedding into $\mathbb C$
has been used.

If the residual representation is absolutely irreducible, SSC1
gives an integral conjugacy with $f_i(\rho_s)$. The triple identifies
the first character with the product of the pushed triple character
and $\psi_L(i)$. Multiplying (SSC.2) by these scalar inverses gives
\[
 W_{\pi}(g)U^{-1}
     =U^{-1}T_{s,f_i,\psi_L(i)}(g).                            \tag{SSC.6}
\]
Here $W_\pi$ is exactly \texttt{Wp\allowbreak O}; inversion of $U$ accounts for
the source's direction of the intertwiner. This proves the full
\texttt{twisted\allowbreak Rho\_\allowbreak conj\_\allowbreak lat\allowbreak Tpsi} implication on the original
integral representations.
% END SSC2

% BEGIN SSC3
\paragraph{SSC3: the actual local bridges transport crystalline cocycles.}
\label{claim:SSC3}
Two embeddings $j_1,j_2:\Omega\to\overline{\mathbb Q}_\ell$
inducing the same valuation ring are related by an automorphism
$\sigma$ over $\mathbb Q_\ell$. We include the argument needed
to compare the source's independently specified bridges. Equality
of the closed unit balls gives equality of the open unit balls by
applying it to $x^{-1}$. For a finite tuple $x_a\in\Omega$,
the condition
$\|\sum c_a j_t(x_a)\|<1$ is clopen as a condition on
$(c_a)\in\mathbb Q_\ell^n$. The two conditions agree on rational
tuples, so density of $\mathbb Q^n$ makes them agree everywhere.
If a linear combination is zero for $j_1$ and nonzero for $j_2$,
multiply all coefficients by a scalar of sufficiently large norm
to contradict this agreement. Hence the two maps
$\mathbb Q_\ell\otimes_{\mathbb Z}\Omega\to
\overline{\mathbb Q}_\ell$ have the same kernel. The induced
isomorphism of their image fields extends to an automorphism of
the algebraic closure, as asserted.

When $\Omega$ is algebraically closed, its image is dense in
$\overline{\mathbb Q}_\ell$. For completeness, approximate the
coefficients of a monic polynomial of an algebraic $\ell$-adic
element by rational coefficients. Factoring the approximating
polynomial shows that one root is arbitrarily close to the chosen
element: otherwise its value there, a product of root differences,
could not tend to zero. Such roots belong to the image of
$\Omega$. All $\mathbb Q_\ell$-automorphisms are isometries, so
agreement on $j(\Omega)$ implies agreement everywhere. The two
decomposition-group identifications in the bridges therefore obey
the exact conjugacy
$g_2=\sigma g_1\sigma^{-1}$, rather than an unspecified
abstract group isomorphism.

Here is the cochain transport. A crystalline cocycle $c$ has a
witness $x\in V\otimes_{\mathbb Q_\ell}B$ with
$c(g)\otimes1=gx-x$, where $B$ is the period ring of the supplied
\texttt{Loc\allowbreak Prime}. If $q:V_1\to V_2$ is a
$\mathbb Q_\ell$-linear map intertwining the representations
under the above group conjugacy, use
\[
 q\otimes\sigma_B:V_1\otimes B\longrightarrow V_2\otimes B.
\]
On pure tensors the conjugacy identity proves equivariance;
linearity proves it everywhere. Applying this map to $x$ proves
that $q(c(g_1))$ is crystalline, with the displayed image of $x$
as witness. Its continuity follows from finite-dimensional
continuity of $q$ and continuity of the group maps. This proves
the full original \texttt{crystalline\allowbreak Cocycle\_\allowbreak transport\_\allowbreak conjugation}
with the actual supplied period-ring action.

The needed coefficient maps automatically respect
$\mathbb Z_\ell$. In a Noetherian local target with $\ell$ in its
maximal ideal, two homomorphisms from $\mathbb Z_\ell$ agree:
approximate $z$ by an integer modulo $\ell^n$; their values then
differ by an element of every power of the maximal ideal, whose
intersection is zero. This supplies the precise scalar compatibility
for the coefficient maps used below.
% END SSC3

% BEGIN SSC4
\paragraph{SSC4: integral coefficient extension and conjugacy preserve the local conditions.}
\label{claim:SSC4}
Let $O_0\hookrightarrow O$ be the original injective map of
coefficient DVRs finite over $\mathbb Z_\ell$. Let $B_0,B$ be
local bridges for the same valuation ring and the same
\texttt{Loc\allowbreak Prime}. Then for any integral matrix homomorphism $M_0$
(with no continuity assumption on $M_0$)
the extension, in the cocycle-span sense above, of its crystalline
local condition is exactly the crystalline condition of $f(M_0)$.

One inclusion follows by pushing a crystalline lift and its witness
through the coefficient map, then using SSC3 for the bridges.
To prove the reverse inclusion, note that $O$ is finite
torsion-free over $O_0$: finite generation over $\mathbb Z_\ell$
and SSC3 give finite generation over $O_0$. It is therefore free.
Choose an integral basis $b_a$ and fraction-field coordinate maps
$p_a:K\to K_0$. They satisfy
\[
 p_a(O)\subseteq O_0,\qquad
 p_a(f(y)x)=y p_a(x),\qquad x=\sum_a b_a f(p_a(x)).              \tag{SSC.7}
\]
They consequently induce maps $A_O^r\to A_{O_0}^r$ which
commute with the actions. Keep the source's literal witnesses for a
continuous quotient cocycle $c$ in its crystalline condition:
\[
 c(g)=[z(g)+M(g)v-v],
\]
where $z$ is a continuous crystalline $K$-valued cocycle and
$v\in K^r$. Apply each $p_a$ and SSC3 to $z$ and its period-ring
witness, and apply $p_a$ algebraically to $v$. The continuous
quotient coordinate map sends $c$ to a continuous cocycle, and
\[
 p_a(c(g))=[p_a(z(g))+M_0(g)p_a(v)-p_a(v)].
\]
Thus $p_a(c)$ has precisely the allowed original pair of witnesses
$(p_a(z),p_a(v))$. No continuity of the rational coboundary is
needed. Formula (SSC.7)
expresses the original quotient cocycle as a finite $O$-linear
combination of their pushes. This is exactly membership in
\texttt{ext\allowbreak Cond}, proving equality, including ramified coefficient
extensions.

Likewise any integral intertwiner $U\in\operatorname{GL}_r(O)$
induces inverse maps $[v]\mapsto[Uv]$ and $[v]\mapsto[U^{-1}v]$
on the quotient lattices and on continuous cochains. They commute
with restrictions and preserve inertia coboundaries, hence preserve
the unramified conditions. For the crystalline conditions use
$U\otimes1$ on the continuous lift's period-ring witness, and
apply $U$ separately to its rational boundary vector. Thus they give an actual
isomorphism of the full Selmer modules and preserve their lengths,
including infinite length.

For the source's $s,i$ and $\psi_L(i)$, the latter character is
trivial on every decomposition group above $\ell_i$. The local
representation there is exactly the coefficient extension of
$T_s$, and the first assertion applies. The source identities
$\ell_{\mathrm{twOf}(i)}=\ell_i$ and
$\mathrm{setS}_{\mathrm{twOf}(i)}=S_i$ identify all other local
conditions. Applying the integral conjugacy (SSC.6) yields the
full original equality
\[
        \texttt{s\allowbreak Pi\allowbreak Top}(D,i,\rho_i,\pi)
             =\texttt{s\allowbreak TLoc}(D,i,s).                          \tag{SSC.8}
\]
The local field bridges and the integral lattices have remained
those in the original datum.
% END SSC4

\subsection{Arbitrary coefficient DVRs and the same character}

% BEGIN SSC5
\paragraph{SSC5: compatible coordinates in every torsion quotient.}
\label{claim:SSC5}
Let $g:O\hookrightarrow O'$ be any injective local homomorphism
of DVRs. There is no finiteness or completeness hypothesis here.
Choose a uniformizer $\varpi$ of $O$ and lift a vector-space basis
$(\bar b_a)_{a\in I}$ of $O'/g(\varpi)O'$ over $O/\varpi O$
to elements $b_a\in O'$. The index set may be infinite. For every
$0\ne d\in O$ these lifts give a basis of $O'/g(d)O'$ as an
$O/dO$-module.

It suffices to take $d=\varpi^n$. For generation, expand an
element modulo $\varpi$, subtract the finite expansion, divide
the remainder by $g(\varpi)$, and repeat $n$ times. Only finitely
many basis indices occur. For independence, reduce a putative
finite relation modulo $\varpi$. Every coefficient is divisible
by $\varpi$; cancel $g(\varpi)\ne0$ and repeat. Each coefficient
is divisible by $\varpi^n$. Multiplication by a unit proves the
assertion for general $d$.

Every element of $K'=\operatorname{Frac}(O')$ has a denominator
from $O$: $g(\varpi)$ has positive $O'$-valuation, so a large
power clears any denominator. Passing to the quotients by $O'$
and using the just proved generation and independence gives
\[
 \bigoplus_{a\in I} A_O^r \;\xrightarrow{\ \sim\ }\; A_{O'}^r,
 \qquad (x_a)\longmapsto\sum_a b_a\,\mathrm{push}(x_a).        \tag{SSC.9}
\]
Write $p_a$ for its $O$-linear coordinate maps. For matrices over
$O$ they commute with the action, and every vector has finite
coordinate support. These statements hold uniformly in every
denominator; a new basis has not been chosen at each torsion level.
% END SSC5

% BEGIN SSC6
\paragraph{SSC6: the full scalar-change theorem for normalized Selmer lengths.}
\label{claim:SSC6}
For the exact source's fixed $D,i,s,Q$, any injective
$f:O_{0,i}\to O$ and any injective local map $g:O\to O'$ of
DVRs, and any character $\psi:G\to O^\times$, one has
\[
 sT(D,i,s,Q,gf,g\psi)=sT(D,i,s,Q,f,\psi).                     \tag{SSC.10}
\]
Neither $O'$ nor $O$ is required to be finite over the other.

Apply (SSC.9) to $T_{s,f,\psi}$. Continuous cocycles with values
in a discrete module have finite image because $G$ is compact.
Thus each cocycle uses only finitely many of the coordinates $p_a$,
and the coordinate functions are continuous cocycles. The same
statement holds for a coboundary, using the finite support of its
defining vector. Consequently cohomology admits the corresponding
coordinate decomposition, compatible with restrictions to every
decomposition or inertia group. The latter groups are compact too.

It remains to check the exact crystalline-span condition, since a
coordinate map need not be $O'$-linear. A defining original cocycle
has finite image and hence is killed by one $0\ne d\in O$.
For any scalar $c\in O'$, expand $c$ modulo $g(d)O'$ using SSC5.
On all values of that cocycle, multiplication by $c$ is therefore
a finite combination of multiplication by the $b_a$ with
coefficients in $O$. Applying a coordinate to any finite sum
of such scalar multiples remains in the original defining
$O$-span. Pushes of allowed cocycles remain in the new span as
well. Intersecting with the cocycle modules and passing to
cohomology respects these maps. Inertia coboundaries are preserved
by direct equivariance, and the unrestricted local conditions add
no constraint. Thus, for the two Selmer modules $S,S'$, the natural
map is an isomorphism
\[
                     O'\otimes_O S\simeq S'.                 \tag{SSC.11}
\]
For clarity this assertion also holds for an arbitrary torsion
$S$ in the coordinate argument: any finite tensor sum is killed
by one $d$, and the bases of $O'/g(d)O'$ give both generation and
independence. This proves injectivity as well as surjectivity;
no exchange with an infinite product of local conditions is needed.
Each coordinate individually respects every local condition.

Let $e=v_{O'}(g(\varpi))>0$. Flatness of the torsion-free
$O$-module $O'$ and a composition series give
\[
 \operatorname{length}_{O'}(O'\otimes_O S)
          =e\operatorname{length}_O S.                        \tag{SSC.12}
\]
For infinite length, there are finite-length torsion submodules of
arbitrarily large length: otherwise maximal such length forces
the whole module to equal one finite submodule. Tensoring their
inclusions proves (SSC.12) in the infinite case. The denominator
$\operatorname{length}(O/\ell O)$ scales by the same $e$.
Cancellation proves (SSC.10). If the source's finite denominator
is zero, both finite normalized lengths are zero; if the numerator
has infinite length, both \texttt{nlen} values are infinity before
any division. This verifies the source conventions in all cases.
% END SSC6

% BEGIN SSC7
\paragraph{SSC7: the same-character relation produces a common DVR.}
\label{claim:SSC7}
Suppose $O_0,A,B$ are DVRs, $f_A:O_0\hookrightarrow A$ is
integral, $f_B:O_0\hookrightarrow B$ is local, and
$\psi_A:G\to A^\times$, $\psi_B:G\to B^\times$ satisfy the
original \texttt{Same\allowbreak Char}: every finite $O_0$-linear relation
among the first character values holds if and only if the
corresponding relation among the second values holds. There are a
DVR $O_1$, a map $f:O_0\to O_1$, injective local maps
$g_A:O_1\to A$, $g_B:O_1\to B$, and one character
$\psi:G\to O_1^\times$ whose images are exactly $\psi_A,\psi_B$,
with $g_Af=f_A$ and $g_Bf=f_B$.

Evaluate the group algebra $O_0[G]$ at the two characters. The
same-character condition says that the two kernels agree. Its
image $R\subset A$ therefore embeds in $B$. Let
$K=\operatorname{Frac}(R)$ and denote the induced field embeddings
by $j_A,j_B$. Put $O_1=K\cap A$, meaning the inverse image of $A$
under $j_A$. If $x\in O_1$, then $j_A(x)$ is integral over $O_0$
because $A$ is. A monic polynomial witnessing this also holds
for $j_B(x)$, so integral closedness of $B$ implies $j_B(x)\in B$.
This constructs both injections.

The injection into $A$ reflects divisibility: if $j_A(x)$ divides
$j_A(y)$, then $y/x\in K\cap A$. Hence divisibility in $O_1$
is total, and strict divisibility cannot descend indefinitely,
as its image cannot do so in a DVR. More concretely the valuation
of $A$ restricts to a subgroup of $\mathbb Z$ on $K^\times$;
$O_1$ is its nonnegative valuation ring. It is not a field:
a uniformizer of $O_0$ stays a nonzero nonunit in $B$ by locality
of $f_B$. Thus the restricted value group is a nonzero discrete
group, making $O_1$ a DVR. The map to $A$ is local by divisibility
reflection. If a uniformizer of $O_1$ became a unit in $B$, then
so would the nonzero nonunit from $O_0$, by its DVR factorization.
The map to $B$ is local as well. Character values and their inverses
already lie in $R$, so they define the required units of $O_1$.
Every stated equality follows from evaluation on group-algebra
generators. This proves the whole original common-DVR theorem.
% END SSC7

% BEGIN SSC8
\paragraph{SSC8: the actual same twist preserves Selmer length and local triviality.}
\label{claim:SSC8}
Apply SSC7 with $A=O_i$, $B=O_{\mathrm{twOf}(i),k}$,
$f_A=f_i$ and $f_B=\texttt{to\allowbreak Ok}$. These maps have the required
properties. The first is integral because its target is finite
over $\mathbb Z_\ell$ and its two $\mathbb Z_\ell$ structures
agree by SSC3: a monic polynomial over $\mathbb Z_\ell$ is also
one over the source. For locality of the second map, use the
following elementary test. If $f:O\to B$ maps a nonzero natural
number $\ell$ to a nonunit, factor $\ell=u\varpi^n$. Then
$f(\varpi)$ is a nonunit. Every source nonunit is divisible by
$\varpi$, and hence cannot map to a unit. At the actual coefficient
rings $\ell$ is nonzero and remains a nonunit by integral finite
extension. This establishes locality.

The common character and common DVR of SSC7 give two applications
of SSC6, one to $g_A$ and one to $g_B$. Their middle normalized
length is identical. Thus the original same-character premise
gives the full equality
\[
 sTLoc(D,i,s)
      =sT(D,\mathrm{twOf}(i),s,Q_i,\texttt{to\allowbreak Ok},\psi).        \tag{SSC.13}
\]
In addition, \texttt{Same\allowbreak Char} preserves equality of a character
value to $1$. Apply its defining relation to $[g]-[1]$ in the
group algebra. Therefore $\psi_L(i)=1$ on a decomposition group
implies $\psi=1$ there. In particular the transported character
is trivial at all places above $\ell$, exactly as required by the
following comparison. No assertion that the two original coefficient
rings are subrings of one another has been made.
% END SSC8

\subsection{From local crystalline classes to the signature Selmer module}

% BEGIN SSC9
\paragraph{SSC9: a cocycle-level inclusion into the original signature condition.}
\label{claim:SSC9}
For every $D,i,s,k,Q$ and $\psi:G\to O_{i,k}^\times$ trivial
on the decomposition groups above $\ell_i$, there is an injective
map from the source's \texttt{sel\allowbreak T} to the signature Selmer module
$\texttt{Gen.\allowbreak sel\allowbreak T}$. Both use the same $\psi$.

The coefficient map underlying this comparison is the elementary
isomorphism
\[
 K^r/O^r\longrightarrow(K/O)^r,
       [v]\longmapsto([v_1],\ldots,[v_r]).                    \tag{SSC.14}
\]
Its inverse takes any coordinate representatives. They differ
only by an integral vector. The formula
$[Mv]_a=\sum_b M_{ab}[v_b]$ shows equivariance. Inserting the
definitions of $T_{s,f,\psi}$, \texttt{chi\allowbreak S} and \texttt{Trep}
shows that these are the literal two representations in the source.

Here is the local verification, including the lift which its
definition requires. Let $c$ be an original quotient cocycle in
the crystalline condition. Its class is represented by the
reduction of a crystalline rational cocycle $z$, up to a
coboundary. Absorb the coboundary in $z$. For each $g$, choose a
representative $u_g\in K_0^r$ of $c(g)$. Under the local bridge,
the difference between $u_g$ and $z(g)$ is an integral vector.
Subtract its preimage in $O_0^r$ from $u_g$; call the result $y_g$.
Then $[y_g]=c(g)$ and the bridge sends $y_g$ exactly to $z(g)$.
This construction makes no assertion that an arbitrary choice
of representatives is continuous: continuity belongs to $c$ and
$z$, which are the objects required in the definition.

Push the $y_g$ to the new fraction field. To verify the source's
explicit fraction witness, write any coordinate $y=n/d$ with
$n,d\in O_0$, $d\ne0$. The identities
$dy=n$ survive both the bridge and the fraction-field map. These
are precisely the two denominator equations in the definition of
\texttt{Gen.\allowbreak bk}. The crystalline cocycle there is $z$, whose
representation is \texttt{Gen.\allowbreak rep\allowbreak V}: both are the specified
inverse-character twist of $\rho_s$. Local triviality of $\psi$
ensures the coefficient push respects this local action.
Equation (SSC.14) now puts each defining original cocycle into
\texttt{Gen.\allowbreak bk}. Linear spans and the passage to cohomology preserve
this inclusion.

At the unrestricted places there is nothing to check; at all other
places (SSC.14) preserves inertia coboundaries. The original sets
$\mathrm{setS}$ and $\mathrm{Spl}$ coincide. Thus (SSC.14) induces
the claimed map of the global Selmer modules. It is injective
because its inverse is equivariant on the entire coefficient module
and hence on $H^1$. This proves the source's whole
\texttt{length\_\allowbreak le} comparison with no finite-length assumption.
% END SSC9

% BEGIN SSC10
\paragraph{SSC10: finite and infinite normalized lengths.}
\label{claim:SSC10}
Under the hypotheses of SSC9, write $S$ for the first Selmer module
and $S_{\mathrm{sig}}$ for the second. Then
\[
 \begin{split}
 \texttt{s\allowbreak Len}(S_{\mathrm{sig}})<\infty
      &\ \Longrightarrow\ n_O(S)<\infty,\\
 n_O(S)=q\in\mathbb Q
      &\ \Longrightarrow\ \operatorname{ofReal}(q)
                    \le\texttt{s\allowbreak Len}(S_{\mathrm{sig}}).
 \end{split}                                                   \tag{SSC.15}
\]
The injection of SSC9 gives the inequality of module lengths. If
the first length were infinite, the second would be infinite too.
Its extended nonnegative normalized length would then be infinite,
because the denominator is a finite natural number. This proves
the first assertion. For finite first length, division of the
length inequality by the common positive denominator proves the
second. This is the actual coefficient situation. At the broader
scope of the source's elementary \texttt{nlen\_\allowbreak le\_\allowbreak len\allowbreak N} lemma,
the denominator may be zero. The finite rational normalization
on the left is then zero, so the second assertion is automatic;
the infinity argument for the first remains valid. This covers
both original conventions rather than silently identifying them.

Combining (SSC.8), (SSC.13), and (SSC.15) gives the full source
\texttt{s\allowbreak T\_\allowbreak le\_\allowbreak s\allowbreak Len} implication and its required application to
the actual point Selmer module. If that module is infinite, a finite
signature length is impossible. If it is finite, its rational
\texttt{s\allowbreak Pi} value is exactly the rational value in (SSC.1).
% END SSC10

% BEGIN SSC11
\paragraph{SSC11: one attached signature proves the full Selmer core.}
\label{claim:SSC11}
Assume the source's \texttt{Cheb}, \texttt{Triples\allowbreak Good},
\texttt{Good\allowbreak Pt\allowbreak Row}, and \texttt{Res\allowbreak Abs\allowbreak Irr}. At a point
$\pi\in Z_i^{\mathrm{good}}$, let $\rho_i(\pi)$ satisfy
\texttt{Is\allowbreak Rho\allowbreak Pt}, and let the target character satisfy the full
original \texttt{Same\allowbreak Char} relation with $\psi_L(i)$ at the
specified coefficient maps. There is a signature $s$ which is
good, has the prescribed side, is attached to this actual point,
and has both conclusions (SSC.15) for its point Selmer module.

Choose the row $\sigma,S_g$ supplied by \texttt{Good\allowbreak Pt\allowbreak Row}.
Apply \texttt{Triples\allowbreak Good} to that same row and parameter list.
This gives one $s$, its goodness, and its \texttt{Is\allowbreak Triple}.
Residual absolute irreducibility applies to the same
$\rho_i(\pi)$. SSC2 and SSC4 give (SSC.8); SSC8 gives
(SSC.13) and triviality of the target character above $\ell$;
SSC9--SSC10 give the two length conclusions. The triple's side
field gives the required side, and the collected row, parameter,
point and triple data give the original \texttt{Attached\allowbreak In}.
All five output fields therefore use the same $s$. This proves
the whole original \texttt{sel\_\allowbreak at} theorem, including arbitrary
target index $j$, its equality with $\mathrm{twOf}(i)$, and the
specified coefficient map transported by that equality.

Now take the original $C:\texttt{NChoice}$ and suppose
\texttt{NCh\allowbreak Tie} and \texttt{Same\allowbreak Twist}. At every $i$, the latter
provides the identical coefficient map after the prescribed
index identification and the character
\[
 \psi_G(C,w)=(C.\psi'(w))\circ\operatorname{art},
       \qquad w=D.\mathrm{nd}.w(i).                           \tag{SSC.16}
\]
Use the preceding construction with $k=C.k'(w)$. The tie identifies
$Q_i=C.Q(w)$ and the side with $C.\mathrm{sd}(w)$.
It gives every clause of \texttt{Sel\allowbreak Core}, including
\texttt{Attached\allowbreak In}, for every actual family of \texttt{Is\allowbreak Rho\allowbreak Pt}
representations. In particular it applies to the \texttt{Pt\allowbreak Fam}
constructed in DFD7, because $Z_i^{\mathrm{good}}\subseteq Z_i$.
% END SSC11

\subsection{Tame conductors, residual irreducibility, and the final join}

% BEGIN SSC12
\paragraph{SSC12: normalized valuations agree in the two coefficient realizations.}
\label{claim:SSC12}
Let a field $L$ with valuation $w$ embed into fraction fields of
DVRs $R_1,R_2$, and suppose both embeddings identify the condition
$w(x)\ge0$ with integrality. Let $\ell$ be nonzero and a nonunit
in both DVRs. If $x\in L$ has integral images $a,b$, then
\[
             \frac{v_1(a)}{v_1(\ell)}
                  =\frac{v_2(b)}{v_2(\ell)}.                 \tag{SSC.17}
\]
At $x=0$ use the source's $\mathrm{toNat}(\infty)=0$ on both
sides. Otherwise, for all natural $p,q$,
\[
 \ell^q\mid a^p
   \ \Longleftrightarrow\ x^p/\ell^q\text{ is integral}
   \ \Longleftrightarrow\ \ell^q\mid b^p.
\]
Write $A_j=v_j(\ell)>0$, $B_1=v_1(a)$, $B_2=v_2(b)$.
This says $qA_1\le pB_1$ if and only if $qA_2\le pB_2$.
Use $(p,q)=(A_1,B_1)$ and $(A_2,B_2)$ to obtain both
inequalities between $B_1A_2$ and $B_2A_1$, proving (SSC.17).

Also
\[
 \operatorname{lenUnit}_{\ell}(R/(x))
       = nval_{R,\ell}(x),\qquad
 nval_{R,\ell}(ux)=nval_{R,\ell}(x)\ (u\in R^\times).         \tag{SSC.18}
\]
These follow from the uniformizer factorization and the length of
$R/\varpi^nR$. For $x=0$ both sides of the first equality use the
same zero convention; it is not a claim that $R$ has finite length.
The same valuation argument proves invariance under local
coefficient extensions, by canceling the common ramification index.
% END SSC12

% BEGIN SSC13
\paragraph{SSC13: tame generators and inversion of the character.}
\label{claim:SSC13}
For the actual $v\in Q_i$, the finite-order character $\psi_L(i)$
is trivial on wild inertia. Here is an elementary verification of
the exact input needed. A wild automorphism $\tau$ satisfies
$|\tau(y)/y-1|_v<1$ for every nonzero $y$, and so do all its
nonnegative powers, by expanding $(1+a)(1+b)-1$.
If $\tau^k x=x$ with $|k|_v=1$, then $\tau x=x$. Otherwise put
$d=\tau x-x\ne0$; telescoping gives
\[
 k=-\sum_{j=0}^{k-1}\left(\frac{\tau^j(d)}d-1\right),
\]
contradicting $|k|_v=1$. The open kernel of $\psi_L(i)$ contains
the subgroup fixing some finite extension. A power $\tau^N$
fixes that extension. Write $N=\ell^a m$ with $\ell\nmid m$.
Since $v\nmid\ell$, the preceding observation shows that
$\tau^m$ already fixes it. Thus $\psi_L(i)(\tau)^m=1$.
This value also has $\ell$-power order by the original character
field, so it equals $1$.

Let $\gamma=\gamma_{i,v}$ be the exact supplied tame generator.
The inverse image of $\langle\psi_L(i)(\gamma)\rangle$ is an open,
and hence closed, subgroup. It contains $\gamma$ and wild inertia.
The supplied tame-generator density therefore puts all inertia
inside this inverse image. Thus each inertia value is an integral
power of $\psi_L(i)(\gamma)$. The shared character $\psi_0$
and both injective coefficient realizations in \texttt{NCh\allowbreak Tie}
transport the identical equality of powers to $\psi_G(C,w)$.
This proves the original full \texttt{Tame\allowbreak Cyc} conclusion.

For any unit $u$ and integer $m$ with $u^m\ne1$, $u-1$ divides
$u^m-1$: use the geometric sum for $m\ge0$ and multiply by
$-u^m$ for negative $m$. Hence
$nval(u-1)\le nval(u^m-1)$. The source's tame value is the
infimum of these normalized valuations over nontrivial inertia
values. If the character value at $\gamma$ is nontrivial, it
belongs to that set and gives a lower bound for every element;
thus
\[
 nval(\psi_G(\gamma)-1)\le\operatorname{tameVal}(\psi_G).
                                                                    \tag{SSC.19}
\]
If it is trivial, all inertia values are trivial and the defining
set is empty; the source uses $\inf\varnothing=0$ in $\mathbb R$,
so the same inequality holds. Finally
$u^{-1}-1=-u^{-1}(u-1)$ implies
\[
           \operatorname{tameVal}(\psi^{-1})
                     =\operatorname{tameVal}(\psi).           \tag{SSC.20}
\]
This includes the trivial-character case and retains the actual
generator and coefficient embeddings.
% END SSC13

% BEGIN SSC14
\paragraph{SSC14: the full conductor-root bound at the tied auxiliary places.}
\label{claim:SSC14}
Assume \texttt{NCh\allowbreak Tie} for $D,C$. For every $i$ and every
$v\in Q_i$,
\[
 \operatorname{lenUnit}_{\ell_i}
        \bigl(O_i/(\zeta_{i,v}-1)\bigr)
  \le \operatorname{tameVal}_{\ell_i}
          (\psi_W(C,w),v),\qquad w=D.\mathrm{nd}.w(i).         \tag{SSC.21}
\]
Indeed $Q_i$ is disjoint from $S_i$ and from the primes above
$\ell_i$. The first twist character is unramified at $v$, since
all its possible ramification places belong to $S_i$. Its value
on $\gamma_{i,v}$ is therefore $1$. The exact character product
in the datum gives
$\zeta_{i,v}=\psi_L(i)(\gamma_{i,v})$.

The two realizations of $\psi_0$ in \texttt{NCh\allowbreak Tie} identify
$\psi_L(i)(\gamma)-1$ and the image of
$\psi_G(C,w)(\gamma)-1$ in $C.O_e(w)$ as images of one element
of the original number field. Both embeddings induce the same
specified place. The prime $\ell_i$ is nonzero and a nonunit in
all these rings: first in the original coefficient rings, and then
in the finite integral extensions defining $O_{i,k}$ and $C.O_e$.
SSC12 thus identifies their normalized valuations, and its last
assertion removes the extra coefficient extension. By definition
\[
 \psi_W(C,w)=\operatorname{psiAt}(C.\mathrm{sd}(w),C.\psi'(w))
                         \circ\operatorname{art}.
\]
It is either $\psi_G$ or $\psi_G^{-1}$. Apply (SSC.19)--(SSC.20)
and (SSC.18) to obtain (SSC.21). Summing it and using the literal
tie $Q_i=C.Q(w)$ gives the full second clause of
\texttt{Signature\allowbreak Choice\allowbreak Condition}. The zero and empty-set cases
have already been included; no strict positivity of a conductor
root is assumed.
% END SSC14

% BEGIN SSC15
\paragraph{SSC15: the original residual hypothesis covers every chosen point representation.}
\label{claim:SSC15}
The source's \texttt{Residual\allowbreak Irreducible\allowbreak Rank\allowbreak R} implies its exact
\texttt{Res\allowbreak Abs\allowbreak Irr}. To see this, fix $\pi\in Z_i^{\mathrm{good}}$
and any \texttt{Is\allowbreak Rho\allowbreak Pt} realization $\rho$. Its actual rational
block decomposition gives the characteristic-polynomial identity
\[
 \operatorname{charpoly}(R(g))
    =\operatorname{charpoly}(\rho(g))
           \prod_{c=1}^{t}(X-\chi_{i,c}(g)).                   \tag{SSC.22}
\]
Both sides have coefficients in $O_i$ and the fraction-field map
is injective, so (SSC.22) is an integral equality. Reduce it modulo
the maximal ideal and use the specified residual Hecke-polynomial
tie at the permitted Frobenius elements. Together with the open
kernel of the continuous residual representation, these are
exactly the fields of \texttt{Is\allowbreak Rho\allowbreak Bar}. The good point also gives
the source's \texttt{HG} premise. The original residual hypothesis
therefore makes the representation irreducible over the given
algebraically closed residual coefficient field.

This is absolute irreducibility over the original residue field.
Explicitly embed its algebraic closure into that given field.
A nonzero invariant subspace extends to a nonzero invariant
subspace there. A proper subspace remains proper: choose a
nonzero linear functional killing it and extend that functional
by the field embedding. Thus any nonzero proper invariant subspace
would contradict the stated irreducibility. The rank is positive
by the original ranks record. This proves \texttt{Abs\allowbreak Irr\allowbreak Sub} on the
full space and hence \texttt{Res\allowbreak Abs\allowbreak Irr}, uniformly for every chosen
\texttt{Is\allowbreak Rho\allowbreak Pt} lattice, not only a preferred representative.
% END SSC15

% BEGIN SSC16
\paragraph{SSC16: the whole original Selmer-signature supplier.}
\label{claim:SSC16}
The following is now proved at its full original scope:
\[
\begin{gathered}
 \texttt{Cheb},\quad \texttt{Pt\allowbreak Fam}(D,R,\rho),\quad
 \texttt{NCh\allowbreak Tie}(D,C),\\
 \texttt{Triples\allowbreak Good}(D),\quad\texttt{Same\allowbreak Twist}(D,C),\quad
 \texttt{Good\allowbreak Pt\allowbreak Row}(D),\quad\texttt{Res\allowbreak Abs\allowbreak Irr}(D)\\
 \Longrightarrow\quad
            \texttt{Signature\allowbreak Choice\allowbreak Condition}(\rho,C).        \tag{SSC.23}
\end{gathered}
\]
Apply SSC11 to each $\pi\in Z_i^{\mathrm{good}}$ and use the
\texttt{Pt\allowbreak Fam} field on that point. Its signature has side
$\mathrm{side}(s)=C.\mathrm{sd}(w)$. By the two-case definition
of \texttt{psi\allowbreak At},
\[
 \operatorname{psiAt}(C.\mathrm{sd}(w),\psi_W(C,w))
                        =\psi_G(C,w).                         \tag{SSC.24}
\]
For the true side nothing is inverted; for the false side the
same character is inverted twice. The side equality is essential:
two unrelated Boolean choices could not be canceled. Therefore
both the finiteness implication and the normalized-length
inequality of SSC11 are exactly the first clause of the required
condition, with its specified signature character. SSC14 supplies
its second clause. These statements hold at every location and
therefore at every original \texttt{Tw\allowbreak Loc}; no enlargement of the
location set is required by the theorem. This proves the entire
original \texttt{Joint\allowbreak Six.\allowbreak main} supplier, retaining the supplied
\texttt{Cheb} premise and every other original input.

With \texttt{Tw\allowbreak Good\allowbreak At} on the same $C$ and $I$, DFD11 now gives
$\texttt{Kill\allowbreak At}(\rho)$ directly. The budget remains exactly
$1/\ell^2$ and the local term is the same tame term in (SSC.21).
This conclusion is about the actual chosen representation family.
It neither replaces the family nor chooses new twists after
the length estimates.
% END SSC16

% BEGIN SSC17
\paragraph{SSC17: the exact family setup now needs only the coherent Wach supplier.}
\label{claim:SSC17}
In DFD15 retain the original family, constants, \texttt{Tw\allowbreak Good},
\texttt{Facts\allowbreak F}, \texttt{Cheb}, and just its first extra supplier:
for every relevant $D$, actual compatible
$HW,F_W,\texttt{Cmp\allowbreak Of}$ on $D.\mathrm{loc},D.\mathrm{locA}$.
Then the full original output of \texttt{exists\_\allowbreak setup\_\allowbreak large}
follows. Its separate Selmer-signature supplier is no longer an
extra premise.

Choose once the maximum of the parameter, small-prime, and
\texttt{Tw\allowbreak Good} cutoffs in DFD13--DFD15. For any later
$x,h_x,i$, retain the specified $D,L,C_D=\texttt{nch}$ and twist
index $I$. Choose the coherent $HW,F_W,\mathrm{Cmp}$ using the
remaining supplier. DFD9 constructs the actual canonical
$X$ with $X.I=I$, error function zero and exceptional-set function
constantly the arbitrary given $\mathrm{exc}$. In particular
$X.\mathrm{exc}=\mathrm{exc}$, and its pair $(X.R,X.\rho)$ has
the original \texttt{Pt\allowbreak Fam}.

The member of \texttt{Facts\allowbreak F} supplies \texttt{NCh\allowbreak Tie},
\texttt{Triples\allowbreak Good}, and \texttt{Same\allowbreak Twist}. DFD10 supplies
\texttt{Good\allowbreak Pt\allowbreak Row} from that member's actual row conclusions.
Its residual rank-$r$ hypothesis gives \texttt{Res\allowbreak Abs\allowbreak Irr} by
SSC15. SSC16 applies to this precise $X.R,X.\rho,C_D$, so DFD11
gives the killing field. The facts, canonical, two parameter,
small-scale, split-bound and supplied \texttt{Sha\allowbreak Fin\allowbreak Odd} fields
are the other seven fields already assembled in DFD15.
The same parameter implication gives \texttt{Adm\allowbreak J}.

The cutoff is chosen before $x,h_x,i$, before the
\texttt{Sha\allowbreak Fin\allowbreak Odd} premise, and before every finite exceptional
set. All objects in the output are the ones just constructed.
This proves the exact conditional setup statement with only the
coherent local comparison supplier still displayed. DFD12 proves
its finite-coefficient Bhatt--Scholze reduction, but construction
of the whole coherent Wach system remains. The later local row
bounds and the complete common \texttt{Concl\allowbreak F} family also remain;
\texttt{GLOBAL\_\allowbreak DENOMINATOR} is consequently still open.
% END SSC17

\subsection{Source alignment and limits of the result}

The pinned source is commit
\texttt{7cbf7a0c32fca1badf05931879a37c3f8a2b03b4}.
The accompanying source manifest records the exact files, ranges,
bytes and hashes. In the \texttt{Inertia\allowbreak Eigenvalues} proofs the
main endpoints are the trace and twisted-conjugacy theorems,
\texttt{SPi\allowbreak Top\allowbreak Eq\allowbreak STLoc}, \texttt{Is\allowbreak Local\allowbreak Hom\allowbreak Of\allowbreak Nat\allowbreak Cast},
\texttt{Nlen\allowbreak Le\allowbreak Len\allowbreak N}, \texttt{Same\allowbreak Char\allowbreak Eq\allowbreak One},
\texttt{Conductor\allowbreak Root\allowbreak Sub\allowbreak One\allowbreak Length\allowbreak Le\allowbreak Tame\allowbreak Value},
\texttt{Residual\allowbreak Abs\allowbreak Irred\allowbreak Good\allowbreak Point\allowbreak Of\allowbreak Residual\allowbreak Irred}, and
\texttt{Psi\allowbreak At\allowbreak Psi\allowbreak At}. The foundation files for integral trace
conjugacy, common coefficient rings and arbitrary injective local
coefficient changes are reconstructed in SSC1 and SSC5--SSC8;
their names are not used as substitute proofs.

Classical algebraic foundations used here are the
finite-dimensional density theorem, field-embedding extension,
basic DVR structure and module length, and the usual compactness
and valuation properties of absolute Galois groups and finite
$\ell$-adic fields. The original local period-ring and bridge
structures, original tame-generator density, and original
\texttt{Cheb} are retained where the source assumes them.
No new automorphy, crystalline existence, or Wach comparison
theorem is inserted among these foundations.

\endgroup
\end{document}
